\section{F. The later Montpellier topoi}
\label{sec:F}

We distinguish three constructions throughout this section: a pro-object of
full neighbourhood topoi, the classifying topos of covering germs, and a topos
of sheaves constructible for a fixed stratification. They have different
universal properties. We write $B G$ for the topos of continuous discrete
sets with an action of a profinite group $G$, and $B\mathcal G$ for its
groupoid version. A morphism over a base topos includes its commutation
\emph{isomorphism}; forgetting that isomorphism changes several answers below.

\subsection{F(a): the six presentations and their rigidifications}

\Needspace{5\baselineskip}
\paragraph{F.1 (Proved). The six presentations, with the information retained.}
\textit{Manuscript: MPL 149; section1; folios1-1v and8v-9; scans3-4 and18-19.}

Fix a profinite group $\Gamma$ and its one-object profinite groupoid
$B_{\mathrm{gr}}\Gamma$. The following are presentations of the same
unpointed object, provided the indicated two-dimensional data are retained.
\begin{enumerate}
\item Connected point-admitting profinite Galois topoi with a geometric morphism to
$B\Gamma$ whose induced group homomorphism has open image. No point is
selected. Morphisms are geometric morphisms together with a specified
isomorphism over $B\Gamma$; two-morphisms respect that isomorphism.
\item Connected profinite groupoids over $B_{\mathrm{gr}}\Gamma$ with open image, with functors,
commutation isomorphisms, and compatible natural isomorphisms.
\item One-object presentations of (2), modulo inner conjugacy, \emph{together
with the conjugating elements as two-morphisms}. If only exterior classes
are retained, this is the homotopy category, not the original two-category.
\item Homomorphisms $p:E\to\Gamma$ of profinite groups with open image. A morphism from
$p':E'\to\Gamma$ to $p:E\to\Gamma$ is a pair $(u,\gamma)$ with
$u:E'\to E$ continuous and $p u=\operatorname{Ad}(\gamma)p'$.
A two-morphism $e:(u,\gamma)\Rightarrow(v,\delta)$ satisfies
$v=\operatorname{Ad}(e)u$ and $\delta=p(e)\gamma$ (the equivalent inverse
convention gives the same two-category).
\item If $H=p(E)$ is open, the extension $1\to\pi\to E\to H\to1$, together
with its embedding $H\hookrightarrow\Gamma$. If $Z(\pi)=1$, one can instead
retain the continuous outer action $H\to\operatorname{Out}(\pi)$, with its
extension topology; morphisms require the intertwiners and their coherence.
\item A profinite groupoid with continuous weak $\Gamma$-action, whose
homotopy quotient is connected and whose object-component $\Gamma$-set is
finite; equivariant functors and equivariant natural isomorphisms are its
morphisms. Strict actions may be used after the usual transport-groupoid
strictification, without discarding its coherence.
\end{enumerate}
Without the open-image restriction, replace the finite component set in (6)
by an arbitrary profinite homogeneous $\Gamma$-space, and omit the word \emph{open} in (5).
These descriptions are equivalent for the data arising from a finitely
generated field $K/\mathbb Q$. In (5), $H=\Gamma_k$, where $k$ is the
algebraic closure of $\mathbb Q$ in $K$, embedded in a fixed
$\overline{\mathbb Q}$; in (6) the groupoid is that of
$K\otimes_{\mathbb Q}\overline{\mathbb Q}$. If $K$ is finite over
$\mathbb Q$, then $\pi=1$: the subgroup $H$, or equivalently the finite
$\Gamma$-set $\Gamma/H$, must still be retained. The trivial geometric group
alone does not distinguish number fields.

\begin{proof}
For a profinite groupoid, evaluation at each object is a fibre functor on
continuous actions. Equivariant natural transformations of evaluation are
exactly the continuous arrows of the groupoid: finite transitive actions
detect each finite quotient, and their inverse limit detects the arrow.
Conversely the fibre functors of a pointed connected profinite Galois topos
recover the finite Galois category and its continuous-action topos. This
proves (1)$\simeq$(2), including all natural isomorphisms, without choosing
one object globally. Choosing an object of a connected groupoid gives its
one-object skeleton; changing that object is an isomorphism torsor, which
is exactly the inner-conjugacy datum in (3)--(4). On one-object groupoids
the naturality equation is the displayed conjugacy equation; compatibility
with the base is its second displayed equation. Consequently forgetting
the conjugating elements loses the centralizer of $u(E')$ intersected with
$\ker p$, and is a biequivalence only on Hom-groupoids where that group is
trivial.

The geometric fibre of $B_{\mathrm{gr}}E\to B_{\mathrm{gr}}\Gamma$ is the
transport groupoid with objects $\Gamma$ and arrows given by right
translation by $E$ through $p$. Its components are $\Gamma/H$, its vertex
groups are $\pi$, and left translation supplies the $\Gamma$-action.
Taking the transport homotopy quotient reverses this construction. This is
also a direct verification of effective Galois descent of actions, so it
proves (4)$\simeq$(6), rather than identifying a quotient group with a
quotient groupoid.

For a centreless kernel, conjugation identifies an extension with
\[
 E\ \simeq\ \operatorname{Aut}(\pi)
       \mathbin{\times}_{\operatorname{Out}(\pi)} H.
\]
Indeed its kernel is $Z(\pi)$, and any pair on the right differs from the
conjugation action of a lift of its $H$-coordinate by a unique inner
automorphism. This proves both injectivity and surjectivity and also the
uniqueness of the extension isomorphism. In the profinite situation all
maps are continuous; the extension topology is retained, so the assertion
does not silently replace continuous extensions by abstract extensions.
For a general kernel the outer action alone has the usual central
extension ambiguity: already $H=C_2$, $\pi=C_2$ and trivial outer action
admit $C_4$ and $C_2\times C_2$. This proves exactly what description (5)
would otherwise lose. Finally the components of
$K\otimes_{\mathbb Q}\overline{\mathbb Q}$ are indexed by embeddings
$k\hookrightarrow\overline{\mathbb Q}$, and their stabilizer is $\Gamma_k$.
\end{proof}

\Needspace{5\baselineskip}
\paragraph{F.2 (Proved). Openness and the centralizer used at scan 20.}
\textit{Manuscript: MPL 149; section1; folio9v; scan20.}

Let $F\subset F'$ be finitely generated extensions of an algebraically
closed characteristic-zero field. Then the image of
$G_{F'}\to G_F$ is open. Every open subgroup of $G_F$ has trivial
centralizer in $G_F$ when $\operatorname{trdeg}F>0$. For transcendence
degree zero, $G_F=1$ and its centralizer is again the trivial group.
Consequently the geometric exterior homomorphism at MPL~149, fol.~9,
scan~20 determines a geometric morphism of classifying topoi up to unique
isomorphism. This assertion is about the geometric groups; openness alone
is not a centralizer theorem for arbitrary profinite groups.

\begin{proof}
The relative algebraic closure $L$ of $F$ in $F'$ is finite over $F$.
The regular extension $F'/L$ is linearly disjoint from every finite
separable extension of $L$. Every automorphism of $L^{\mathrm{sep}}/L$
therefore extends to the compositum with $F'$, and then to a separable
closure of $F'$. Thus the image is $G_L$, an open subgroup of $G_F$.

We give the needed slimness argument. For every finite Galois extension
$L/F$ with group $Q$ and every prime $\ell$, there exists a finite Galois
extension $M/F$ above it such that the action of $Q$ on
$N=\operatorname{Gal}(M/L)$ is faithful. Choose a normal model of $F$ and
shrink it so that the normalization in $L$ is finite etale. A closed point
in its smooth locus has a full $Q$-orbit above it, since the constant field
is algebraically closed. A divisorial valuation centred there (use the
exceptional divisor of the blow-up when the dimension exceeds one) has
$|Q|$ distinct prolongations to $L$. Approximation for these finitely many
independent discrete valuations gives $a\in L^\times$ with valuation one
at one prolongation and zero at all the others. Its conjugates are
independent in $L^\times/L^{\times\ell}$, as their valuation vectors are
the standard basis. Since $\mu_\ell\subset F$, adjoining all their
$\ell$th roots gives $N\simeq(C_\ell)^Q$, with faithful permutation action.
If $z$ is central in $G_F$, its image in $Q$ acts trivially on $N$,
so faithfulness gives that image equal to one. Since $L/F$ was arbitrary,
$z=1$. The same argument applies to every open subgroup, because its fixed
field is again finitely generated over the same algebraically closed
constant field. Finally, if $z$ centralizes an open subgroup $H$, then
$J=\langle H,z\rangle$ is open, $H$ is normal in $J$, and $z$ lies in
$Z(J)$. The already proved centrelessness of $J$ gives $z=1$. Thus every
open subgroup has trivial centralizer in $G_F$. Natural automorphisms of a restriction functor are precisely
the centralizer of the image, proving the final assertion.
\end{proof}

\Needspace{5\baselineskip}
\paragraph{F.3 (Proved). Faithfulness of fields in every characteristic.}
\textit{Manuscript: MPL 149; section1 Theorem2; folios10-10v; scans21-22.}

Let $k$ be any field and $K,K'$ finitely generated extensions. Two
$k$-embeddings $f,g:K\to K'$ inducing isomorphic morphisms of profinite
Galois groupoids \emph{over} $\Pi_k$ are equal. Thus the characteristic-zero
margin to Theorem~2, MPL~149, fol.~10, scan~22, is correct; a
positive-characteristic counterexample with the same relative data does
not exist.

\begin{proof}
For each integer $n$ invertible in $k$, the coefficient $\mu_n$ is pulled
back from $B_k$. The specified relative isomorphism therefore identifies
the pullbacks in Kummer cohomology, with the specified identification of
coefficients, so
\[
 f(a)/g(a)\in K'^{\times n}\qquad(a\in K^\times).
\]
The same remains true after any algebraic extension of the constants and
on each connected component of the geometric base change. Choose a
component over an algebraic closure $\bar k$ of $k$ and write $F'$ for its
function field. For $a$ transcendental over $k$, both $f(a)$ and $g(a)$
remain nonconstant in $F'$. On a normal proper integral $\bar k$-model of
$F'$, every prime-divisor valuation of their ratio is divisible by all
integers prime to the characteristic, and hence zero. A rational function
with zero divisor on a normal proper integral variety over $\bar k$ is a
constant: regularity in codimension one and normality make it an everywhere
regular unit, and properness makes it an element of $\bar k^\times$.
Consequently $g(a)=\lambda f(a)$. Applying the same argument to $a-1$
gives $g(a)-1=\mu(f(a)-1)$. Nonconstancy gives $\lambda=\mu=1$.

For $a$ algebraic over $k$, the finite separable part of $k(a)/k$ is an
object of the finite etale category of $k$. The relative commutation datum
recovers its embedding in $K'$ by the universal property of the associated
slice topos. A purely inseparable extension has at most one embedding
extending that separable embedding. Thus $f(a)=g(a)$ also in this case.
The map from $K'$ into the chosen geometric component field is injective
(it is a map from a field), so equality there implies equality in $K'$.
\end{proof}

\Needspace{5\baselineskip}
\paragraph{F.4 (Proved). Relative faithfulness of reduced schemes.}
\textit{Manuscript: MPL 149; section2; folio25v; scan52.}

Let $S$ be locally noetherian and let $X,Y$ be locally of finite
presentation over $S$, with $X$ reduced. If $f,g:X\to Y$ induce isomorphic
morphisms $X_{\mathrm{et}}\to Y_{\mathrm{et}}$ in the two-category over
$S_{\mathrm{et}}$, then $f=g$. No characteristic-zero hypothesis and no
reducedness hypothesis on $Y$ are necessary. In particular the requested
positive-characteristic refutation of this relative statement cannot be
given.

\begin{proof}
First suppose $S=\operatorname{Spec}k$. Isomorphic geometric morphisms
induce the same map on subterminal objects and hence the same map on
underlying Zariski spaces. At a closed point $x\in X$ their common image
$y\in Y$ is closed, and both residue fields are finite over $k$. The two
maps on residue fields have the same effect on the maximal separable
subextension of $\kappa(y)/k$: a finite separable extension represents a
finite etale object of $B_k$, and the slice-topos universal property
recovers its actual map, not merely a conjugacy class. The remaining
purely inseparable extension has a unique embedding. Thus $f$ and $g$
agree on every closed residue point. Locally on a common inverse image of
an affine open in $Y$, their difference on any coordinate function vanishes
in every closed residue field. The source is reduced and Jacobson, so the
intersection of its maximal ideals is zero. Hence the two ring maps, and
therefore the morphisms, agree.

For general $S$, the square of etale topoi obtained by base change along
$\operatorname{Spec}\kappa(s)\to S$ is cartesian. One can verify this by
writing the local scheme as the inverse limit of affine open
neighbourhoods of $s$ and then taking its closed residue subtopos; this is
also \cite[Lemma~2.9]{CHW2024}. Pull the relative isomorphism to this square,
and then restrict to $(X_s)_{\mathrm{red}}$. The field case applies to this
reduced finite-type source and proves equality on every reduced fibre.
The morphism
$\coprod_s(X_s)_{\mathrm{red}}\to X$ is schematically dominant for reduced
$X$: a function vanishing on every such fibre vanishes in every residue
field and hence is zero. Since the underlying point maps already agree,
the same affine-local argument proves $f=g$. All reductions are local on
$S,X,Y$, so no global quasi-compactness is required.

For comparison with \cite[Theorem~3.9]{CHW2024}, the extra reduced-fibre
hypothesis in its stated version can be removed here by explicitly
restricting to the reduced fibres, as above. This argument asserts
faithfulness, not fullness.
\end{proof}

\Needspace{5\baselineskip}
\paragraph{F.5 (Refuted). Unrestricted
full faithfulness is false; pinned reconstruction is the correct theorem.}
\textit{Manuscript: MPL 149; section2; folio25v; scan52.}

For every characteristic-zero field $k$, morphisms of etale topoi over
$B_k$ need not arise from scheme morphisms, even after universal
homeomorphisms are inverted. If $k$ is an infinite finitely generated
field, the functor becomes fully faithful into the category of
\emph{pinned} geometric morphisms, which send closed points to closed
points. The same holds for sub-$p$-adic characteristic-zero fields.

\begin{proof}
The $t$-adic decomposition subgroup gives a section of
$G_{k(t)}\to G_k$: explicitly embed $k(t)$ in $k((t))$, use the splitting
of $G_{k((t))}\to G_k$ obtained by acting on coefficients and fixing
compatible roots of $t$, and restrict to a separable closure of $k(t)$.
The resulting morphism $B_k\to B_{k(t)}\to(\mathbb P^1_k)_{\mathrm{et}}$
over $B_k$ has image the generic point. A $k$-point of $\mathbb P^1_k$
has closed image, and universal homeomorphisms do not change that
obstruction. This is the explicit obstruction in
\cite[Example~2.30]{CHW2024}.

For an infinite field finitely generated over its prime field,
\cite[Theorem~0.1]{CHW2024} gives an equivalence on the Hom-sets of the
localization at universal homeomorphisms and the pinned Hom-groupoids.
The latter are equivalent to sets by its Proposition~2.22. The source
localization is the category of absolute weak normalizations: in
characteristic zero this is seminormalization, and in positive
characteristic it is perfection. Thus the theorem applies to the
manuscript's reduced finite-type schemes with exactly this localization
and this target restriction. The assertion for sub-$p$-adic fields is
\cite[Theorem~A, equivalently Theorem~4.1]{CarlsonStix2024}; its
identification of the localization with seminormal schemes gives the
same reduction. These statements do not classify arbitrary base fields.
\end{proof}

\Needspace{5\baselineskip}
\paragraph{F.6 (Refuted). Even pinned reconstruction fails over every
algebraically closed characteristic-zero field.}
\textit{Manuscript: MPL 149; section2; folio25v; scan52.}

Let $k$ be algebraically closed of characteristic zero. Put $j=i$ and
$j'=-i$, and let
\[
 E_j:\quad y^2=x^3-\frac{3j}{j-1728}x-\frac{2j}{j-1728}.
\]
Then $E_j$ and $E_{j'}$ are nonisomorphic smooth proper $k$-curves but
have equivalent etale topoi over $B_k\simeq\mathbf{Set}$. This remains a
counterexample after universal homeomorphisms are inverted.

\begin{proof}
The displayed Weierstrass equation has nonzero discriminant and
$j$-invariant $j$; the symbolic computation is in
\path{code/montpellier/montpellier_checks.py}. Extend the automorphism
$i\mapsto-i$ of $\overline{\mathbb Q}$ to $k$, by choosing and fixing a
transcendence basis and then extending to an algebraic closure. Applying
this automorphism to coefficients gives an isomorphism of the underlying
schemes and hence an equivalence of their small etale topoi. It is over
$\mathbf{Set}$ automatically and is pinned because it is a homeomorphism.
A $k$-isomorphism of the curves would preserve the $j$-invariant,
contradicting $i\ne-i$. The curves are normal, hence seminormal, so their
objects in the localized characteristic-zero category remain themselves.
\end{proof}

\Needspace{5\baselineskip}
\paragraph{F.7 (Known). The experiment with a proper
curve over a number field.}
\textit{Manuscript: MPL 149; section2; folio25v; end of scan52.}

Let $k$ be a number field, $C/k$ a smooth proper geometrically connected
curve of any genus, $K=k(C)$, and
$\pi=G_{K\bar k}$. The $k$-automorphisms of $C$ are in bijection with
continuous outer automorphisms of $G_K$ over $G_k$, and equivalently with
continuous outer automorphisms of $\pi$ commuting with the outer
$G_k$-action. Requiring preservation of the conjugacy classes of the
inertia groups of closed points does not change this answer. Smoothness
and normality are essential if the curve, rather than its normalization,
is to be recovered from its function field.

\begin{proof}
Apply \cite[Theorem~B, Theorem~17.1 in the body]{Mochizuki1999} to the
regular function field $K/k$ on both sides; $k$ is sub-$p$-adic. An
isomorphism of groups is open, so it and its inverse give field embeddings.
Their composites are the identity by the theorem's uniqueness, hence
are inverse $k$-automorphisms. A birational map between smooth proper
curves extends uniquely at every discrete valuation ring by the valuative
criterion of properness, and the inverse extends as well. Thus it is an
automorphism of $C$. Conversely any such automorphism induces the stated
outer group automorphism and permutes the local valuations and their
inertia. F.2 makes $\pi$ centreless; the extension pullback in F.1 therefore
turns a $G_k$-equivariant outer automorphism of $\pi$ into the required
outer automorphism of $G_K$. This verifies the proposed experiment at
MPL~149, \S2, fol.~25v, scan~52, including genus zero and genus one; it uses the
birational group, not just the proper curve's fundamental group.
\end{proof}

\subsection{F(c): quasi-Galois gluing and ordinary exit categories}

\Needspace{5\baselineskip}
\paragraph{F.8 (Proved). A precise finite quasi-Galois characterization.}
\textit{Manuscript: MPL 147; section12; folio85; scans179-180; MPL149 folio67 scan135.}

Suppose a Grothendieck topos $\mathcal T$ has a finite filtration by open
subtopoi, each successive locally closed stratum being the presheaf topos
of a groupoid with finitely many connected components. At each step write
$F=i^*j_*:\mathcal T_{\rm open}\to\mathcal T_{\rm closed}$ for the
accessible left-exact gluing functor. Then $\mathcal T$ is, compatibly with
the filtration, a presheaf topos on an EI-category with finitely many
isomorphism classes if and only if every $F$ preserves arbitrary products.
Here an EI-category is a category in which every endomorphism is
invertible, and its isomorphism classes are partially ordered by the
existence of a morphism. Without the finiteness of the component sets of
the strata, the corresponding conclusion uses a small EI-category and
need not have finitely many isomorphism classes.

\begin{proof}
We prove one induction step. Write the open topos as
$\widehat A=\operatorname{Fun}(A^{\rm op},\mathbf{Set})$ and the closed
one as $\widehat B$, with $B$ a groupoid. A sheaf in the glued topos is a
triple $(P,Q,q:Q\to F(P))$. If $F$ preserves products, left exactness
makes it preserve every small limit. Accessibility and the adjoint
functor theorem give a left adjoint $L$ to $F$. For each representable
$h_b\in\widehat B$ put $M_b=L(h_b)\in\widehat A$. Yoneda and adjunction
give
\[
 F(P)(b)=\operatorname{Hom}_{\widehat B}(h_b,F(P))
          \simeq\operatorname{Hom}_{\widehat A}(M_b,P).
\]
Define a category $C$ by adjoining to $A$ the objects of $B$, putting
$\operatorname{Hom}_C(a,b)=M_b(a)$, no morphisms from $B$ to $A$, and
using the actions supplied by the presheaves $M_b$ and functoriality in
$b$ for composition. A presheaf on $C$ is exactly the triple above.
This proves $\mathcal T\simeq\widehat C$. All new endomorphisms lie in
$A$ or $B$, so the EI property follows by induction. No cycle can cross
the finite filtration; after passing to isomorphism classes the resulting
preorder is a partial order.

Conversely, for an EI-category with the given filtration, crossing arrows
form the bimodule $M_b(a)=\operatorname{Hom}_C(a,b)$. The gluing functor
has the displayed representable form and preserves all limits,
particularly products. The numbers of isomorphism classes add under the
finite construction. A single stratum $\mathbf{Set}^{\mathbb N}$ shows
why that finite-component condition cannot be omitted: its indecomposable
projective generators have infinitely many nonisomorphic types.
This proves the categorical case asserted in STRATA B.1--B.2 rather than
assuming that arbitrary left-exact gluing is representable by crossing
sets.
\end{proof}

\Needspace{5\baselineskip}
\paragraph{F.9 (Proved). A two-stratum gluing which is not of presheaf EI type.}
\textit{Manuscript: MPL 147; section12; folio85; scan179; task F(c).}

Let $C=\{0,1\}^{\mathbb N}$ be the Cantor space, and define
$F(A)=\operatorname{LC}(C,A)$ for a set $A$ with the discrete topology.
The Artin gluing of two copies of $\mathbf{Set}$ along $F$ is a
Grothendieck topos with two Galois strata, but it is not the presheaf
EI-topos asserted in F.8.

\begin{proof}
Every locally constant map from compact $C$ has finite image and factors
through a finite clopen partition. Consequently
$F=\operatorname*{colim}_{\mathcal P}(-)^{\mathcal P}$, where
$\mathcal P$ runs through the filtered set of finite clopen partitions.
It is accessible and preserves finite limits, so its Artin gluing is a
Grothendieck topos. For the countable family $A_n=\{0,1\}$, the coordinate
projections $C\to A_n$ are locally constant. They define an element of
$\prod_nF(A_n)$ which does not come from $F(\prod_n A_n)$: the only map
with these coordinates is the identity of the underlying set of $C$,
which is not continuous from Cantor space to that set with the discrete
topology. Thus $F$ fails to preserve this product, and F.8 proves the
claim. This example also specifies the missing gluing condition in the
undefined term ``topos quasi-galoisien'' at MPL~147, fol.~85, scan~179.
\end{proof}

\Needspace{5\baselineskip}
\paragraph{F.10 (Proved). Constructible sets and the crossing double cosets.}
\textit{Manuscript: MPL 147; sections11-12; folios56-64 and85; scans114-130 and179-180; MPL149 scans133-135.}

Let $(X,D)$ be a complex manifold with a locally finite normal-crossings
divisor, or an analytic Deligne--Mumford stack with such an etale local
presentation. Assume the usual paracompact local-conical hypotheses for
constructible exodromy. Write $\mathcal E=\operatorname{Exit}(X,D)$,
with arrows from closed to more open strata. Then
\[
 \operatorname{Sh}_{\rm str}(X;\mathbf{Set})
       \simeq\operatorname{Fun}(h\mathcal E,\mathbf{Set})
       =\widehat{(h\mathcal E)^{\rm op}}.
\]
This is a quasi-Galois topos in the sense of F.8 (without finite component
restriction if necessary). For one connected incidence link $L$ between
connected strata $S,T$, choose endpoints and put
$A=\pi_1S$, $B=\pi_1T$, $G=\pi_1L$. The crossing-arrow set, with its
$A$--$B$ actions, is
\[
             (A\times B)/\operatorname{im}(G\to A\times B),
\]
with the harmless opposite convention on one factor dictated by endpoint
transport. If several link components occur, take the disjoint union of
these orbit sets. Composition is induced by concatenation on the full
flag spaces, not by composing their sets of components without the flag
maps.

\begin{proof}
A constructible sheaf transports germs along paths within each stratum.
In a conical chart a generization from a closed stratum is transport along
an outgoing radial arc; homotopic exit arcs induce the same map of germs.
The sheaf axiom on conical charts and their intersections gives precisely
compatibility under concatenation, and produces a functor on
$h\mathcal E$. Conversely, for such a functor assign on a conical open
chart the compatible germs indexed by its exit objects. Restriction is
restriction of compatible families. The charts form a basis stable under
refinement, and compatibility on a cover is exactly the sheaf condition;
the stalk at an exit object is the assigned set. On each stratum the
construction is the usual locally constant sheaf associated to a
fundamental-groupoid action. These two constructions are inverse on
stalks, hence on sheaves, and on natural transformations. For the stack
apply this construction on the entire etale atlas nerve and descend:
compatible charts include the stabilizer actions. This is the
set-valued reduction of the exodromy theorem, with all transition data
retained.

The exit mapping space for this incidence is the homotopy fibre of
$L\to S\times T$ at the chosen endpoints. The long exact sequence in
components of that fibre gives the orbit-set formula. Its residual
endpoint actions give the two module actions. Two composable exits are
represented by a flag with two steps; the face forgetting the middle
step is their composition. Three-step flags verify associativity. Within
a stratum every exit is invertible, and a strict exit cannot return to its
source stratum, so all endomorphisms are invertible. Finally the bimodule
formula for restrictions shows that its gluing functor is represented as
in F.8. This verifies the relevant STRATA C.1 and C.6 calculation in the
geometric setting; it does not identify the full sheaf topos of $X$ with
this constructible subcategory.
\end{proof}

\Needspace{5\baselineskip}
\paragraph{F.11 (Refuted). The Boolean source axiom does not imply thinness.}
\textit{Manuscript: MPL147; section11; folios56-59; scans114-121.}

The weighted Boolean axiom (174)c at MPL~147, fols.~56--57,
scans~115--116, implies (174)a and b. It does \emph{not} imply that an
arrow is determined by its source and target. The additional geometric
transversality hypothesis later imposed at scans~120--121 gives this
property in the intended branch-incidence examples.

\begin{proof}
The coslice of a rank-$d$ object has the weighted Boolean lattice of
subsets of a $d$-element set. Its ranks are between zero and $d$, and it
has a unique object of rank $d$, namely the identity. This proves a and b.
For the false extra implication use the finite category with objects
$x,y,z$ of ranks $2,1,0$, arrows
\[
 a,b:x\longrightarrow y,\qquad c:y\longrightarrow z,
 \qquad d:x\longrightarrow z,
 \qquad ca=cb=d,
\]
and identities. Its coslices are the Boolean diamond, the two-element
chain, and the one-point poset, respectively. Thus it satisfies the full
weighted Boolean coslice condition, but $a\ne b$ have identical source
and target. Regard all these finite sets as discrete zero-dimensional
analytic spaces to obtain an internal analytic category. The program
\texttt{montpellier\_checks.py} verifies every composition and every
coslice relation.

The exact correction to the axiom is injectivity of the target map on
every source coslice, or equivalently monicity of
$(s,t):D_{**}\to D_*\times D_*$. For genuine normalized branch germs, an
object is a specified set of local branches, so an arrow exists exactly
when one specified set contains the other, and its inclusion is unique.
Under the later transverse geometric realization, distinct subsets give
distinct intersection germs by independent local coordinates, so this
extra condition holds. It is the geometric realization, not the
abstract Boolean axiom alone, that supplies it.
\end{proof}

\subsection{F(d): germs, full sheaves, and descent}

\Needspace{5\baselineskip}
\paragraph{F.12 (Proved). The punctured germ is a pro-topos; its covering
reflection is a different object.}
\textit{Manuscript: MPL147; section11; folios70-73; scans143-148; LM8 E.10-E.11.}

For a locally closed analytic embedding $Z\hookrightarrow Y$, let
$\mathfrak V_{Z,Y}$ be the pro-object of full sheaf topoi
$(\operatorname{Sh}(V))_V$ indexed by open neighbourhoods of $Z$ in an
ambient open where $Z$ is closed. With a divisor $D$, put
$\mathfrak V^*_{Z,Y}=(\operatorname{Sh}(V\setminus D))_V$.
These definitions are invariant under cofinal refinement. Their
covering-germ categories are the filtered two-colimits of the covering
categories at the stages. For a normal-crossing point with $r$ branches
this category is $\mathbb Z^r$-sets, and its finite part is finite
continuous $\widehat{\mathbb Z}(1)^r$-sets. In contrast the literal
inverse two-limit of the full sheaf topoi of shrinking punctured discs
is degenerate.

\begin{proof}
Cofinal systems present the same pro-object by the definition of
pro-morphisms. Polar coordinates identify a cofinal system of punctured
normal polydiscs with a contractible radial factor times $(S^1)^r$;
shrinking acts identically on the coordinate winding numbers. Covering
classification therefore gives the asserted filtered category, including
its morphisms after common restriction. This verifies all steps of
LM8 E.10 independently.

For the last assertion let $D_n^*=\{0<|z|<1/n\}$ and suppose a topos
$\mathcal A$ maps compatibly to every $\operatorname{Sh}(D_n^*)$. Viewed
as opens of $D_1^*$, the annuli $U_k=\{1/k<|z|<1\}$ cover $D_1^*$.
Each inverse image of $U_k$ is initial because the morphism factors
through $D_k^*$, disjoint from $U_k$. Inverse images preserve unions,
so the terminal object of $\mathcal A$ is initial. The degenerate topos
has the unique such cone, proving the universal property. This also
verifies LM8 E.11. The manuscript's arcs and $C^\infty$ embeddings at
scans~143--148 identify these as classical analytic germs; a profinite
etale replacement is a separate construction, not a new reading of the
word ``germe''.
\end{proof}

\Needspace{5\baselineskip}
\paragraph{F.13 (Proved). The precise full-topos gluing theorem.}
\textit{Manuscript: MPL147; sections11-12; folios74 and76; scans150 and155.}

Let $X$ be paracompact and locally compact, $D$ a locally finite
normal-crossings divisor, and $\Sigma\subset D$ a closed union of
strata. A cofinal system of controlled tubular neighbourhoods of $\Sigma$
can be covered by tubular neighbourhoods of its normalized branches,
with all their intersections. The resulting full Cech diagram, or its
full flag subdivision, has two-colimit
$\mathfrak V^*_{\Sigma,X}$ in pro-Grothendieck topoi. This is the correct
full-sheaf meaning of both versions of the gluing theorem at MPL~147,
fols.~74 and 76, scans~150 and 155. The same theorem holds for analytic
Deligne--Mumford stacks presented by proper etale groupoids with
paracompact locally compact orbit space, or for an explicitly controlled
atlas whose descent conditions admit the common-stage refinement used
below. Arbitrary nonproper atlases, even with finite isotropy, require
that extra hypothesis: F.51 refutes passage from objectwise atlas germs to
an intrinsic stack neighbourhood without it. For a closed $D$ one also has
\[
 \operatorname{Sh}(X)\ \simeq\
 \operatorname{Sh}(X\setminus D)
       \amalg_{\mathfrak V^*_{D,X}}\mathfrak V_{D,X}
 \quad\text{in pro-topoi},
\]
where the left side is constant. This proves the full-topos version of
(254) and (266), with their germ objects correctly typed.

\begin{proof}
At a fixed neighbourhood stage $V$, choose open tubular sets $V_a$ whose
union is $V$ (and intersect them with $X\setminus D$ for the punctured
version). Sheaves on $V$ are exactly sheaves on the $V_a$ with
isomorphisms on double intersections satisfying the triple cocycle
condition. Here is the topoi-level universal property: for any target
$\mathcal A$, a geometric morphism $\operatorname{Sh}(V)\to\mathcal A$
is a finite-limit- and colimit-preserving functor
$\mathcal A\to\operatorname{Sh}(V)$. Restricting that functor produces
such functors on every $V_a$ with compatible isomorphisms. Conversely
objectwise sheaf descent glues these functors and natural transformations;
finite limits and colimits are checked after the conservative open
restriction family. Thus the gluing is a bicategorical colimit of
\emph{topoi}, not merely a statement about covering categories.

For finitely many local branches, subdividing the Cech index category by
all nonempty nested intersections does not change descent data: a
subdivided datum is recovered by restriction from its one-step data, and
the flag triangles are precisely the cocycle relations. For completeness the passage from local germs to a global pro-object
requires a uniformization argument. Fix a target topos $\mathcal A$ and a
compatible germ datum. On a locally finite, relatively compact shrinking
of the coordinate cover choose a representative neighbourhood for every
chart datum, every overlap isomorphism, and every triple cocycle equality.
Near any given point only finitely many of these conditions occur.
After shrinking the coordinate sets once more, their allowed radii have
a positive lower bound on each compact member of that shrinking. A
partition of unity gives a positive continuous function smaller than all
these locally finite bounds. Restricting to that variable-radius tube
realizes every datum and every compatibility at one stage. The same
argument applied to two representatives, and then to their two-morphisms,
proves full faithfulness as well as essential surjectivity. Thus finite
chart descent passes to germs and global open descent is legitimate;
we have not commuted a filtered colimit with an arbitrary infinite
limit. For a pro-target first apply this argument to each fixed target
stage, and then take the limit over target stages in the pro-mapping
formula. This proves the first assertion and its version for $\Sigma$.

For the second assertion, at every stage $V$ the pair
$X\setminus D,V$ is an open cover of $X$, with intersection
$V\setminus D$. Apply the just-proved two-open universal property and
pass through the same cofinal system. For non-simple normal crossings,
choose an etale chart atlas and all its fibre products; the normalized
local branches are handled on this atlas, not assumed globally embedded.
For a proper etale groupoid the local orbifold charts have finite
stabilizers. Choose a locally finite shrinking on the orbit space and
take the minimum over the finitely many stabilizer translates of each
chart radius. Over compact chart closures, properness makes the relevant
transition conditions finite after an etale refinement. The preceding
uniformization therefore applies to the overlap maps and their cocycles
as well as to the branch data. Tubular constructions and intersections
pull back etale locally, and these synchronized descent diagrams may be
composed. The result retains stabilizer groups and requires no invariant
tubular embedding. For a nonproper groupoid this finite-control step
must be supplied separately; F.51 gives a case where it fails. The
intrinsic pro-object indexed by actual stack neighbourhoods still makes
sense, and its two-open pushout is proved at each actual common stage;
it is not thereby identified with the objectwise atlas-germ diagram.
\end{proof}

\Needspace{5\baselineskip}
\paragraph{F.14 (Refuted). Covering topoi cannot replace full elementary
topoi in an ordinary two-colimit.}
\textit{Manuscript: MPL147; section12; folios85-86; scans179-180; LM8 E.11.}

For discrete groupoids $\mathcal G_i$, the ordinary bicategorical colimit
of the topoi $B\mathcal G_i$ is
$B(\operatorname*{2\!colim}_i\mathcal G_i)$. It generally has the wrong
higher shape for the geometric gluing. In particular
\[
 \mathbf{Set}\amalg_{B\mathbb Z}\mathbf{Set}\simeq\mathbf{Set},
 \qquad *\amalg_{S^1}^{h}*\simeq S^2.
\]
Thus neither ``aspherical strata'' nor ``all flags'' repairs the
replacement of full elementary topoi by covering topoi inside the
ordinary two-category. One must instead take the homotopy colimit of
classifying \emph{spaces}, or use the corresponding infinity-topoi.

\begin{proof}
Present a topos $\mathcal A$ by a small site. A geometric morphism
$B\mathcal G\to\mathcal A$ is a flat continuous site functor valued in
$\mathcal G$-sets. Evaluation at the objects of $\mathcal G$ identifies
this with a $\mathcal G$-diagram of points of $\mathcal A$, with natural
transformations: flatness and covering epimorphisms are checked
objectwise. Hence
$\operatorname{Geom}(B\mathcal G,\mathcal A)
\simeq\operatorname{Fun}(\mathcal G,\operatorname{Pt}\mathcal A)$.
The bicategorical colimit universal property for groupoids now proves
the first assertion. The groupoid pushout of $1\leftarrow\mathbb Z\to1$
is the trivial groupoid, whereas the homotopy pushout of its classifying
spaces is the suspension of $S^1$. The pair $(\mathbb P^1,\infty)$ is the
geometric instance. This correction is consistent with F.13: its full
neighbourhood topoi have not been replaced by their covering reflections.
\end{proof}

\subsection{Finite flag truncations and van Kampen}

\Needspace{5\baselineskip}
\paragraph{F.15 (Refuted). The two-depth diagram (246).}
\textit{Manuscript: MPL147; section11; folio75; scan152; (246).}

For the coordinate divisor in $\mathbb C^3$, the bicategorical colimit
of the depth-one and depth-two groupoids in (246), with the incidence
isomorphisms but without the depth-three flag relations, has vertex
group $\mathbb Z^4$. The punctured normal boundary has fundamental
group $\mathbb Z^3$. These statements concern discrete groups; their
finite-covering analogues have the indicated profinite completions.

\begin{proof}
In polar coordinates the boundary of the real oriented blow-up is
\[
 L=(S^1)^3\times R,\qquad
 R=\{(r_1,r_2,r_3)\in\mathbb R_{\geq0}^3:\min r_i=0\}.
\]
The factor $R$ contracts to its vertex. Removing depth three deletes
that vertex, and radial normalization gives $R\setminus\{0\}\simeq
\partial\Delta^2\simeq S^1$. Thus $L_{\leq2}\simeq(S^1)^4$.
The index category has three singleton and three two-element subsets,
with six incidences: its nerve is a hexagon. Each total-link group is
$\mathbb Z^3$ and each incidence identifies its three coordinate
meridians. A spanning tree eliminates five incidence generators; the
remaining generator commutes with the three meridians, giving
$\mathbb Z^3\times\mathbb Z$. Adding the three-element subset and all
flags makes the index nerve a cone on that hexagon. Its six triangular
relations eliminate the six non-tree generators in the enlarged
presentation. This gives $\mathbb Z^3$.

The attached program constructs the subset posets, their edges and
triangles, chooses a spanning tree, and performs single-occurrence
Tietze eliminations; the outputs are respectively $(6,6,0)$ and
$(7,12,6)$ for vertices, edges, triangles, with one and zero surviving
index generators. This checks the actual bicategorical presentation,
not just its abelianization. The extra generator is precisely the
coherence lost by a strict colimit of one-object groups.
\end{proof}

\Needspace{5\baselineskip}
\paragraph{F.16 (Proved). Sharp uniform depth bound.}
\textit{Manuscript: MPL147; section11; folio75; scan152; EXIT P5a2.}

Let $L$ be the boundary of the real oriented blow-up of a smooth
normal-crossings pair, and let $L_{\leq r}$ be its part over points
with at most $r$ local branches. Include every incidence flag among
the retained depths. For $r\geq1$ the map
$L_{\leq r}\to L$ is $(r-1)$-connected, with the convention that an
$m$-connected map induces isomorphisms on $\pi_i$ for $i<m$ and a
surjection for $i=m$. Hence depths at most $k+2$ compute
$\tau_{\leq k}L$. This is the smallest uniform depth bound. The same
assertion holds on the classifying spaces of analytic
Deligne--Mumford stacks. In this last statement $L$ means the actual
oriented-blow-up boundary stack; identifying it with an intrinsic
punctured-neighbourhood germ still requires the hypotheses of F.19.

\begin{proof}
At a depth-$a$ point, the resolved boundary has $a$ angular coordinates
and a radial corner with at least one coordinate zero. Its depth-$a$
locus has real codimension $a-1$ in the boundary. Round corners and
triangulate compatibly. The deleted strata for $L_{\leq r}$ therefore
have codimension at least $r$. Relative general position moves a map
of a compact polyhedron of dimension less than $r$ off them, fixing
any closed subpolyhedron already disjoint. Apply this first to
representatives of homotopy groups and then to nullhomotopies. It
gives precisely the asserted connectivity. Local finiteness suffices:
the image of each such map meets only finitely many charts.

For sharpness take $m=r+1$ coordinate hyperplanes in $\mathbb C^m$.
The radial boundary is contractible, whereas its complement of the
vertex retracts onto $\partial\Delta^{m-1}$; consequently
\[
 L\simeq(S^1)^{r+1},\qquad
 L_{\leq r}\simeq(S^1)^{r+1}\times S^{r-1}.
\]
For $r=1$ the components differ. For $r>1$ the factor $S^{r-1}$
changes $\pi_{r-1}$. Thus $r\leq k+1$ cannot determine every $k$-type.
Exact integral boundary-matrix ranks for the simplex boundaries of
dimensions zero through four were computed independently; the radial
argument proves the formula in all dimensions. For a stack apply the
manifold result to every term of an etale atlas nerve. Realization
preserves an $m$-connectivity bound: relative CW models have cells only
above degree $m$, and realization adds nonnegative simplicial degrees.
This proves the assertion with isotropy retained.
\end{proof}

\Needspace{5\baselineskip}
\paragraph{F.17 (Proved). Corollary 3
and the topoi that are truncated.}
\textit{Manuscript: MPL147; section11; folio75; scan152; (247).}

Depths one, two and three, with every flag between them, compute
$\Pi_1(L)$. If $L\to X^*=X\setminus D$ is a one-type equivalence,
formula (247) has target $\Pi_1(X^*)$. The unstarred $\Pi_1(X)$ visibly
written in scan 152 is false under that hypothesis. For the full
\emph{sheaf pro-topos} gluing of F.13 all nonempty depths must be
retained; no dimension-independent cutoff computes it. The depth
bound in F.16 applies to the associated shape, not to that topos.

\begin{proof}
The fundamental-groupoid functor is the reflection from spaces to
one-types and preserves homotopy colimits, interpreted as
bicategorical colimits of groupoids. Apply F.16 with $k=1$. The
counterexample to the unstarred formula is $(\mathbb C,\{0\})$:
$L\simeq\mathbb C^*\simeq S^1$, whereas $\mathbb C$ is contractible.
This identifies a mathematical correction, not an alternative reading
of the ink. For full topoi a deleted deeper stratum is detected by
sheaves supported on it; even its locale of opens has changed. Thus
arbitrarily high depths cannot be removed in arbitrary dimension.
Ordinary sheaf descent uses objects, arrows and the cocycle condition
on triple overlaps; higher simplicial degrees are then forced. This
fact about \emph{simplicial coherence degree} does not permit omission
of high-depth objects: a depth-three object already supplies the
missing triangular relations in F.15.
\end{proof}

\Needspace{5\baselineskip}
\paragraph{F.18 (Proved). Formula (254) in every dimension.}
\textit{Manuscript: MPL147; section12; folio74 as transcription; scan156; (254).}

With the germ conventions of F.12, put
$\mathcal V=\{\operatorname{Sh}(V)\}_{V\supset D}$ and
$\mathcal V^*=\{\operatorname{Sh}(V\setminus D)\}_{V\supset D}$.
Then in the pro-two-category of full Grothendieck topoi
\[
 \operatorname{Sh}(X)\simeq
 \operatorname{Sh}(X^*)\amalg_{\mathcal V^*}\mathcal V .
\]
For locally conical pairs the analogous gluing holds for sheaves of
spaces, and its shape is the homotopy pushout of the shapes of the
three terms. In particular
\[
 \Pi_1X\simeq\Pi_1X^*\amalg_{\Pi_1\mathcal V^*}\Pi_1\mathcal V.
\]
The two germ terms may be computed by the full flag diagrams of their
elementary pieces; for the punctured boundary one-type F.17 gives the
sharp depth-three replacement. For curves the closed term is the
disjoint union of the boundary points and the punctured term the
disjoint union of their meridian circles. This is the geometric form
of STRATA C.2 and C.5, not merely its presheaf special case.

\begin{proof}
For every neighbourhood $V$, $X^*$ and $V$ form an open cover of $X$.
Their sheaves are glued by identifying their restrictions to
$V\setminus D$, with the ordinary cocycle condition; F.13 verifies
this universally for geometric morphisms and passes to cofinal germs.
For space-valued sheaves use the entire Cech nerve: descent identifies
mapping spaces, not just their components. Applying this to a constant
space $K$ identifies maps from the shape with the homotopy limit of
maps from the three shapes. This is exactly the homotopy-pushout
universal property. Passing to one-types gives the displayed groupoid
formula. The tube retraction to a point and the punctured-disc
retraction to a circle give the curve case directly, without an
asphericity assumption on $X$. The intrinsic two-open argument also
applies to actual stack neighbourhoods. Computing those germs through
an atlas requires the proper or common-stage controlled class of F.13;
F.51 shows why the unrestricted atlas-germ substitution is invalid.
\end{proof}

\subsection{Normal directions, tangential points and arithmetic tubes}

\Needspace{5\baselineskip}
\paragraph{F.19 (Proved). The canonical functor (243).}
\textit{Manuscript: MPL147; section11; folios68-73; scans139-148; (231),(234),(239),(243).}

Let $Y$ be a paracompact complex manifold with normal-crossings divisor
$D$, and let $D'\to Y$ be a normalized locally immersed intersection
of branches. Put $S=D'^*$, its open stratum, and let
$P_S=\prod_{j\in J,S}N_j^\times$ be the product of the punctured
normal line bundles. There is an equivalence
\[
 \Pi_1^{\mathrm{arcs}}(P_S)\longrightarrow
 \Pi_1^{\mathrm{rev}}(\mathcal V^*_{D',Y}),
\]
canonical as a functor with its specified natural comparison
isomorphisms. It is compatible with the generizations (231) and with
(234), including the coherence for every longer flag. No tubular
embedding is part of its definition. The intrinsic covering-germ
equivalence descends to the proper or common-stage controlled analytic
etale stacks of F.13. For an arbitrary etale presentation the sector
construction still defines a canonical coherent functor on its
chartwise germs, but these need not compute the intrinsic stack-open
neighbourhood germ. F.51 gives a parabolic quotient for which even
the latter covering equivalence fails. The quotient
$P_S/(\mathbb R_{>0})^J$ is the
open face over $S$ of the real oriented blow-up of $(Y,D)$.

\begin{proof}
Define the oriented blow-up intrinsically: a point over $y$ assigns
a unit complex argument to each local defining equation of a branch,
subject to $\arg(uf)=u(y)/|u(y)|\,\arg(f)$. A normal vector has the
same transformation rule. Its modulus is forgotten by the quotient
by positive real scalars; this gives the asserted face identification
without coordinates or a metric.

Here is the functor on coverings, which makes canonicity explicit.
Given a covering of a sufficiently small deleted neighbourhood and
a normal direction $v\in P_S$, choose a simply connected angular
sector about its arguments and a sufficiently small base neighbourhood
in $S$. The fibre at $v$ is the set of germs of sections on that
sector with all normal moduli positive and sufficiently small.
Shrinking sectors and radii does not change this set: in a normal
crossings chart these regions are contractible, and each covering is
a disjoint union of such regions. These sets form a covering of
$P_S$. A path in $P_S$ is covered by finitely many sectors; overlap
identification transports the germ of a section. Subdivision of the
path and homotopies give the same transport by uniqueness of lifting.
Thus evaluation gives a fibre functor on covering germs at every
$v$, and transport gives exactly (243).

In a chart $\Delta^a\times(\Delta^*)^{|J|}$ this operation identifies
both covering categories with $\mathbb Z^{|J|}$-sets, with the same
coordinate monodromies. It is fully faithful and essentially
surjective there. Coverings and their isomorphisms satisfy descent,
so these equivalences glue along all chart intersections. If $D'$
meets further branches, apply the same construction with their
coordinates also punctured. Pushing their positive radial coordinates
slightly into the interior shows that its covering-germ category is
unchanged upon replacing $D'$ by $S$; in each chart both categories
have exactly the same punctured coordinates and monodromy generators.
This proves (239) in the stated, not merely smooth-divisor, case.

For a nested flag use sectors with all its branch coordinates at once.
Successive restriction of a section germ is literal restriction, so
both orders agree; this proves (231), (234) and all longer
compatibilities before taking fundamental groupoids. On an etale
atlas all these constructions pull back and retain the isotropy
coherence. To identify the result with coverings of intrinsic stack
neighbourhoods, use the synchronized descent of F.13 in its stated
proper or controlled range. Without that hypothesis this argument
constructs the atlas-germ functor only; it does not silently exchange
an arbitrary infinite descent diagram with shrinking neighbourhoods.
In dimension one the construction evaluates a cover on a germ of a
sector tangent to a nonzero vector. The finite algebraic version is
Deligne's specialization to the punctured tangent line
\cite[\S15.13--15.27]{Deligne1989}. In a parameter $t$, it selects
compatible roots of $t$ and the fibre over the specified tangent
vector; changing $t$ by $u t$ multiplies the root coordinate by the
corresponding roots of $u(0)$. This is precisely the sector transition
rule. Hence the analytic, profinite, and tangential-point functors
coincide, with a path torsor rather than a falsely chosen canonical
basepoint identification. The specified cote~140-4 pages of
\emph{Longue Marche} were subsequently recovered inside the LM8
background archive; F.48 supplies the literal manuscript cross-check.
\end{proof}

\Needspace{5\baselineskip}
\paragraph{F.20 (Refuted). An invariant tubular
embedding need not exist.}
\textit{Manuscript: MPL147; section11; folio72; scan146.}

The matrix in MPL~147, fol.~72, scan~146 must have upper-right entry
$1$ for the example to work. Let the discrete group $\mathbb Z$ act
on $\mathbb C^2$ by $(x,y)\mapsto(x+y,y)$ and let $D=\{y=0\}$.
There is no $\mathbb Z$-equivariant differentiable tubular embedding
from the normal line bundle to $\mathbb C^2$ which induces the
identity on its normal quotient. This is an example for an analytic
etale quotient stack with infinite isotropy, not for a proper
finite-stabilizer orbifold.

\begin{proof}
The action on $D$ and on its normal line is trivial. Differentiating
such an embedding at $(0,0)$ would give an equivariant splitting of
\[
 0\longrightarrow\mathbb C(1,0)\longrightarrow\mathbb C^2
       \longrightarrow\mathbb C\longrightarrow0.
\]
A fixed vector $(a,b)$ satisfies $(a+b,b)=(a,b)$, hence $b=0$;
none projects to $1$. This rules out even a real differentiable
splitting. The attached exact matrix calculation verifies this
obstruction. Entry $0$ would give the identity matrix and a splitting,
so it would not be a counterexample. For a finite group, averaging
an invariant metric supplies an equivariant smooth tubular
neighbourhood; the infinite stabilizer is material. The sector
construction in F.19 needs no such splitting and therefore survives.
\end{proof}

\Needspace{5\baselineskip}
\paragraph{F.21 (Proved). The arithmetic
normal tube and (37).}
\textit{Manuscript: MPL149; section4; folios65v-67; scans132-135; (36)-(38).}

Let $(X,D)$ be a smooth normal-crossings pair over a characteristic-zero
field $K$, with $X$ of finite type, and let $S=D_J^*$ be a connected
open stratum on the normalization that labels its $r=|J|$ branches.
Define its tame etale tube by the finite-covering category of the
punctured henselian neighbourhood along $S$; equivalently use the
Kummer etale normal tube. This is independent of the chosen etale
charts. Its geometric profinite shape is that of
$P_S=\prod_{j\in J,S}N_j^\times$. With geometric basepoints and
paths specified, its fundamental groups have exact sequence
\[
 \pi_2^{\mathrm{et}}(S_{\bar K})
 \xrightarrow{\partial}
 T^J=\widehat{\mathbb Z}(1)^J
 \longrightarrow\pi_1^{\mathrm{et}}((\mathcal V_S^*)_{\bar K})
 \longrightarrow\pi_1^{\mathrm{et}}(S_{\bar K})\longrightarrow1.
\]
Here $\pi_2^{\mathrm{et}}$ is the profinite homotopy group of the
geometric etale shape, not the abstract completion of an unspecified
ordinary group. The map $\partial$ is the tuple of normal-line Chern
transgressions. Thus (37) is exact with $D_J$ meaning the open stratum;
its initial $1$ is valid exactly when $\partial=0$. Injectivity of
$T^J\to\pi_1^{\mathrm{et}}(U_{\bar K})$ implies that vanishing.
Arithmetic groups fit into the additional extensions over $G_K$;
Galois acts cyclotomically on each labelled factor of $T^J$ and
permutes factors if branch labels have not been split.

\begin{proof}
Locally along $S$ choose parameters $t_1,\ldots,t_r$ for its branches.
The henselian normal comparison sends a tame covering to its normal
specialization. A finite covering is trivialized locally by adjoining
roots $t_j^{1/n}$: normalization after these adjunctions is unramified,
by the valuation form of Abhyankar's lemma. On overlapping charts
$t'_j=u_jt_{\sigma(j)}$, the transition is the Kummer torsor of the
unit $u_j|_S$. These are exactly the transitions of $N_j^\times$.
This gives mutually inverse functors on finite covers and on their
isomorphisms, with the same descent cocycles; it proves the assertion
about the tube and its functor to $U$.

For the higher exact sequence, the root gerbes of the normal lines
provide a precise profinite model. For each $n$, the product of gerbes
of $n$th roots of the $N_j$ has fibre $B\mu_n^J$ over $S$ and is
classified by the Kummer Chern classes $c_1(N_j)\bmod n$. Their
compatible inverse system has fibre $B\widehat{\mathbb Z}(1)^J$.
The preceding local calculation, now for finite locally constant
truncated spaces by descent and their finite Postnikov towers,
identifies it with the profinite normal-tube shape. Equivalently this
is the local normal-bundle instance of the log comparison theorem
\cite[Theorem~A and Corollary~5.5, preprint numbering]{CSST2019}.
The homotopy exact sequence of this fibre sequence gives the displayed
sequence, and the boundary maps of the root gerbes are precisely the
Chern transgressions. Over $\mathbb C$ their evaluation on a sphere
is the integer degree of the pulled-back normal line bundle; reduction
modulo every $n$ gives the profinite statement. Descent from a finitely
generated field of definition and geometric extension of constants
preserves the finite-covering and finite-coefficient data, so the same
argument applies over arbitrary characteristic-zero $K$.

The composite $T^J\to\pi_1\mathcal V_S^*\to\pi_1U$ kills
$\operatorname{im}\partial$. If this composite is injective then
$\partial=0$, proving the requested direction. It need not vanish:
for the zero section in $\operatorname{Tot}\mathcal O_{\mathbb P^1}(1)$,
$\partial$ is an isomorphism and the punctured bundle is simply
connected; for a smooth conic in $\mathbb P^2$, its normal bundle has
degree four and the tube has fundamental group $\mathbb Z/4$, with
kernel $4\widehat{\mathbb Z}$ in its meridian group. Finally take
$X=\mathbb A^2$, $D=\{xy=0\}$ and the branch $D_1=\{x=0\}$.
The punctured tube has group $\widehat{\mathbb Z}^2$, its normal
meridian has rank one, and $D_1^*=\mathbb G_m$ supplies the other
factor. Replacing this open stratum by the closed $\mathbb A^1$
would make the last group trivial and break exactness. This also
confirms the mathematical force of (36) and the definition of $U_J$
as the trace of $U$ on the localized neighbourhood in scans~132--133.
\end{proof}

\Needspace{5\baselineskip}
\paragraph{F.22 (Refuted). Finite roots do
not have a constant inertia kernel.}
\textit{Manuscript: MPL149; section4; folios66-67; scans133-135; EXIT task5.}

Write $G=\pi_1^{\mathrm{tame}}(U_{\bar K})$ and
$\phi:T^J\to G$. The Kummer etale site of $(X,D)$ and the infinite
root stack have this same group and hence the same kernel
$\ker\phi$. For a finite root stack with orders $n_i$ along the
boundary components, put $N_{\mathbf n}$ equal to the closed normal
subgroup of $G$ generated by all $n_i$th powers of its corresponding
boundary inertia groups. Then
\[
 \pi_1(\sqrt[\mathbf n]{(X,D)})=G/N_{\mathbf n},\qquad
 \ker(T^J\to\pi_1\sqrt[\mathbf n]{(X,D)})
       =\phi^{-1}(N_{\mathbf n}).
\]
These kernels can be strictly larger than $\ker\phi$. Their
intersection over all divisible root orders is $\ker\phi$.
A finite Kummer \emph{cover} likewise cannot remove a nonzero kernel,
but this is a different assertion from equality for every finite
root stack.

\begin{proof}
A finite etale cover of the finite root stack is a tame cover of $U$
whose inertia orders divide the specified $n_i$. The assertion is
checked in a root chart $[\operatorname{Spec}A[z_i]/(z_i^{n_i}-t_i)
/\prod\mu_{n_i}]$: unramifiedness on the chart is precisely that
inertia bound; the charts and their covers descend. Therefore its
finite quotient actions are exactly those killing
$N_{\mathbf n}$, proving the displayed formulas. Every finite quotient
of $G$ kills $N_{\mathbf n}$ for some common multiple of its finite
inertia orders; thus $\bigcap_{\mathbf n}N_{\mathbf n}=1$.
All tame covers extend on the Kummer site by the same local
normalization argument, proving the infinite-level assertion.

For an explicit counterexample take $(\mathbb A^1,\{0\})$.
Here $G=T=\widehat{\mathbb Z}(1)$ and $\phi$ is the identity, while
$\sqrt[n]{(\mathbb A^1,0)}=[\mathbb A^1/\mu_n]$ has group $\mu_n$.
The kernel has changed from zero to $n\widehat{\mathbb Z}(1)$.
Finite quotients of this calculation were checked by the attached
program; the Kummer classification proves it in the inverse limit.

For a Kummer cover, the local lattice is a finite-index sublattice
$T'\hookrightarrow T$, and its global group embeds as an open subgroup
of $G$. Its kernel is the inverse image of $\ker\phi\cap T'$.
For any nonzero closed subgroup $A\subset\widehat{\mathbb Z}^{r}$,
and any finite-index sublattice $T'$, choose a positive integer $m$
with $mT\subset T'$. Multiplication by $m$ is injective on $T$, so
$0\ne mA\subset A\cap T'$. Thus the kernel persists. This proves
the valid content of EXIT\_ONE\_TYPE P1e-roots and refutes the stronger
finite-root-stack sentence added in the present prompt.
\end{proof}

\Needspace{5\baselineskip}
\paragraph{F.23 (Proved). The stable cusp quotient changes the model.}
\textit{Manuscript: MPL149; section4; folios66-67; scans133-135; supplemental EXIT P1e-cusp-stack test, not a manuscript assertion.}

Over a characteristic-zero field containing $\mu_6$, put
$F=\{Y^2=X^3+1\}$ and
\[
 \mathcal X=[F\times\mathbb A^1_s/\mu_6],\qquad
 \zeta(X,Y,s)=(\zeta^{-2}X,\zeta^{-3}Y,\zeta s).
\]
Its boundary is $[F/\mu_6]$ at $s=0$; its open part is the family
$y^2=x^3+t$, $t\ne0$, via $t=s^6,x=s^2X,y=s^3Y$.
The normal meridian injects both discretely and profinitely into
the open fundamental group. Here groups are geometric; the discrete
calculation uses the complex realization of this model defined over
$\mathbb Q(\mu_6)$. The same holds for the full rank-two corner when
the infinity section in the projective elliptic completion is retained
as boundary. The original embedded cusp resolution has the exact
nonzero node kernels $\mathbb Z(-3,1)$ and $\mathbb Z(-2,1)$,
and their profinite counterparts, computed below. Thus the stable
quotient removes explicit local meridian kernels by changing the
model. It is not a finite root stack of that fixed resolved divisor.

\begin{proof}
The displayed substitutions verify the equation and the action is
free when $s\ne0$, so the asserted identification of the open part
is an isomorphism. A loop in the normal coordinate $s$ has degree six
under $t=s^6$. Its composite with the degree homomorphism to
$\pi_1(\mathbb G_m)=\mathbb Z$ is multiplication by six, which is
injective; the same calculation on $\widehat{\mathbb Z}(1)$ proves
profinite injectivity. No residual-finiteness inference from the
discrete injection is being used.

The affine curve $F$ is a once-punctured elliptic curve; its completion
is elliptic. Both are aspherical, and their finite quotient stacks
are aspherical by the fibration over $B\mu_6$. In the affine case
the coarse boundary coordinate $w=Y^2=X^3+1$ has cone orders two
and three, giving group $C_2*C_3$; the projective case adds the
order-six point at infinity. Thus the quotient retains finite
monodromy in boundary isotropy while retaining the infinite normal
meridian in the open part. F.22 shows why merely extracting finite
roots of the old boundary could not eliminate an existing kernel.
Here is the explicit kernel in the old resolved model, so that the
comparison with the stable quotient does not rely on the supplied
EXIT calculation. In the affine plane put $f=y^2-x^3$. Three successive
blow-up charts give
\[
 f=x^2(v^2-x),\qquad
 f=u^2v^3(v-u),\qquad
 f=u^6w^3(w-1).
 \tag{F23.1}
\]
The substitutions are respectively $y=xv$, $x=uv$, and $v=uw$.
In the second substitution the preceding coordinate $v$ is retained;
in the third the preceding $u$ is retained. After the third blow-up,
the strict cusp and the first two exceptional components meet the
last exceptional component $E_3$ at three distinct points. The first
two exceptional components $E_1,E_2$ each have only that one boundary
intersection. Their self-intersections are $-3,-2$, since $E_1$ is
blown up twice after its creation and $E_2$ once. The last component
has square $-1$. The multiplicities of $f$ on $E_1,E_2,E_3$ are
$2,3,6$.

Write $b,a,h$ for the positive normal meridians of these three
components. The plumbing relation at the two exceptional tails is
\[
                         h=b^3,\qquad h=a^2.              \tag{F23.2}
\]
To verify the signs and the relation, consider a rational component
of square $-d$ meeting the remaining divisor only at $q$.
Its normal line is $\mathcal O_{\mathbb P^1}(-d)$. With $u$ a
coordinate at $q$ and $z_\infty$ its normal coordinate there, a
trivialization over $\mathbb P^1\setminus\{q\}\simeq\mathbb A^1$
has normal coordinate $z_0=u^d z_\infty$. Thus the positive loop in
$u$ with $z_\infty$ fixed winds $d$ times in $z_0$. Its base loop
contracts in $\mathbb A^1$, so in the punctured normal tube it is
the $d$th power of the component's normal meridian. The ordinary
smooth tubular-neighbourhood theorem identifies these maps with
the local plumbing maps into the divisor complement. Taking
$d=3,2$ gives (F23.2); no equivariant tubular embedding is used.

The map $f:U\to\mathbb G_m$ sends $b,a,h$ to $2,3,6$ in
$\pi_1(\mathbb C^*)=\mathbb Z$. Hence $a$ and $b$ have infinite
order, and the two full node maps have the exact kernels
\[
 \begin{split}
 \ker\bigl(\mathbb Z^2\longrightarrow\pi_1U,
                      (r,s)\longmapsto b^r h^s\bigr)
       &=\mathbb Z(-3,1),\\
 \ker\bigl(\mathbb Z^2\longrightarrow\pi_1U,
                      (r,s)\longmapsto a^r h^s\bigr)
       &=\mathbb Z(-2,1).
 \end{split}                                             \tag{F23.3}
\]
Indeed (F23.2) gives one inclusion, and composition with $f_*$
gives the reverse one. The same argument is exact after profinite
completion: multiplication by $2$ or $3$ is injective on
$\widehat{\mathbb Z}$, so the kernels are respectively
$\widehat{\mathbb Z}(-3,1)$ and
$\widehat{\mathbb Z}(-2,1)$. Thus the kernel that the stable
quotient removes is explicit, even without computing the entire
trefoil group. The polynomial identities, self-intersection updates,
and finite congruence certificates for these kernel equalities are
executed in \path{code/montpellier/F23_cusp_certificates.py}, with its output
attached.

There are two different uses of the projective completion.
If only its vertical divisor is retained, the preceding
multiplication-by-six argument proves the asserted rank-one
injection. If the infinity section is also retained as boundary,
the finite atlas is
$(\overline F\times\mathbb A^1,
  \overline F\times\{0\}\cup\{\infty\}\times\mathbb A^1)$,
with open part $F\times\mathbb G_m$. Its rank-two corner map is
\[
 \mathbb Z^2\longrightarrow \pi_1(F)\times\mathbb Z,
                  (r,s)\longmapsto (\delta^r,s),           \tag{F23.4}
\]
where $\pi_1(F)$ is free on two generators and $\delta$ is their
commutator, the meridian of the missing point. A reduced commutator
has infinite order in this free group, so (F23.4) is injective.
Its profinite cyclic factor is injective as well: for every $n$
send the two free generators to the elementary generators of the
finite upper-unitriangular Heisenberg group modulo $n$; their
commutator has exact order $n$. Consequently the map
$\widehat{\mathbb Z}\to\widehat{\pi_1(F)}$ generated by $\delta$
has zero kernel. The product lattice therefore injects profinitely.
The finite atlas is etale over the quotient stack and over its open
part, so injection into the open atlas group remains injection into
the open quotient group. This proves the full corner assertion for
the marked projective model, rather than identifying a pointed
completion with its unmarked version.


The construction uses inversion of $2$ and $3$; nothing here asserts
its validity for wild reduction at those primes.
\end{proof}

\Needspace{5\baselineskip}
\paragraph{F.24 (Proved). The tame
constructible topos and its descriptive system.}
\textit{Manuscript: MPL147; section12; folios85-85v; scans179-180; MPL149 section4 folio67 scan135.}

For a finite normal-crossings stratification over a characteristic-zero
field, finite etale sheaves constructible on its strata are finite
continuous representations of the profinite completion of its ordinary
exit category. The Ind-category of these finite representations is the
corresponding coherent constructible topos. Its stratum groups are the
geometric etale fundamental groups, with arithmetic descent over $G_K$;
its crossing sets and compositions are the profinite versions of the
coset and flag constructions in F.10. The groupoids of punctured
normal tubes are exactly the extensions in F.21, and the maps obtained
by forgetting branch flags are (38). These identify the claimed
``1-descriptif'' system without suppressing its kernels or its
commutation isomorphisms.

The arbitrary-product condition in F.8 is a characterization for
\emph{discrete} groupoid presheaf topoi. It is false as a necessary
condition for this profinite variant.

\begin{proof}
Over $\mathbb C$, finite constructible sheaves are recovered by their
finite stalks, path transport on each stratum, and the sector
specialization maps of F.19. The inverse construction glues constant
finite sheaves on the sectors; compatibility on double intersections
and on triple flags gives ordinary sheaf descent. Thus it is exactly
the set-valued case of F.10 with finite fibres. Every such finite
representation factors through a finite quotient of the category:
for each pair of the finitely many stratum basepoints, identify arrows
with the same map on all the specified finite fibres, and identify
endomorphisms with the same permutations. Composition descends because
these are a congruence. Taking all finite representations and their
finite products gives the cofiltered system of finite quotient
categories and proves the continuous-representation assertion.
The same construction for a finite family of representations proves
that the category is a pretopos, whose Ind-category is the asserted
coherent topos.

Analytification identifies the algebraic finite constructible sheaves
with these analytic ones, including their gluing functors, by
\cite[Proposition~12.6.4 and Corollary~12.6.7]{FA_BGH2020}. More
explicitly, comparison holds on each smooth stratum for finite local
systems, and on its punctured normal charts by the Kummer calculation
of F.21; induction over the finite closed filtration identifies the
recollement triples, so it holds on the whole pair. Descent from a
field of definition and finite separable extension of constants gives
the arithmetic groupoids and their weak $G_K$-action as in F.1.
The descriptions of crossing sets and compositions follow by applying
finite stalk functors to the same sector and flag constructions.
In positive characteristic the same assertion concerns tame sheaves
and prime-to-characteristic local normal inertia; no wild cover is
silently included.

To see why products must be treated differently, let $G$ be a
profinite group and $H\subset G$ a closed subgroup which is not open.
In continuous discrete $G$-sets the product of a family is the subset
of the set-theoretic product whose points have open stabilizer. Choose
open subgroups $G_\alpha\supset H$ with intersection $H$. Each coset
$G_\alpha$ in $G/G_\alpha$ is fixed by $H$. Their tuple is not in the
continuous product, since its stabilizer is $H$. Therefore
\[
 \left(\prod_{\alpha}^{\mathrm{cont}}G/G_\alpha\right)^H
 \longrightarrow\prod_\alpha(G/G_\alpha)^H
\]
is not surjective. This is the actual cusp gluing functor for a
hyperbolic curve whenever its inertia subgroup has infinite index,
for example the thrice-punctured projective line. Nevertheless its
finite continuous exit description has just been proved. Finite
representations and coherent descent, rather than an unchanged
arbitrary-product test, give the valid profinite statement.
\end{proof}

\Needspace{5\baselineskip}
\paragraph{F.25 (Refuted). Components of the analytic flag category
need not be a category.}
\textit{Manuscript: MPL147; section11; folios56-59; scans114-121; EXIT P4a1-P4a2.}

Even when the analytic incidence spaces satisfy all Segal identities,
their sets of connected components need not do so. The correct object
is the semisimplicial set of components of \emph{all} strict flags,
with repeated ranks adjoined if a simplicial object is desired.

\begin{proof}
Let $\tau=(12)(34)$ act freely on $\mathbb C^4\times\mathbb C^*$
by permuting the coordinates and sending $w$ to $-w$. Coordinate
hyperplanes descend to a normal-crossings divisor in a smooth
quotient. Every coordinate intersection times $\mathbb C^*$ is
connected, so connected flag components are the $\tau$-orbits of
branch-set flags. Put $B=\{1,2\}$. The upper pairs $(123,B)$ and
$(124,B)$ are in the same orbit, and the lower pairs $(B,1)$ and
$(B,2)$ are in the same orbit. However $(123,B,1)$ and $(123,B,2)$
are different orbits: applying $\tau$ changes $123$ to $124$.
Their composites $(123,1)$ and $(123,2)$ are also different orbits.
Thus a composable pair of component arrows has two possible
composites. The attached orbit enumeration checks these claims.
Forgetting entries of a complete flag always gives well-defined
face maps and satisfies the semisimplicial identities before and
after taking components; that is the stated replacement. Locally
its realization is the barycentric subdivision of the simplex of
branches, so its global realization is the usual coarse dual complex.
The normalized flag spaces themselves, with stabilizers and their
full fibre products, remain the correct geometric descent data.
\end{proof}

\subsection{Elementary topoi and the two shapes}

The intrinsic neighbourhood-germ statements in F.26--F.30 concern
analytic spaces and the proper or common-stage controlled stack
presentations specified in F.13 and F.19. The abstract exit-category
and bar-construction statements are separate from that geometric
identification. No assertion here replaces intrinsic germs by an
uncontrolled nonproper atlas, as ruled out by F.51.

\Needspace{5\baselineskip}
\paragraph{F.26 (Proved). The four elementary pieces and their products.}
\textit{Manuscript: MPL147; section12; folios78-81; scans163-169; (257),(264)-(277).}

Over $(\mathbb C,\{0\})$ write
$P=\operatorname{Sh}(\{0\})$, $A=\operatorname{Sh}(\mathbb C^*)$,
$D=\{\operatorname{Sh}(\Delta_\epsilon)\}_\epsilon$ and
$D^*=\{\operatorname{Sh}(\Delta_\epsilon^*)\}_\epsilon$.
The canonical incidence maps \emph{over the ambient line} are
$P\to D$, $D^*\to D$, $D^*\to A$, their identities and composites
with the ambient inclusions. These generate exactly diagram (265).
There is no incidence $P\to A$ over the line. Structural maps to
the terminal topos exist, but they are not these ambient incidences.
For $I$ finite the elementary pieces are
\[
 E_{I_P,I_A,I_D,I_*}
 =P^{I_P}\times A^{I_A}\times D^{I_D}\times(D^*)^{I_*},
 \qquad I=I_P\amalg I_A\amalg I_D\amalg I_*.
\]
Their canonical incidence maps are the coordinatewise products of
the preceding three maps. For $|I|=2$ the complete list is
\[
\begin{array}{c|cccc}
 &P&A&D&D^*\\ \hline
 P&P P&P A&P D&P D^*\\
 A&A P&A A&A D&A D^*\\
 D&D P&D A&D D&D D^*\\
 D^*&D^*P&D^*A&D^*D&D^*D^*
\end{array}
\]
where juxtaposition denotes the product, not an identification.
The shape of such a piece is $(S^1)^{|I_A|+|I_*|}$, although its
full sheaf pro-topos is generally not the classifying topos of that
group.

Every canonical deleted germ (257), defined as the neighbourhood
germ of $\Sigma'$ in $\Sigma''$ followed by deletion of its trace
of $\Sigma$, is the bicategorical colimit of its contained elementary
pieces with their full incidence flags. This assertion uses full
sheaf pro-topoi. For shapes it is a homotopy colimit, and not the
shape of an arbitrary ordinary two-colimit of covering topoi.
For stack presentations this intrinsic-germ assertion has the proper
or common-stage controlled scope of F.13 and F.19.

\begin{proof}
The list is the independent choice of one of four states for each
coordinate: fixed zero, globally nonzero, near zero, and near zero
but nonzero. An inclusion over the ambient line can only fix zero
inside its germ, or include a punctured germ into that germ or the
nonzero line. This proves the generators and their coordinatewise
product description; it also explains why an apparent arrow from
$P$ to $\mathbb C^*$ in the typographical list (264) cannot be an
ambient inclusion. Polar retraction gives the displayed shapes.

At each positive radius the cover
$\mathbb C=\mathbb C^*\cup\Delta_\epsilon$ has intersection
$\Delta_\epsilon^*$. Its products form a finite cubical open cover,
and their intersections are exactly the stated punctured and
unpunctured coordinate pieces. Restriction to a coordinate closed
subspace replaces its zero coordinates by $P$. Restriction to an
admissible union, deletion of another admissible union, and taking
neighbourhoods of a third are computed in the same coordinates;
finite unions of tubes and their intersections give the corresponding
subdiagram. Apply full sheaf descent, including its triple-overlap
relations, and pass to the cofinal system of positive radius
functions. This is F.13 applied to this explicit cover, and proves
the assertion for (257). Applying space-valued descent gives the
homotopy-colimit assertion by the mapping-space argument of F.18.
On non-simple crossings in a Hausdorff analytic space use normalized
branch charts and their full etale nerve; their fibre products supply
the flag incidences and descend the same calculation. A stack atlas
additionally requires the synchronization hypothesis in F.13; no
unrestricted nonproper atlas-germ replacement is used. The order of taking a germ and deleting
its support matters: the punctured germ of a nonempty stratum is
not the germ of the empty stratum, exactly as scan~158 observes.
\end{proof}

\Needspace{5\baselineskip}
\paragraph{F.27 (Proved). The Scholie and (295).}
\textit{Manuscript: MPL147; section12; folio84; scans177-178; (295).}

For a general flag $d_*$ and two of its ranks $d'_0\geq d''_0$,
the elementary germ $\mathcal V^{d'_0}_{d_*,d''_0}$ has the
homotopy type of the punctured graded normal bundle over
$D_{d_*}^*$ with fibre $(\mathbb C^*)^{d'_0-d''_0}$.
Its fundamental groupoid is canonically equivalent, with the sector
comparison data, to that of this punctured bundle. This is (295).
Normal-line twisting and branch monodromy must be retained.

\begin{proof}
The two nested normal subbundles give an exact sequence
$0\to N'\to N''\to Q\to0$ as in (287). Remove the coordinate
hyperplanes from $Q$ and take its inverse image in $N''$.
The resulting map to $Q^\times$ is an affine bundle under the
pullback of $N'$; its fibres are contractible. Local trivializations
and a numerable cover show that this map is a homotopy equivalence,
without choosing a global splitting. Restrict to $D_{d_*}^*$.
The sector construction of F.19 compares its covering category with
the elementary germ and is compatible with refinement of the flag.
It therefore supplies the canonical equivalence on fundamental
groupoids claimed in (295). The homotopy equivalence need not be
a specified global analytic map, consistent with the warning at
(292). A line bundle of nonzero Chern class already shows why one
must keep the torus bundle rather than replace it by a product with
the base. For a stack the normal sequence and its punctured quotient
are formed on an atlas and descend; no invariant splitting is used.
That vector-bundle descent is distinct from identifying the result
with an intrinsic neighbourhood germ, for which the proper or
common-stage controlled hypotheses of F.19 remain necessary.
\end{proof}

\Needspace{5\baselineskip}
\paragraph{F.28 (Proved). Exactly what the Remarque recovers.}
\textit{Manuscript: MPL147; section12; folios85-86; scans179-180; Remarque andRemords.}

Let $\mathcal E=\operatorname{Exit}(X,D)$ be the full exit
infinity-category, with exits directed from deeper to more open
strata. The system of fundamental groupoids of the elementary
pieces, with every flag and its comparison isomorphisms, always
recovers $\Pi_1X$ and the topos of constructible \emph{sets},
$\operatorname{Fun}(h\mathcal E,\mathbf{Set})$.
If all open incidence strata are aspherical, the homotopy colimit
of the \emph{classifying spaces} of that same coherent system
recovers the homotopy type of $X$ and of every canonical germ.
Finally $\mathcal E$ is equivalent to an ordinary category if and
only if all these strata are aspherical and every full local
meridian lattice injects into $\pi_1(U)$. Under both hypotheses the
shape of the ordinary constructible topos is also the correct shape.
Asphericity alone does not imply this last conclusion.

\begin{proof}
The first groupoid assertion follows from F.18 and the preservation
of homotopy colimits by $\Pi_1$. For constructible sets, F.10 obtains
crossing Hom-sets as the components of the endpoint homotopy fibres;
only the fundamental groups of links and endpoints enter these
components. Full flag spaces give composition. This reconstructs
$h\mathcal E$ and its representation topos.

By F.27 every elementary link is a torus bundle over an incidence
cover of an open stratum, with possible additional contractible
factors. A finite cover of an aspherical space is aspherical, and
the long exact homotopy sequence of a torus bundle shows that its
total space is aspherical. Thus the natural maps from all diagram
values to their $B\Pi_1$ are equivalences, coherently with every
map and homotopy. F.26 now proves the second assertion.

For the criterion, endomorphism spaces in $\mathcal E$ are loop
spaces of strata; they are homotopy-discrete exactly when the strata
are aspherical. For incident $S,T$, choose its incidence cover
$F_{ST}\to S$ and normal link $L_{ST}$ with torus fibre. Normal
sectors give
\[
 \operatorname{Map}_{\mathcal E}(s,t)
 \simeq\operatorname{hofib}_{(s,t)}(L_{ST}\to S\times T).
\]
If the strata are aspherical, so is $L_{ST}$. A component of this
homotopy fibre is the classifying space of the kernel of
$\pi_1L_{ST}\to\pi_1S\times\pi_1T$. The incidence cover injects
on $\pi_1$, so this kernel is precisely that of the relative normal
lattice in $\pi_1T$. Hence all mapping spaces are zero-types exactly
when all relative meridian maps are injective.

It remains to justify that full injection suffices for every relative
map. Partition the local branches as $I=J\amalg K$, with the $J$
branches remaining zero on $T$. Push the full torus into the normal
torus bundle over $T$. If $v\in\mathbb Z^K$ dies in $\pi_1T$,
exactness lifts it to the image of some $w\in\mathbb Z^J$ in that
normal bundle. Then $v-w$ dies in $\pi_1U$. Full injection forces
$v-w=0$, and disjoint coordinate supports force $v=w=0$. This step
does not presuppose asphericity or injection of the intermediate
normal lattice. Taking $T=U$ proves the converse. It proves the
criterion, including all incidence covers and isotropy.
\end{proof}

\Needspace{5\baselineskip}
\paragraph{F.29 (Proved). Shapes of the two constructible topoi and
the first obstruction.}
\textit{Manuscript: MPL147; section12; folios85-86; scans179-180; EXIT P1d3.}

The genuine constructible infinity-topos has shape
$B\mathcal E\simeq X$; the infinity-topos associated with the
ordinary constructible-set topos has shape $B(h\mathcal E)$.
Their comparison
$q:X\to B(h\mathcal E)$ is two-connected. If every mapping-space
component is $(r-1)$-connected, $r\geq1$, it is $(r+1)$-connected.
For its homotopy fibre $F$, its first possible homotopy group is
\[
 \pi_{r+1}F\cong
 \operatorname{coker}\left[
 H_r(\widetilde V_2;\mathbb Z)
 \xrightarrow{d_0-d_1+d_2}H_r(\widetilde V_1;\mathbb Z)\right].
\]
Here $V_p$ is the bar construction of a well-based enriched model
of $\mathcal E$ with discrete objects, and tildes denote pullback of
the fundamental-groupoid local system of paths on the chosen
component of $B(h\mathcal E)$. The formula is equivariant for that
component's fundamental group.

\begin{proof}
In a presheaf infinity-topos on a category $C$, global sections of
the constant space $K$ are $\lim_CK=\operatorname{Map}(BC,K)$;
this proves the shape assertions. The exit category groupoidifies to
$X$ by conical sector descent, or by F.26 with its space-valued
coherence. Replace the exit category by a cofibrant enriched model
and put
\[
 V_p=\coprod_{x_0,\ldots,x_p}
 \prod_{i=1}^p\operatorname{Map}(x_{i-1},x_i),\qquad W_p=\pi_0V_p.
\]
The degree-zero map is an isomorphism. In positive simplicial degree,
$V_p\to W_p$ is $r$-connected. A relative cell first occurring there
has dimension at least $r+1$, and realization adds the positive
simplicial degree; the first relative cells thus have dimension
$r+2$. This proves $(r+1)$-connectivity, and $r=1$ always applies.

Pull back the universal path local system of $|W|$ to both bars;
faces transport the lift along the corresponding edge. Its relative
homology spectral sequence is
\[
 E^1_{p,j}=H_j(\widetilde W_p,\widetilde V_p;\mathbb Z)
 \Longrightarrow H_{p+j}(\widetilde{|W|},\widetilde{|V|};\mathbb Z).
\]
Column zero vanishes and all rows $j\leq r$ vanish. In the first
possible total degree $r+2$, the only position is $(1,r+1)$,
with $H_{r+1}(\widetilde W_p,\widetilde V_p)=H_r(\widetilde V_p)$.
Its incoming differential is the three-face map displayed above;
higher differentials have zero source or target. Relative Hurewicz
and the exact sequence of the homotopy fibre identify this cokernel
with $\pi_{r+1}F$. This supplies the topos version of
EXIT\_ONE\_TYPE P1d3 rather than merely referring to it.
\end{proof}

\Needspace{5\baselineskip}
\paragraph{F.30 (Proved). The gluing question
and loss of homotopy.}
\textit{Manuscript: MPL147; section12; folio85; scans179-180; closing question.}

For a two-stratum conical space with closed stratum $S$, open stratum
$T$ and link $L$, the genuine constructible shape is
$S\amalg_L^hT$. When $S,T,L$ are aspherical with groups $A,B,G$,
the ordinary constructible shape replaces $G$ in this pushout by
$\operatorname{im}(G\to A\times B)$. Therefore the two shapes can
differ. For $(\mathbb P^1,\infty)$ they are respectively $S^2$ and
a point. For $(\mathbb P^n,\mathbb P^{n-1})$ the comparison fibre
is $\mathbb P^n$; for the square of $(\mathbb P^1,\infty)$ it is
$S^2\times S^2$. The first nonzero comparison-fibre group can occur
in arbitrarily high degrees, even for smooth algebraic pairs.

For a presheaf topos on $C$ with a full sieve $A$ and complementary
full subcategory $B$, let $E$ be the category of crossing arrows
$a\to b$, with commuting squares as morphisms. Then its shape is
\[
 BC\simeq BA\amalg_{BE}^{h}BB.
\]
This is the unconditional categorical answer to the question at
scan~180. For full geometric topoi the link is an actual topos if a
specified collared link exists; otherwise the deleted-neighbourhood
pro-topos is the precise object, as in F.12--F.18.

\begin{proof}
The conical mapping cylinder of $L\to S$ is a neighbourhood of $S$;
open descent with $T$ gives the first pushout. In the aspherical
case the crossing space is the homotopy fibre of
$BG\to B(A\times B)$. Its components are cosets of the image,
and their homotopy quotient by $A\times B$ is the classifying space
of the image stabilizer. Passing to the ordinary exit category
keeps only these components, which proves the replacement rule.
For the projective-line pair both endpoints are contractible and
the link is $S^1$. The ordinary category is $[1]$ and has
contractible nerve. For the hyperplane pair the open and closed
strata have trivial fundamental groups, and the Hopf-link endpoint
fibre is connected; again $h\mathcal E=[1]$. Products of exit
categories are computed coordinatewise, giving $[1]\times[1]$
for the square. These prove the three fibre assertions.

For arbitrary high degree take the smooth affine quadric
$Q_m=\{\sum_{i=0}^{m}z_i^2=1\}$, $m\geq2$, and the pair
$(Q_m\times\mathbb C,Q_m\times\{0\})$. Writing $z=x+iy$
identifies $Q_m$ with $TS^m$ up to the harmless positive rescaling
$|x|^2=1+|y|^2$, $x\cdot y=0$; scaling $y$ to zero retracts it
onto $S^m$. Its ordinary fundamental group is trivial.
The second factor's ordinary exit category has contractible
classifying space, so $B(h\mathcal E)$ is contractible while the
genuine shape is $S^m$. The first discrepancy is therefore degree
$m$, as asserted.

For the categorical statement, a nondegenerate string of arrows
in $C$ crosses from the sieve $A$ to $B$ at most once. Divide the
nerve into strings lying in either side and strings with their
unique crossing marked. The latter form the simplicial prism
resolution of $NE$; the two prism ends are its source and target
maps. Realizing this decomposition is the double mapping cylinder,
which proves the pushout, including disconnected and empty cases.
This checks STRATA C.3 and iterates to C.5. Its space-valued
coefficient form is the lax pullback of restriction functors:
objects are $(F_A,F_B,F_B\to i^*j_*F_A)$, and mapping spaces are
homotopy pullbacks. Restriction and the units of the two adjunctions
construct inverse functors, objectwise, proving STRATA C.6.
For locally constant coefficients the gluing arrow is an equivalence;
functors out of the just-proved homotopy pushout give the ordinary
infinity-categorical pullback of STRATA C.8. One cannot first take
homotopy categories: for the arrows $S^0\to *$ and $*\to S^1$,
the mapping-space components are
$\operatorname{coker}(\mathbb Z\xrightarrow{(1,1)}\mathbb Z^2)
=\mathbb Z$, whereas the ordinary comma category of homotopy
categories has a single morphism. This directly checks STRATA C.7
and locates the lost higher coherence.
\end{proof}

\subsection{F(h): local stratifications, flags, and controlled links}

The topology used in this subsection is essential. A stratum has its
subspace topology in a global chart; a global leaf of a local
stratification has the topology obtained by gluing its local plaques.
A neighbourhood germ remains a pro-object of \emph{full} sheaf topoi.
Neither taking connected components nor taking covering categories is
part of that definition.

\Needspace{5\baselineskip}
\paragraph{F.31 (Proved). With a counterexample to the sheaf assertion:
the presheaf is separated, and its sheafification is the atlas of local
stratifications.}
\textit{Manuscript: MPL 147; section 13, folios 86-87; scans 181-184.}

Use conditions a)--d) of MPL~147, fol.~86, scans~181--182: a locally
finite partition into locally closed, connected, locally connected strata;
the frontier condition; and, for every stratum $S$ and $x\in\overline S$,
a fundamental system of neighbourhoods $N$ of $x$ for which $N\cap S$
is nonempty and connected. Restrictions split strata into connected
components. These restrictions define a separated presheaf
$\operatorname{Strat}_{\rm glob}$. It is not a sheaf, even if all the
charts satisfy e). A section of its associated sheaf is exactly an open
atlas of such partitions agreeing, after component splitting, on
overlaps. Its leaves and all its local flag spaces are intrinsic to that
section, rather than to a chosen atlas.

\begin{proof}
We first verify the restriction assertion, which is needed before asking
whether the presheaf is separated. Let $U\subset X$ be open. A component
$C$ of $S\cap U$ is open in $S\cap U$ because $S$ is locally connected,
and is therefore locally closed in $U$. For $x\in U$, only finitely many
original strata meet a sufficiently small neighbourhood. For each such
stratum whose closure contains $x$, choose $N_S\subset U$ with
$N_S\cap S$ connected. This set lies in just one component of $S\cap U$.
Intersect the finitely many $N_S$, and avoid the closures of the remaining
strata. The resulting neighbourhood meets at most one component of each
original stratum, proving local finiteness. If $x\in\overline C^{,U}$,
every sufficiently small connected intersection $N\cap S$, with
$N\subset U$, meets $C$ and is contained in it. This proves condition d)
for $C$.

Let $C\subset S\cap U$ and $E\subset R\cap U$ be components and suppose
$x\in E\cap\overline C^{,U}$. The original frontier condition gives
$R\subset\overline S$. Choose $N\subset U$ about $x$ with $N\cap S$
connected, hence contained in $C$. For every $y\in E\cap N$, every
neighbourhood of $y$ inside $N$ meets $S$, hence $C$. Thus
$E\cap\overline C^{,U}$ is open in $E$; it is also closed and nonempty.
Connectedness of $E$ gives $E\subset\overline C^{,U}$. This is the
frontier condition for the restricted partition. Successive component
splitting agrees with direct restriction. Condition e), when imposed as
a condition on local product germs, is preserved by the same restriction.

Suppose global partitions $P,Q$ have the same restriction to each member
of an open cover. Fix a stratum $S$ of $P$. Its intersections with the
cover split into connected plaques, open in $S$. Each plaque is also a
plaque for $Q$. Two overlapping plaques belong to the same stratum of
$Q$. The graph whose vertices are these plaques and whose edges are
nonempty intersections is connected: a partition into two graph
components would partition $S$ into two nonempty disjoint open subsets.
Consequently $S$ lies in one $Q$-stratum. Interchanging $P,Q$ proves
equality of the partitions. This proves separatedness, answering the
margin of scan~182 affirmatively. It also explains why the deleted word
``separated'' cannot be restored in the preceding sentence.

For failure of gluing, choose distinct $p,q\in S^1$ and put
$U=S^1\setminus\{q\}$, $V=S^1\setminus\{p\}$. Stratify $U$ by $p$ and
the two components of $U\setminus\{p\}$, and stratify $V$ by its single
stratum. Both restrictions to $U\cap V$ are its two connected
components. Both charts satisfy a)--e). A global partition with these
restrictions must have strata $\{p\}$ and $S^1\setminus\{p\}$: the
second follows from its connectedness and the preceding plaque argument.
The germ of $S^1\setminus\{p\}$ at $p$ has two components, violating d).
This is also the boundary phenomenon in the marked-disc picture of
scan~181.

Because the presheaf is separated, its sheafification is its usual
``plus'' construction: a compatible atlas modulo common refinement.
To construct the leaves, generate an equivalence relation on $X$ by
membership in a connected local plaque. A leaf is given the final
topology for all its plaques. Plaques agree on open subplaques, so they
are an atlas of locally closed embeddings into $X$. The leaf inclusion
is a local immersion, and the disjoint union of leaves is the bijective
local immersion $D_0^*\to X$ of (296)--(297). This construction commutes
with open restriction. If $X$ is compact Hausdorff, a finite collection
of charts, each reduced to finitely many plaques, covers $X$; every leaf
contains one of these plaques. There are therefore finitely many leaves,
and they are locally compact. The Hausdorff condition is included here
to justify the local compactness assertion of scan~184.
\end{proof}

\Needspace{5\baselineskip}
\paragraph{F.32 (Proved). The germ flags form the nerve of an internal
topological category in every degree.}
\textit{Manuscript: MPL 147; section 13, folios 88-90; scans 185-189.}

For a chart $V$ with strata $S_i$, order them by
$S_i\leq S_j\iff\overline S_i^{,V}\subset\overline S_j^{,V}$.
The flag spaces of (299)--(303), with the literal topology specified
there, are
\[
 D_r|_V=
 \coprod_{i_0\leq\cdots\leq i_r}\overline S_{i_0}^{,V},
 \qquad
 D_r^*|_V=
 \coprod_{i_0\leq\cdots\leq i_r}S_{i_0}.
 \tag{F.h.1}
\]
These formulas glue canonically for every local stratification of F.31.
Every projection $D_r\to X$ is a local immersion; $D_r^*\subset D_r$
is open. Deleting or repeating entries gives continuous face and
degeneracy maps. For every $r\geq2$ the Segal map is a homeomorphism
\[
 D_r\ \xrightarrow{\ \sim\ }\
 \underbrace{D_1\times_{D_0}\cdots\times_{D_0}D_1}_{r\text{ factors}}.
 \tag{F.h.2}
\]
Thus $D_\bullet$ is exactly the nerve of an internal category in
topological spaces. Its fibre over $x$ is the nerve of the finite poset
of incident stratum germs at $x$. The first-vertex map $D_r\to D_0$ is
a finite covering locally on $D_0$; in particular it has the asserted
property over $D_0^*$. Under the literal model (F.h.1), this property
over all of $D_0$ is not special to normal crossings.

\begin{proof}
Let $W\subset V$ be open and $x\in\overline S_i\cap W$. Condition d)
selects a unique component $C_i$ of $S_i\cap W$ whose closure contains
$x$: a small connected $N\cap S_i$ lies in such a component, and any
other component accumulating at $x$ would meet that same $N$.
The choice is locally constant as $x$ moves in $\overline S_i\cap W$.
Indeed $N\cap S_i\subset C_i$ implies
$N\cap\overline S_i\subset\overline C_i$.
For a chain $S_{i_0}\leq\cdots\leq S_{i_r}$ these uniquely selected
components form a chain: near $x$, the lower component is contained in
the closure of the higher one, and the restricted frontier condition
proved in F.31 extends this inclusion over the entire connected lower
component. This gives a canonical homeomorphism between the restriction
of (F.h.1) and the formula computed after component splitting on $W$.
The identifications compose exactly, because they select the same germs.
They therefore glue on an arbitrary atlas and prove independence of
that atlas. This also gives the precise sense in which the associated
sheaf is described by local flags: an atlas defines compatible germ
flags, and the degree-zero strict pieces recover its local plaques.
One must retain the topology and the maps to $X$; bare incidence posets
do not determine a local stratification.

Deleting the first entry of a chain uses the inclusion
$\overline S_{i_0}\hookrightarrow\overline S_{i_1}$; deleting any other
entry uses the identity on the first closure. Repetition also uses the
identity. All are continuous and commute with chart changes. The
simplicial identities are the identities for deleting and repeating
entries, before any topological quotient is taken.

In the iterated fibre product in (F.h.2), composable arrows have the
same origin $x$ and the same middle \emph{germ}. Their concatenation
is a unique chain. On the component indexed by this chain, the fibre
product is
\[
 \overline S_{i_0}
 \times_{\overline S_{i_1}}\overline S_{i_1}
 \times_{\overline S_{i_2}}\cdots
 \times_{\overline S_{i_{r-1}}}\overline S_{i_{r-1}}
 \cong\overline S_{i_0}
\]
with its subspace topology. This proves a homeomorphism, not just a
bijection, and works for every $r$. Composition is concatenation followed
by forgetting the intermediate entry; associativity and units follow
from the same identities. The strict spaces alone do not form this
category, since deleting the initial stratum need not preserve strictness.

Fix a chart stratum $S_i$. Every stratum above it has closure containing
any selected point of $S_i$, so local finiteness implies that there are
only finitely many such strata. Over the copy of $\overline S_i$ in
$D_0$, the first-vertex map is a disjoint union of a finite number of
identity maps, one for every chain beginning with $i$. This proves the
covering assertion, including at boundary points of that copy. Thus the
aside in scan~189 claiming that extension from $D_0^*$ to all $D_0$ is
special to normal crossings is not a consequence of the definitions
as written. The manifolds and the Boolean incident posets of normal
crossings \emph{are} additional properties. Likewise the phrase about
degeneracies giving isomorphisms at scan~187 means isomorphisms onto the
equal-adjacent-entry summands, not $D_r\cong D_{r+1}$: the two-element
chain already has two vertices and three arrows. Passing to $\pi_0$
before the fibre products would invalidate this proof; F.25 addresses
that separate error.
\end{proof}

\Needspace{5\baselineskip}
\paragraph{F.33 (Refuted). Product germs in e) do not imply a stationary
link or a conical tube.}
\textit{Manuscript: MPL 147; section 13, folios 86 and 90; scans 181-182 and 190.}

There is a compact metrizable space with two strata satisfying a)--e)
whose punctured neighbourhood pro-homology at the closed stratum is not
constant. Hence e) alone cannot justify replacing its elementary
neighbourhood pro-topoi by fixed link models. This does not obstruct
the full pro-topos descent theorem below.

\begin{proof}
Construct an orientable surface $M$ by starting with a disc and attaching
successively compact genus-one cobordisms with two boundary circles.
Choose a compact exhaustion $K_n$ by connected genus-$n$ subsurfaces
with one boundary circle and $K_n\subset\operatorname{int}K_{n+1}$.
Then $M$ has one end and infinite genus. Let
$X=M\cup\{\infty\}$ be its one-point compactification. It is compact
metrizable and is stratified by $M$ and $\{\infty\}$. Both strata are
connected and locally connected; the partition is finite and satisfies
the frontier condition. The sets
$T_n=M\setminus K_n$ are a connected cofinal system of punctured
neighbourhoods of $\infty$, proving d). Along $M$, e) is the product
with a point. Along the point $\infty$, e) is the tautological product
$\{\infty\}\times(X,\infty)$. Its transverse stratum has connected
germ. Thus even the asserted transverse connectedness in scan~182 holds.

Choose the two handle generators $a_j,b_j$ in the $j$th cobordism and
let $c_n$ be its inner boundary class. Cutting off after the $m$th
handle gives a genus-$(m-n)$ surface with two boundary components,
whose first homology has basis
$c_n,a_{n+1},b_{n+1},\ldots,a_m,b_m$. Passing to the union yields
\[
 H_1(T_n;\mathbb Z)
  =\mathbb Z c_n\oplus\bigoplus_{j>n}
       (\mathbb Z a_j\oplus\mathbb Z b_j).
\]
Inclusion $T_m\hookrightarrow T_n$, $m>n$, sends $c_m$ to $c_n$
and preserves the surviving handle generators. Its image is
\[
 \mathbb Z c_n\oplus\bigoplus_{j>m}
       (\mathbb Z a_j\oplus\mathbb Z b_j),
 \tag{F.h.3}
\]
which strictly decreases with every $m$.

For completeness, an inverse system isomorphic in pro-abelian groups
to a constant group has eventually stationary images at each fixed
stage. Write such inverse isomorphisms as
$f:A_n\to B$ and compatible $g_m:B\to A_m$ after choosing a
representative of $f$. The equation $gf=1$ says that for sufficiently
large $k$ the transition $A_k\to A_m$ factors through $g_m$; the
compatibility $g_m=(A_k\to A_m)g_k$ gives the reverse inclusion of
images. Hence the eventual image equals $g_m(B)$. Formula (F.h.3)
contradicts this conclusion. A conical punctured tube
$L\times(0,\epsilon)$ with compatible shrinkings would have constant
pro-homology $H_1(L)$, so no such link model exists here. A finite
truncation of the image calculation is reproduced in the accompanying
program; the argument above proves strict decrease at every stage.
\end{proof}

\Needspace{5\baselineskip}
\paragraph{F.34 (Proved). Elementary pro-topoi, their gluing, and the
precise additional link hypotheses.}
\textit{Manuscript: MPL 147; section 13, folio 90; scan 190; gluing scans 150 and 155.}

There are two different assertions.
First, full sheaf descent holds for neighbourhood \emph{pro-topoi}
without a conical or asphericity assumption. In an ambient open $Y$
where a locally closed set $Z$ is closed, define
\[
 \mathfrak T_{Z,Y}=(\operatorname{Sh}(N))_{N\supset Z},
 \qquad
 \mathfrak L_{Z,Y;W}=(\operatorname{Sh}(N\cap W))_{N\supset Z},
 \tag{F.h.4}
\]
where $W$ is a specified open part of $Y$; in particular
$W=Y\setminus Z$ gives the punctured tube. Restriction to local
stratification charts and all their overlaps defines the same objects
for local plaques. For an incident flag
$\sigma=(S_0<\cdots<S_r)$, its elementary pro-topoi are the iterated
neighbourhood objects
\[
 \mathfrak E_{\sigma,W}
   =\bigl(\operatorname{Sh}
       (W\cap N_0\cap\cdots\cap N_r)\bigr)_{(N_0,\ldots,N_r)},
 \tag{F.h.5}
\]
where $N_j$ ranges over neighbourhoods of the indicated plaque $S_j$
in the chart, and $W$ is an independently specified locally closed
union of strata, with its \emph{subspace topology}. It need not be
open in the ambient chart. The empty flag gives the constant topos
$\operatorname{Sh}(W)$. One may take $W=S_r$, in which case the last
neighbourhood condition is redundant, or take a stratum closure, or
the specified punctured part of such a closure. Ambient open
restrictions make every closed embedding in this expression legitimate.
In a Cech or flag diagram $W$ is held fixed. Refining neighbourhoods
and deleting a neighbourhood condition give the comparison maps;
a specified inclusion $W\subset W'$ gives a further comparison.
There is no inclusion from one disjoint stratum into another when a
terminal vertex is deleted. Conical retractions, when available, are
separate geometric maps and retain their specified control data.
Common refinement makes the restriction identities strict. These
definitions do not assert that
$\mathfrak E_{\sigma,W}$ is constant.

Second, to model these objects by locally trivial link bundles, it is
sufficient to assume a locally compact paracompact Hausdorff space of
finite depth with compatible conical control data: local manifold
plaques $S$, tubes $T_S$, retractions $\pi_S$, and radius functions
$\rho_S$ with $\rho_S^{-1}(0)=S$, locally identified over $S$ with
\[
 S\times C(L_S),\qquad
 T_S\setminus S\cong L_S^{\rm bun}\times_S(S\times(0,\epsilon_S)).
 \tag{F.h.6}
\]
Here $L_S$ is compact and recursively locally contractible and
conically stratified; $L_S^{\rm bun}\to S$ is a locally trivial bundle,
so monodromy is retained. Changes of charts preserve the strata, cone
radius, and these control data. For incident plaques $S<T$ require
$\pi_S\pi_T=\pi_S$ and $\rho_S\pi_T=\rho_S$ wherever defined;
tubes of incomparable plaques are disjoint after shrinking. The data
may be given on a locally finite atlas with its full overlap nerve,
which allows distinct local branches of one global leaf. This is an
explicit sufficient hypothesis, not a conclusion from e). Whitney
stratifications in the usual locally compact smooth setting provide
such controlled tube germs by the control-data and isotopy results of
Mather (1970, \S7 and Corollaries~10.2--10.6)\citep{Mather1970Local}.

With these data, the full Cech diagram of tubes covering a closed
union of strata, or its subdivision by all flags, has the required
two-colimit in pro-Grothendieck topoi. In particular, for a closed
union $Z\subset Y$,
\[
 \operatorname{Sh}(Y)\ \simeq\
 \operatorname{Sh}(Y\setminus Z)
       \amalg_{\mathfrak L_{Z,Y;Y\setminus Z}}
       \mathfrak T_{Z,Y}
 \quad\text{in pro-Grothendieck topoi}.
 \tag{F.h.7}
\]
At the level of homotopy types these are homotopy colimits of the
\emph{full} link and iterated-link bundles. No $K(\pi,1)$ assumption
is required for this gluing. To replace every elementary homotopy type
by its fundamental groupoid, the exact condition is that every such
type be a $1$-type. A sufficient geometric condition is that the base
strata and all the relevant iterated-link fibres be $1$-types; for
connected components this means $K(\pi,1)$ spaces. The reconstruction
then uses the homotopy colimit of their classifying \emph{spaces}, not
an ordinary two-colimit of classifying topoi or groupoids.

\begin{proof}
At a fixed open-neighbourhood stage, an open cover $N=\bigcup N_a$
gives the usual sheaf descent equivalence, including morphisms:
a sheaf on $N$ is a family of sheaves on $N_a$ with isomorphisms on
$N_a\cap N_b$ satisfying the triple cocycle. For any target topos
$\mathcal A$, apply this objectwise to inverse-image functors
$\mathcal A\to\operatorname{Sh}(N_a)$. Their finite-limit and
colimit preservation can be checked on the conservative open
restriction family; hence the glued functor is again the inverse image
of a geometric morphism. Natural transformations glue as well. This
proves the bicategorical universal property for full topoi, not only
for local systems.

The two-open case makes passage to pro-objects completely explicit.
For a fixed target $\mathcal A$, filtered two-colimits of mapping
\emph{categories} commute with finite pseudolimits in
$\operatorname{Cat}$. This follows by representing the finitely many
objects, arrows, invertible comparison arrows, and their coherence
equalities at one common stage; the same argument handles arbitrary
morphisms and equality of morphisms. Apply this to
\[
 \operatorname{Geom}(\operatorname{Sh}(Y),\mathcal A)
 \simeq
 \operatorname{Geom}(\operatorname{Sh}(Y\setminus Z),\mathcal A)
 \mathbin{\times}^{\rm ps}_{
   \operatorname{Geom}(\operatorname{Sh}(N\setminus Z),\mathcal A)}
 \operatorname{Geom}(\operatorname{Sh}(N),\mathcal A).
\]
The left side and first factor are constant in $N$. Taking the filtered
two-colimit gives precisely the mapping-out universal property of
(F.h.7). For a pro-target $(\mathcal A_j)$, take the additional limit
over $j$; pseudolimits commute with the displayed finite pseudolimits.
This proves
(F.h.7) for arbitrary $Z$ as stated, without appealing to a putative
inverse limit of punctured full sheaf topoi. The latter can be
degenerate, as F.12 proves.

We next explain the geometric content of the extra hypothesis. On a
conical chart, compactness of $L_S$ implies that every neighbourhood of
the cone vertex contains a full radial ball of some positive radius.
Over each relatively small open subset of $S$, the same holds with a
positive local bound. A locally finite partition of unity produces a
positive continuous minorant of these bounds. Consequently tubes with
positive continuous radii form a cofinal system of all tube germs.
The control identities make these choices compatible on incident
plaques. Choose the radii successively in the finite depth, reducing
them to satisfy all earlier conditions. Only finitely many conditions
occur near a given point. Incomparable tubes are disjoint, so a
nonempty intersection of distinct tubes is indexed by an incident
flag. It is exactly a stage of (F.h.5). Recursing inside $L_S$ shows
that its angular part is the corresponding iterated-link bundle;
the remaining coordinates are positive radial intervals. This is
proved inductively on depth from (F.h.6): an incident higher plaque
is $S$ times a stratum in $L_S$ times the radial interval, and the
next tube condition is its corresponding link-tube condition. There
is no assertion that the angular bundle is a product of its fibres.

At every simultaneous radius stage, the tubes covering the given
closed union give an actual open cover of their union, and their full
intersections give its Cech nerve. Thus the preceding full-topos
descent applies. Passing from ordered intersections to all incident
flags is subdivision of the same descent datum: its one-step maps
are restrictions and its triangles are the cocycle relations. It
does not discard higher flags. For an infinite locally finite atlas,
one must also justify synchronization of pro-stages. A datum and its
isomorphisms are represented on individual charts, double overlaps,
and triple overlaps. Choose a common refinement on each such finite
family; at any point only finitely many bounds are imposed. A positive
continuous minorant gives one global radius assignment satisfying
all of them. The same argument applies to transformations and equality
of transformations. This proves essential surjectivity and full
faithfulness on the mapping categories, and justifies the pro-colimit
without commuting a filtered colimit with an arbitrary infinite
product. The same argument applies to the Cech groupoid of the stated
locally finite Hausdorff chart cover when plaques have branch
identifications: on compact chart pieces only finitely many overlap
conditions occur. This is a chart presentation of the given space,
not a claim about an arbitrary nonproper etale quotient atlas.
These constructions establish the
extension of both versions of F.13 to the stated controlled local
stratifications.

For shapes, apply the open-cover homotopy-colimit statement at every
stage. It can be proved directly by the thick Cech realization: a
partition of unity subordinate to a numerable cover provides a section
up to homotopy, and the simplex of possible partition weights over
each point is contractible. Refining charts and radii preserves this
comparison. A conical radial interval is contractible and its shrinkings
are homotopy equivalences, so the iterated-link bundles give constant
models for the resulting local \emph{pro-homotopy types}. They do not
give equivalences of full sheaf topoi: even
$\operatorname{Sh}((0,\epsilon))$ is not $\operatorname{Set}$.
The homotopy long exact sequence of each locally trivial link bundle
shows that an aspherical fibre over an aspherical base has aspherical
total space, proving the sufficient $1$-type criterion. Conversely a
space with nonzero homotopy above degree one cannot be replaced by
the classifying space of its fundamental groupoid. This is precisely
the information discarded by a level-one replacement.

The categorical cases are STRATA C.3, C.5, and C.6. The relevant
calculation in C.6 is independently checkable: for a full sieve
$A\subset C$ and complementary cosieve $B$, restriction identifies a
presheaf with $(F_A,F_B,F_B\to i^*j_*F_A)$, where
$(i^*j_*F_A)(b)=\lim_{(A/b)^{\rm op}}F_A$. The inverse sends this
triple to $j_*F_A\times_{i_*i^*j_*F_A}i_*F_B$, as is verified on
$A$ and $B$ separately. With spaces one must take the coherent limit
and pullback. Repeating this argument along a finite filtration gives
the actual crossing-flag data of C.5; its shape is the homotopy
pushout of C.3. This verifies the categorical pattern but is not
substituted for the full-topos geometric descent just proved.

Finally consider the two requested tests. For $(S^2,\{p\})$, take
$N$ a small open disc and $U=S^2\setminus\{p\}\cong\mathbb R^2$.
The punctured tube is $S^1\times(0,\epsilon)$, so (F.h.7) has shape
\[
 \operatorname{hocolim}(*\leftarrow S^1\longrightarrow *)\simeq S^2.
\]
Both open and closed strata are simply connected. The topos of sheaves
constructible for these two strata is the arrow topos: the local-system
gluing functor sends a set to that same set, since the link is connected
and the local system on $U$ is constant. Its classifying shape is
contractible. Similarly the ordinary groupoid pushout of the displayed
diagram is trivial. The full diagram with its link retains the missing
$H_2(S^2;\mathbb Z)=\mathbb Z$. Thus even aspherical elementary
pieces do not license an ordinary level-one colimit.

The drawing at scan~190 is three half-lines in the \emph{plane}:
there are seven strata, namely the origin, three open rays, and three
open sectors. The origin's link is a circle divided into three points
and three open arcs. The incident poset has seven elements, twelve
strict comparable pairs, and six strict three-element chains. Removing
the origin gives the six-vertex, six-edge incidence cycle, with first
Betti number one. Adjoining the origin cones off that cycle, as the
disc tube does. The three ray links are two-point spaces and the
sector links are empty. All these elementary angular pieces are
$1$-types, but the seven-element incident poset is not Boolean, since
seven is not a power of two. Treating the picture as only a tripod
would remove the three sectors and give a three-point link, contradicting
both the written count and the topology. The program checks the flag
counts, the Segal bijections through degree six, and the circle and
sphere homology calculations; the proof of (F.h.2) establishes all
degrees.
\end{proof}


\subsection{Additional arithmetic tests and coherent limits}

\Needspace{5\baselineskip}
\paragraph{F.35 (Proved). Inertia injection and extension as a sub-germ.}
\textit{Manuscript: Additional F(h)(1); MPL149 section4 folios66-67 scans133-135.}

Let $X/K$ be smooth in characteristic zero, $D$ a normal-crossings divisor,
$U=X\setminus D$, and $\bar x$ a geometric point of $D$. Choose a geometric
component of $U$ and a path from its base point to the strict local
punctured neighbourhood at $\bar x$. Write
\[
 H=\widehat{\mathbb Z}(1)^{I(\bar x)},\qquad
 G=\pi_1(U_{\overline K}),\qquad \iota:H\longrightarrow G.
\]
Then $\iota$ is injective if and only if every finite tame covering germ
at $\bar x$ is an open-and-closed sub-germ of the restriction of some
finite etale covering of $U_{\overline K}$. Equivalently, that covering
may be taken over $U$ after allowing a finite extension of the constants
and then composing its covering with $U_{K'}\to U$. The statement does
not assert that the entire restriction of the global covering is the
prescribed germ. This proves the first additional test F(h).

\begin{proof}
We prove the group-theoretic assertion for any continuous homomorphism
$\iota:H\to G$ of profinite groups. If it is injective, identify the compact
$H$ with a closed subgroup of $G$. A transitive finite $H$-set is $H/K$
for an open subgroup $K\subset H$. The compact subset $H\setminus K$
does not contain $1$, so some open normal $N\lhd G$ satisfies
$N\cap H\subset K$. Put $J=KN$. Since $N$ is normal, $J$ is a subgroup;
since it contains $N$, it is open. Also
\[
 J\cap H=K:
 \quad kn\in H\Longrightarrow n\in H\cap N\subset K.
\]
Consequently $hK\mapsto hJ$ embeds $H/K$ into the restriction of the
finite $G$-set $G/J$. Apply this to each of the finitely many orbits of
an arbitrary finite $H$-set. Conversely, if $1\ne h\in\ker\iota$, some
finite continuous quotient of $H$ detects $h$. Its regular $H$-set
cannot embed in a restricted $G$-set, on which $h$ acts trivially.
This proves the equivalence.

Finite tame germs are finite continuous $H$-sets by the strict local
Kummer calculation of F.21. Finite etale coverings of the geometric
component of $U$ are finite $G$-sets. An invariant subset of a finite
covering is precisely an open-and-closed subcover, so the group assertion
has exactly the stated geometric interpretation. The global covering and its finite presentation descend to a finite
extension $K'/K$; in characteristic
zero this extension is separable. Composition with the finite etale map
$U_{K'}\to U$ gives the stated version over $K$. The comparison remains
geometric: for a nonclosed point one must not substitute the whole
arithmetic decomposition group for $H$.
\end{proof}

\Needspace{5\baselineskip}
\paragraph{F.36 (Proved). The coherent etale limit is nondegenerate.}
\textit{Manuscript: Supplemental LM8 coherent-etale-limit question; context MPL147 sections11-12 folios70-81 scans143-148 and157-169; MPL149 section4 folios66-67 scans133-135.}

Let $x$ be a closed point on a smooth curve over an algebraically closed
field $k$ of characteristic zero. Let $R=\mathcal O_{X,x}^{\mathrm{sh}}$,
$F=\operatorname{Frac}R$, and let $(V_i,\bar x_i)$ run through pointed
affine etale neighbourhoods of $\bar x$, shrunk so that a local parameter
$t$ is defined and its only zero in $V_i$ lies over $x$. Put
$V_i^*=V_i[1/t]$. Then
\[
 \varprojlim_i (V_i^*)_{\mathrm{et}}
   \simeq (\operatorname{Spec}F)_{\mathrm{et}}
   \simeq B\widehat{\mathbb Z}(1)
\]
as a bicategorical limit of Grothendieck topoi. The hypotheses of the
coherent-limit assertion in the LM8 intake are therefore sufficient.
The exact general references are SGA~4, VI, Theorems~8.2.3 and~8.3.13,
and VII, Lemma~5.6; the last statement is also the Stacks Project,
Lemma~59.51.2 \citep{SGA4II,StacksEtaleLimits}.

\begin{proof}
The schemes $V_i^*$ are affine, the transition maps are affine, and
$\varinjlim_i\mathcal O(V_i)=R$. Localization commutes with that colimit,
so their scheme limit is $\operatorname{Spec}R[1/t]=\operatorname{Spec}F$.
An affine etale algebra of finite presentation over this limit is given
by finitely many coefficients and finitely many equations; it descends
to some $V_i^*$. The same finite-presentation argument descends morphisms
and equalities. A covering can be replaced by a finite affine covering
of each affine object, and surjectivity of that finite family descends
after increasing $i$. Thus the restricted etale site of the limit is
the filtered colimit of the restricted sites, including the covering
families. Passing to sheaves gives exactly the inverse-topos universal
property of VI~8.2.3.

Each stage is a coherent topos and every transition morphism is coherent.
One can see the extra assertion of VI~8.3.13 directly here: a coherent
object is the quotient of a finite diagram of coherent generators by a
finitely presented equivalence relation. Generators, relation, and its
finite reflexivity, symmetry and transitivity identities all descend to
one stage. Morphisms descend by the same finite presentation. Hence
coherent objects in the limit are the filtered two-colimit of coherent
objects at the stages; no analytic infinite annulus covering is present.

The ring $R$ is a strictly henselian discrete valuation ring. Its residue
field has characteristic zero, so all finite extensions of $F$ are tame.
Every unit of $R$ has an $n$th root by Hensel's lemma. The tame extension
classification therefore gives the separable closure
$\bigcup_n F(t^{1/n})$, with Galois group
$\varprojlim_n\mu_n=\widehat{\mathbb Z}(1)$; the classification follows
by removing the ramification index with $t^{1/n}$, after which the
extension is unramified and hence trivial over a strictly henselian
ring. This also follows from the tame valuation calculation in the
Stacks Project, Lemma~15.116.2 \citep{StacksTameDVR}.
The topos is not degenerate: its terminal finite set is not its empty
set. In residue characteristic $p$, the tame quotient is instead
$\prod_{\ell\ne p}\mathbb Z_\ell(1)$ and the full group has wild inertia.

There is no conflict with F.12 or LM8 E.11. The classical shrinking-disc
system has the annular covering that kills every cone. The etale system
has finite coherent covering descent and the nonempty scheme limit
$\operatorname{Spec}F$. MPL~147, scans~143--148 and 157--169, explicitly
works with analytic neighbourhoods, arcs, and analytic multiplicities;
its $D,D^*$ are the classical pro-objects. The phrase ``tubulures etales''
of MPL~149, scans~133--135, proposes this different arithmetic construction.
\end{proof}

\Needspace{5\baselineskip}
\paragraph{F.37 (Refuted). Pinned fullness over a finite field.}
\textit{Manuscript: MPL149; section2; folio25v; scan52; taskF(b).}

Even after localization at universal homeomorphisms, the etale-topos
functor on finite-type $\mathbb F_q$-schemes is not full onto pinned
morphisms over $B\mathbb F_q$. This does not contradict F.4, which proves
faithfulness. It also explains why the infinite-field hypothesis in
\cite[Theorem~0.1]{CHW2024} must not be silently dropped.

\begin{proof}
Let $X=\mathbb A^1_{\mathbb F_q}$, let $p:X_{\mathrm{et}}\to
B\mathbb F_q$ be the structure morphism, and identify
$\Gamma_{\mathbb F_q}$ with $\widehat{\mathbb Z}$. Since
$\pi_1(X)\to\Gamma_{\mathbb F_q}$ is onto and the latter is abelian,
every $a\in\widehat{\mathbb Z}$ gives a natural automorphism of $p^*$:
on a finite constant-field covering it acts by $a$. Thus
\[
 (\operatorname{id}_{X_{\mathrm{et}}},a):
 X_{\mathrm{et}}\longrightarrow X_{\mathrm{et}}
 \quad\text{over }B\mathbb F_q
\]
is a relative geometric morphism. Its underlying map of spaces is the
identity, so it is pinned.

We claim that the identity geometric morphism of $X_{\mathrm{et}}$ has
no nontrivial natural automorphism. Such an automorphism restricts to
every nonempty open of $X$, and then to its generic point. Every finite
separable extension of $\mathbb F_q(t)$ spreads to a finite etale cover
of some such open, so restriction gives an element of
$Z(G_{\mathbb F_q(t)})$. This restriction is injective. Indeed affine
etale schemes are generators; an automorphism of such a representable
which is generically the identity is the identity, because its source
is reduced and its affine target is separated. Affine generators are
essential here: one need not make a generic-detection assertion for
arbitrary nonseparated etale objects.

To prove $Z(G_{\mathbb F_q(t)})=1$, let a central element $z$ project to
$\sigma\in\Gamma_{\mathbb F_q}$. Write
$N=G_{\overline{\mathbb F}_q(t)}$. For any prime $\ell\ne p$, the Kummer
classes of $t-b$, $b\in\overline{\mathbb F}_q$, are linearly independent
in $\overline{\mathbb F}_q(t)^*/(\overline{\mathbb F}_q(t)^*)^\ell$,
as their valuations at the distinct points $b$ show. The conjugation
outer action of $\sigma$ on $N$ permutes these classes by
$b\mapsto\sigma(b)$, with the scalar factor prescribed by the action
on $\mu_\ell$. If $\sigma$ moves any $b$, this cannot be the identity
on $\operatorname{Hom}(N,\mu_\ell)$: distinct valuation basis vectors
cannot become one another by a nonzero scalar. Centrality therefore
forces $\sigma=1$. Hence $z\in Z(N)$, which is trivial by the valuation
and Kummer proof of F.2, unchanged in characteristic $p$ when its auxiliary
prime is chosen different from $p$. We have proved the claim.

An isomorphism between the two relative maps indexed by $a,b$ would now
have to be the identity automorphism of the underlying identity map;
compatibility with their base isomorphisms forces $a=b$. There are
therefore uncountably many isomorphism classes of pinned relative maps
of this form. In contrast, the category of finite-type schemes over a
finite field has a countable skeleton and countable Hom-sets: all its
affine presentations and finite gluing data use finite strings over a
countable set. Its localization at universal homeomorphisms still has
countable Hom-sets, since morphisms are represented by finite zigzags.
It cannot surject onto this uncountable set of pinned morphisms.
This is a complete cardinality counterexample; it does not require
choosing an unspecified nonintegral $a$ and asserting that particular
one is nonalgebraic. Powers of Frobenius account for some, but not all,
of the relative twists.
\end{proof}

\Needspace{5\baselineskip}
\paragraph{F.38 (Known). The verified reconstruction range over larger bases.}
\textit{Manuscript: MPL149; section2; folio25v; scan52.}

The following are the precise positive ranges used for MPL~149,
scan~52. Over an infinite finitely generated field, use pinned morphisms
and invert universal homeomorphisms (F.5). The same holds over every
sub-$p$-adic field, meaning a subfield of a finitely generated extension
of $\mathbb Q_p$. More generally, let $S$ be noetherian, Nagata and
geometrically unibranch, with sub-$p$-adic function field on every
irreducible component. On finite-type $S$-schemes which are
\emph{componentwise universally open}, the localization at universal
homeomorphisms is fully faithful into topoi over $S_{\mathrm{et}}$ with
\emph{chivalrous} morphisms: the induced continuous map sends
constructible subsets to constructible subsets. In particular normal
flat finite-type schemes over $\operatorname{Spec}\mathbb Z$ are
reconstructed, with all their morphisms. These are Theorem~4.1,
Proposition~6.11 and Theorem~6.12 of the opened September~20, 2026 version
of Carlson--Stix \citep{CarlsonStix2026}; the relative hypothesis is part
of the theorem and is not a theorem about all schemes over $S$.

\begin{proof}
We spell out the reduction from the quoted statements to the formulation
at scan~52. A universal homeomorphism induces an equivalence of small
etale topoi, so the functor factors through the stated localization.
Proposition~6.11 identifies that localization, in the componentwise
universally open category, with absolute weak normalizations. Theorem~6.12
identifies the Hom-set from $X^{\mathrm{awn}}$ to $Y$ with the groupoid
of chivalrous relative topoi morphisms; its rigidity assertion says that
this groupoid is equivalent to that set, with no extra automorphisms.
For a normal $X$ in this category, $X=X^{\mathrm{awn}}$ by their
Lemma~6.10. Thus the asserted Hom-equivalence follows, and applying it
to a topos equivalence and its inverse gives inverse scheme morphisms
by faithfulness. Flat finite-type schemes over $\mathbb Z$ satisfy the
componentwise openness condition, giving the stated example.

Over a field all these schemes are componentwise universally open and
finite-type field morphisms send closed points to closed points. Their
Theorem~4.1, together with the seminormalization description of the
localization in characteristic zero, gives the pinned field assertion.
In the general-base theorem, pinned is the wrong substitute for
chivalrous: for a discrete valuation ring $R$ and nonunit $a\ne0$, the
closed point $(aT-1)$ of $\mathbb A^1_R$ maps to the generic point of
$\operatorname{Spec}R$. This genuine scheme morphism is chivalrous by
constructibility of images, but not pinned. The proof of the relative
theorem reconstructs over the generic fields and spreads out using
normal quasi-sections; this is exactly where geometric unibranchness,
Nagata finiteness and componentwise openness enter. Thus its application
here does not discard any of those hypotheses.

Combining F.5--F.7 and F.37 gives the necessary qualifications to the
provisional question: unrestricted fullness is false even over number
fields; pinned fullness holds in the preceding ranges, fails over finite
fields, and fails over algebraically closed characteristic-zero fields.
Faithfulness of genuine scheme maps remains the all-characteristic
statement F.4. The remaining general-field question is recorded precisely
in F.44, not identified with the false finite-field assertion.
\end{proof}


\Needspace{5\baselineskip}
\paragraph{F.39 (Proved). Stable towers, full inertia and cyclotomic normalization.}
\textit{Manuscript: MPL149 section4 folios66-67 scans133-135; EXIT task6 andP3b rows.}

Let $U\to B$ be an elementary hyperbolic fibration over a field $K$
finitely generated over $\mathbb Q$, with $B$ a connected smooth
finite-type $K$-variety: a smooth proper curve family with
a nonempty finite etale marked divisor removed, of type $(g,n)$ with
$2g-2+n>0$. Then each horizontal geometric puncture group injects into
$\pi_1(U_{\bar K})$. For a recursively compatible tower of these
fibrations, equip its base with a normal-crossings completion whose
strata are aspherical and whose full meridian lattices inject both
discretely and profinitely. Suppose the fibration extends to a stable
pointed family there, after finite root or quotient changes, with a
smooth toroidal total model and fibre incidence type constant on each
base stratum. Then every total full meridian group injects:
\[
 I_{\bar x}=\widehat{\mathbb Z}(1)^{I(\bar x)}
       \lhook\joinrel\longrightarrow\pi_1(U_{\bar K}).
\]
These assertions are natural for etale changes and compatible monomial
maps, with path conjugacy retained. On a splitting field of the
branches, conjugation on each normal factor is the cyclotomic character;
without splitting, it is the product of that action with the branch
permutation action. A monomial map $t'_j=u_j\prod_i t_i^{m_{ji}}$ induces
the integer contact matrix $(m_{ji})$ on inertia. This supplies a common
geometric injection and normalization theorem for the arithmetic issue
raised in the sixth EXIT task. F.48 audits the recovered
\texttt{CYCLOTOMIC\_KERNEL} rows; the remaining literal comparison
with the unlocated \texttt{FLAGGED\_149} rows is recorded in F.45.

\begin{proof}
We give the completion argument, since discrete injection alone is
insufficient. Every stage of the tower has this smooth finite-type
base geometry. After choosing $\bar K\hookrightarrow\mathbb C$, the
proper smooth family, with its finite etale markings, is locally a
trivial differentiable family of marked compact surfaces: a horizontal
distribution tangent to the markings is obtained by a partition of
unity and its local flows give the trivializations. Removing the
markings therefore gives a locally trivial analytic topological bundle.
Its fibre is a punctured surface with group a free
group $N=F_r$, $r\ge2$. For a bundle over any connected base, the image
of $\pi_2(B)\to\pi_1(F)$ is central: pull back over a sphere, write its
clutching loop as $h_t\in\operatorname{Homeo}(F)$, and use the square
$h_t(\gamma(s))$ to commute the evaluation loop $h_t(x)$ with every
$\gamma\in\pi_1(F,x)$. The free group's centre is trivial, so the fibre
injects. A puncture loop is nontrivial: if there are several punctures
it survives in homology; with one puncture it is the reduced nonempty
product of handle commutators. Thus each such loop has infinite order.

We prove that for \emph{every} extension
$1\to F_r\to G\to Q\to1$ the completed sequence is exact on the left.
First $Z(\widehat F_r)=1$. Given a finite quotient $F_r\to A$, with
kernel $L$, choose a prime $\ell\nmid|A|$. The cellular chain complex
of the finite regular covering of the $r$-rose gives
\[
 0\longrightarrow H_1(L,\mathbb F_\ell)
 \longrightarrow\mathbb F_\ell[A]^r
 \longrightarrow\mathbb F_\ell[A]
 \longrightarrow\mathbb F_\ell\longrightarrow0.
\]
Averaging over $A$ splits these representations, so
$H_1(L,\mathbb F_\ell)\simeq
\mathbb F_\ell\oplus\mathbb F_\ell[A]^{r-1}$ is a faithful $A$-module.
The finite quotient $F_r/[L,L]L^\ell$ consequently detects the image in
$A$ of any purported central element of $\widehat F_r$: centrality
forces its action on this module to be trivial, and hence its image
in $A$ is $1$. This for every $A$ proves centrelessness.

The automorphism group of a finitely generated profinite group is
profinite: the intersections of its finitely many open subgroups of
index at most $m$ form a characteristic open basis, and automorphisms
are the compatible invertible maps on the resulting finite quotients.
The conjugation map
$G\to\operatorname{Aut}(\widehat F_r)$ therefore extends to
$\widehat G$. Its composite with
$\widehat F_r\to\widehat G$ is the inner-automorphism map, which is
injective by centrelessness. Thus $\widehat F_r\to\widehat G$ is
injective. Its image is the closure of $F_r$, and the quotient of
$\widehat G$ by that closure is $\widehat Q$, since both classify the
finite quotients killing $F_r$.

A nontrivial cyclic subgroup $\langle w\rangle\subset F_r$ has its full
profinite topology. Indeed represent a cyclically reduced conjugate
of $w$ by an immersed labelled circle in a rose. For every label, pair
the missing incoming and outgoing edges to complete this finite graph
to a finite covering of the rose, without identifying its vertices.
The original circle is an embedded cycle in that graph, so its integral
homology class is primitive. The corresponding finite-index subgroup
$H\subset F_r$ has a homomorphism $H\to\mathbb Z$ sending $w$ to $1$.
For any $m$ the core in $F_r$ of the kernel of $H\to\mathbb Z/m$
gives a finite quotient in which $w$ has order divisible by $m$.
This proves $\widehat{\langle w\rangle}\simeq\widehat{\mathbb Z}
\hookrightarrow\widehat F_r$. Together with the extension argument it
proves horizontal profinite injection over an arbitrary base. By
induction a tower is aspherical, and its discrete group is residually
finite: separate an element either in the completed base or in the
embedded completed free fibre.

For the vertical assertion let $A$ be a total local normal lattice,
$A_B$ its base lattice, and $M:A\to A_B$ the integral monomial map.
At a smooth vertical point $\ker M=0$; at a marking it is generated
by the fibre puncture; at a node $xy=t$ it is generated by $(1,-1)$
in the map $(u,v)\mapsto u+v$. In the latter case its image in the smooth
fibre is the vanishing cycle. Root changes and smooth toroidal charts
replace these lattices by finite-index or primitive sublattices; hence
in general $\ker M$ has rank at most one and maps to a nonzero power
of the corresponding cycle. The cycle is essential. If a neck bounded
an unpunctured disc, the side it cut off would be an unmarked genus-zero
tree attached at one node, impossible for a stable pointed curve:
the sum of the special-point counts on such a tree is less than three
times its number of components. A neck around one marked point remains
a nontrivial puncture loop, and the extra two-punctured bridges of a
toroidal model also have essential meridians.

The diagram
\[
\begin{array}{ccccc}
 \ker M&\longrightarrow&A&\longrightarrow&A_B\\
 \downarrow&&\downarrow&&\downarrow\\
 \pi_1(F)&\longrightarrow&\pi_1(U)&\longrightarrow&\pi_1(B)
\end{array}
\]
has the needed exactness in its rows. An element of $A$ killed in
$\pi_1(U)$ maps to zero in the injective base lattice, then belongs to
$\ker M$, where the essential-cycle map is injective. It is zero.
Complete the lower sequence by the free-kernel result. Smith normal
form proves $\ker\widehat M=\widehat{\ker M}$: the image of $M$ is
free abelian, the sequence onto it splits, and its inclusion into
$A_B$ remains injective after completion. The same diagram chase,
using completed base injection and cyclic fibre injection, proves the
full profinite assertion. The strata are bundles of punctured stable
components or tori over finite incidence covers of base strata; their
homotopy sequences prove their asphericity. These statements give the
induction for every recursively compatible tower in the theorem.

The Riemann existence comparison is SGA~1, XII, Theorem~5.1 and
Corollary~5.2 \citep{SGA1}. It identifies these completed groups and
Kummer meridians with the geometric etale groups, so it transports
the proved injections. In local parameters the automorphisms of
$t_i^{1/n}$ are multiplication by $\mu_n$; Galois conjugation acts on
these roots of unity by the cyclotomic character. Taking all $n$
gives the asserted Tate modules. Pulling back $t'_j$ gives the contact
matrix, while the roots of the units $u_j$ alter the tangential
splittings by their Kummer cocycles. They do not alter the homomorphism
on geometric meridians. Thus neither a preferred generator of
$\widehat{\mathbb Z}(1)$ nor a preferred tangential path has been
asserted.

Finally, if a local arithmetic extension
$1\to I_x\to E_x\to G_{\kappa(x)}\to1$ maps to that of $U/K$ and the
chosen map $G_{\kappa(x)}\to G_K$ is injective, then
\[
 \ker(E_x\to\pi_1(U))=
 \ker(I_x\to\pi_1(U_{\bar K})).
\]
Project a kernel element to $G_{\kappa(x)}$ to prove this equality.
The condition holds at a closed point after its finite residue-field
extension, but is not automatic at a transcendental point.

There are two precise limits on this theorem. First, cyclotomic action
on a subgroup is a necessary normalization of genuine inertia; we have
not asserted that it characterizes inertia among all subgroups.
Second, stable extension alone is insufficient for an arbitrarily
chosen base completion. Blow up a smooth boundary point of the product
completion of $(\mathbb P^1\setminus\{0,1,\infty\})^2$. At the new
double point the exceptional and strict-transform meridians both map
to the old meridian, so $(1,-1)$ is killed. A constant stable curve
family over that base retains this kernel, discretely and profinitely.
The compatibility and full base-lattice assumptions above are therefore
necessary to this induction. This also verifies, rather than merely
quotes, the distinction between EXIT P3b-vertical and its counterexample.
\end{proof}

\subsection{F(h), added item (2): the punctured completion and the end}

\Needspace{5\baselineskip}
\paragraph{F.40 (Proved). Nonarchimedean invariance, with the missing hypothesis.}
\textit{Manuscript: MPL 149; section4; folio66; scan133; additional F(h)(2) tests.}

Let $k$ be any field, equipped with the trivial absolute value, let $X$ be
a separated $k$-scheme of finite type, and let $U\subset X$ be a dense open
subscheme with complement $D$. Denote by $Y^{\beth}$ the generic fibre of
the scheme $Y$ regarded as a formal $k$-scheme with the discrete topology.
The analytic punctured tube is
\[
 W(X,D)=X^{\beth}\setminus\bigl(U^{\beth}\cup D^{\beth}\bigr).
\]
It is the punctured generic fibre associated with $\widehat X_D$.
For every proper morphism $X'\to X$ which is the identity on $U$, the
induced morphism $W(X',X'\setminus U)\to W(X,D)$ is an isomorphism of
$k$-analytic spaces. If $X$ is proper over $k$, then
\[
                 W(X,D)=U^{\mathrm{an}}\setminus U^{\beth}=:U_\infty^{\mathrm{na}}.
\]
Consequently it is independent of the proper compactification, without
smoothness, normal-crossings, or characteristic-zero assumptions. The
phrase ``with no hypothesis'' is false if it also removes properness.

\begin{proof}
A point of $Y^{\beth}$ is a multiplicative seminorm whose valuation ring
has a centre on $Y$; separatedness makes the centre unique. A point of
$W(X,D)$ has its generic point in $U$ and its centre in $D$. The valuative
criterion for a proper $X'\to X$ gives a unique lifted centre, necessarily
in $X'\setminus U$. Thus both tubes are the same analytic domain of
$U^{\mathrm{an}}$, with the same induced analytic structure. If $X$ is
proper, every analytic point has a centre, so $X^{\beth}=X^{\mathrm{an}}$
and $D^{\beth}=D^{\mathrm{an}}$, giving the formula. This is also the
precise statement of Ben-Bassat--Temkin, \S3.3.2 and Proposition~4.5.1
\citep{BenBassatTemkin2012}.

For necessity, take $U=\mathbb G_m$ over an algebraically closed field.
Its embedding in $\mathbb A^1$ has boundary $\{0\}$ and tube
$0<|t|<1$, a single open interval. Its proper compactification
$\mathbb P^1$ has two boundary points and tube
$\{0<|t|<1\}\amalg\{1<|t|<\infty\}$, two open intervals. These are not
even equally connected. The first ambient space is a partial
compactification, so it is not a counterexample to the properly stated
theorem.

The notation $\widehat X_D\setminus D$ on MPL~149, fol.~66, scan~133
must be given this punctured-generic-fibre meaning: literally deleting
$D$ from the underlying topological space of the formal completion gives
the empty space. In particular, the invariant does not assert that the
formal schemes $\widehat X_D$ themselves are isomorphic.
\end{proof}

\Needspace{5\baselineskip}
\paragraph{F.41 (Refuted). Stronger replacement: the archimedean link is invariant too.}
\textit{Manuscript: MPL 149; section4; folio66; scan133; additional F(h)(2) tests.}

The proposed failure of compactification invariance for the archimedean
punctured neighbourhood when $X$ is compact is false. For every compact
Hausdorff compactification $U\subset X$, with $D=X\setminus U$, its
punctured-neighbourhood pro-space is canonically
\[
 \{N\setminus D\}_{N\supset D}
       =\{U\setminus K\}_{K\subset U\ \mathrm{compact}}.
\]
Here $N$ ranges over open neighbourhoods, and the equality is an
isomorphism of indexing systems, not just equality of cohomology.
For a smooth proper complex variety with a simple normal-crossings
boundary, this pro-homotopy type is constant, represented by the boundary
$L$ of its real oriented blow-up, after rounding corners. Thus the
homotopy type of $L$ is independent of that compactification. In general
the natural map from this end to $U$ is not a homotopy equivalence.

\begin{proof}
If $N$ is open and contains $D$, then $K=X\setminus N$ is a compact
subset of $U$ and $N\setminus D=U\setminus K$. Conversely, every compact
$K\subset U$ is closed in the Hausdorff space $X$, so $N=X\setminus K$
is such a neighbourhood. These assignments are inverse and preserve the
refinement maps. The same proof applies to punctured analytic
neighbourhoods, since these are open subsets of the common complex
manifold $U$.

In a normal-crossings coordinate chart $z_1\cdots z_r=0$, write
$z_j=\rho_j e^{i\theta_j}$. The real oriented blow-up replaces the
vanishing coordinates by $\rho_j\geq0$ and their angles. Its local model
is $[0,\infty)^r\times(S^1)^r\times\mathbb R^{2n-2r}$, so it is a
compact manifold with corners with interior $U$. Rounding corners and
taking a boundary collar gives a cofinal family of end neighbourhoods
$L\times(0,\epsilon)$. All refinement inclusions are homotopy
equivalences, which proves the constant-model assertion. The constant
embedding of the homotopy category into its pro-category is fully
faithful; hence two such models have equivalent homotopy types. Weber's
\S2 identifies the corresponding cohomology directly with
$\varinjlim_K H^*(U\setminus K)$ \citep{Weber2012}; the argument here
retains the whole pro-homotopy type. The explicit failures of $L\to U$
are proved in F.42.
\end{proof}

\Needspace{5\baselineskip}
\paragraph{F.42 (Proved). Both requested examples and the three different invariants.}
\textit{Manuscript: MPL 149; section4; folio66; scan133; additional F(h)(2) tests.}

For $(X,D)=(\mathbb P^1,\{0,1,\infty\})$ over $\mathbb C$, one has
\[
 L\simeq S^1\amalg S^1\amalg S^1,\qquad
 U(\mathbb C)\simeq S^1\vee S^1,
\]
while over trivially valued $\mathbb C$ the underlying topological space
of $U_\infty^{\mathrm{na}}$ is three disjoint open intervals. For
$(X,D)=(\overline M_{0,5},\partial\overline M_{0,5})$ the invariants are
\[
\begin{array}{c|rrrr}
 &b_0&b_1&b_2&b_3\\\hline
 U(\mathbb C)&1&5&6&0\\
 L&1&11&11&1\\
 U_\infty^{\mathrm{na}}\text{ (underlying topology)}&1&6&0&0.
\end{array}
\]
All displayed integral homology groups are free. Both end constructions
are independent of a proper compactification in the precise senses of
F.40--F.41; neither the archimedean link nor the underlying topological
space of the nonarchimedean tube is the homotopy type of $U(\mathbb C)$.
The last row concerns Berkovich topology, not its etale topos.

\begin{proof}
For the three-punctured sphere, small punctured disks retract to three
circles and a pair of pants retracts to a wedge of two circles. A
multiplicative seminorm in the formal punctured disk satisfies
$|\sum_{n\geq m}a_nt^n|=r^m$, where $a_m\ne0$ and $0<r<1$; a unit and
its inverse are bounded, so they have norm one. Thus $r$ parametrizes an
open interval. The three boundary points give the asserted disjoint
union. Its normalized quotient by rescaling valuations is three points.

Normalize three marked points of a smooth five-pointed curve to
$0,1,\infty$. Then
\[
 U(\mathbb C)=\{(x,y)\in(\mathbb C\setminus\{0,1\})^2:x\ne y\}.
\]
Projection to $x$ is a locally trivial bundle, with base the plane minus
two points and fibre the plane minus three points. A local trivialization
is obtained by an isotopy moving $x$ in a small disk disjoint from $0,1$.
Both base and fibre retract to graphs, so the homotopy exact sequence
proves that $U$ is a $K(\pi,1)$. Monodromy preserves each oriented
puncture-meridian in first homology. The integral Serre spectral sequence
therefore has only its $0$ and $1$ rows and columns, all free, with ranks
given by $(1+2t)(1+3t)$. It has no possible differential, yielding the
first row of the table.

The boundary divisors are indexed by two-element subsets of
$\{1,\ldots,5\}$. Two distinct divisors meet exactly when these subsets
are disjoint, as follows by drawing the unique stable three-component
tree with the two prescribed nodes. They intersect transversely once;
there are no triple intersections. Thus the incidence graph is the
Petersen graph, with $10$ vertices and $15$ edges, and its first Betti
number is $15-10+1=6$. Every component is a projective line of
self-intersection $-1$. These geometric descriptions also follow from
the identification with the blow-up of $\mathbb P^2$ in four points,
where the boundary consists of four exceptional curves and the six
transforms of their joining lines; see Dolgachev--Farb--Looijenga,
\S\S2.1--2.2 \citep{DolgachevFarbLooijenga2018}.

For the nonarchimedean row, the normal-crossings skeleton theorem gives
a deformation retraction of the tube onto
$|\Delta(D)|\times\mathbb R_{>0}$: in a chart its coordinates are the
nonnegative numbers $-\log|z_j|$, not all zero, and normalization divides
by their sum. The coordinate retractions glue under etale changes of
normal-crossings charts. This is Thuillier's Theorem~3.25 and
Proposition~4.5 in the opened 2006 preprint \citep{Thuillier2006}.
Consequently the tube has the homotopy type of the Petersen graph.

For the archimedean row, let $N$ be a closed regular neighbourhood of
$D$ in $X$. It is obtained by plumbing the disk normal bundles of the
ten components, and retracts to $D$. It follows that
$H_1(N;\mathbb Z)=\mathbb Z^6$, $H_2(N;\mathbb Z)=\mathbb Z^{10}$,
and $H_3(N;\mathbb Z)=0$. The boundary $L=\partial N$ is a connected
closed oriented three-manifold. Let $Q$ be the intersection matrix of
the ten components. Its diagonal is $-1$ and its off-diagonal entries
are the adjacency matrix of the Petersen graph. Poincare--Lefschetz
duality and the homology exact sequence of $(N,L)$ give
\[
 0\longrightarrow\operatorname{coker}Q\longrightarrow H_1(L;\mathbb Z)
       \longrightarrow\mathbb Z^6\longrightarrow0.
\]
To compute the cokernel integrally, use the basis $H,E_1,\ldots,E_4$ of
$H_2(X;\mathbb Z)$ with form $\operatorname{diag}(1,-1,-1,-1,-1)$.
The boundary classes are $E_i$ and $H-E_i-E_j$. They generate this
unimodular lattice, so their column map $V:\mathbb Z^{10}\to\mathbb Z^5$
is onto. Since $Q=V^t\operatorname{diag}(1,-1,-1,-1,-1)V$, its image is
the primitive rank-five sublattice $\operatorname{im}V^t$. Hence
$\operatorname{coker}Q=\mathbb Z^5$ and $H_1(L)=\mathbb Z^{11}$.
Duality now gives $H_2(L)=\mathbb Z^{11}$ and $H_3(L)=\mathbb Z$.
The attached program independently computes the graph incidence rank
and Smith normal form of $Q$.

Already connectedness rules out $L\simeq U$ in the curve example; in
the surface example the nonzero fundamental class in $H_3(L)$ rules it
out. Neither fact contradicts compactification invariance. The
manuscript's cautious assertion on scan~133 is therefore supported
when ``punctured completion'' denotes the nonarchimedean tube or the
archimedean punctured germ; its identification with the open variety is
not asserted by that sentence and is false in these examples.
\end{proof}


\Needspace{5\baselineskip}
\paragraph{F.43 (Proved). Finite-homotopy exodromy for $\overline M_{0,5}$.}
\textit{Manuscript: Supplemental PartE(c) n=5 question; LM3 F.18 and LM3b H.9; manuscript context MPL147 section12 folio85 scan179.}

Let $X=\overline M_{0,5,\mathbb C}$ with its normal-crossings boundary
stratification, and let $C=h\operatorname{Exit}(X^{\mathrm{an}})$.
The continuous functors from the strict profinite completion of $C$ to
$\pi$-finite spaces classify the etale sheaves of $\pi$-finite spaces on
$X$ constructible for these strata. Stalks and all generization maps are
preserved. Thus the $n=5$ question of LM3b H.9 and Part E(c) is decided;
one must retain the category and its flags, not just its groupoidification.

\begin{proof}
Use
\[
 U=M_{0,5}=\{(x,y):x,y\in\mathbb P^1\setminus\{0,1,\infty\},\ x\ne y\}.
\]
The compactification is the blow-up of $\mathbb P^1\times\mathbb P^1$
at $(0,0),(1,1),(\infty,\infty)$. Forgetting $y$ is a locally trivial
bundle with base the thrice-punctured sphere and fibre the
four-times-punctured sphere: small disjoint discs around the moving
puncture give local trivializations by an isotopy supported in those
discs. The homotopy exact sequence proves that $U$ is aspherical and
that
$1\to F_3\to\pi_1(U)\to F_2\to1$ is exact. All ten one-dimensional
boundary strata are thrice-punctured spheres and the fifteen deepest
strata are points, so every stratum is aspherical.

Let $a,b,c,d,e$ denote the homology classes of meridians of
$x=0,x=1,y=0,y=1,x=y$. They are a basis of $H_1(U,\mathbb Z)$: project
the five complex affine line complements to their defining linear
functions to see independence, and attach the normal disc boundaries
to the simply connected affine plane to see generation. The ten
boundary vectors are
\[
\begin{array}{c|c@{\qquad}c|c}
x_0&a&y_0&c\\
x_1&b&y_1&d\\
x_\infty&-a-b-e&y_\infty&-c-d-e\\
\Delta&e&E_0&a+c+e\\
E_1&b+d+e&E_\infty&-a-b-c-d-e.
\end{array}
\]
The infinity formulas are the product-one meridian relations on a
projective line; an exceptional meridian is the sum of the three
meridians through the blown-up point. The intersecting pairs are
$(x_i,y_j)$ for $i\ne j$, and $(E_i,x_i),(E_i,y_i),(E_i,\Delta)$.
For each of these fifteen pairs the $5\times2$ matrix has a $2\times2$
minor $1$ or $-1$ after an integral change of rows; equivalently the gcd
of its maximal minors is $1$. The attached exact computation lists all
fifteen gcds and obtains Smith invariants $(1,1)$ in each case.
Consequently every local meridian lattice $\mathbb Z^r$, $r\le2$,
embeds as a direct summand of $H_1(U,\mathbb Z)$, hence injects into
$\pi_1(U)$ and, after completion, into $\widehat{\pi_1(U)}$.

For any incident strata $S,T$, let $L$ be the normal link with its two
maps. The exit mapping space is the homotopy fibre of
$L\to S\times T$. Its higher homotopy groups vanish because these
three spaces are aspherical. Its fundamental group is
$\ker(\pi_1L\to\pi_1S\times\pi_1T)$; an element in that kernel is a
normal meridian and is zero by the injection just established (with the
relative-face argument of F.28). Thus every component of every exit
mapping space is contractible. The exit infinity-category is therefore
equivalent to its ordinary category $C$, with composition specified by
the actual three-stratum flags. This is the extra step not furnished by
asphericity of strata alone.

We next justify completion; it is not a formal assertion that profinite
completion preserves all homotopy pullbacks. A finitely generated free
group is good: the cellular two-term resolution of a rose computes its
cohomology with every finite module, and the same free resolution in
continuous modules computes the cohomology of the free profinite group.
The full profinite topology on $F_3$ inside $\pi_1(U)$ is induced by
conjugation into $\operatorname{Aut}(\widehat F_3)$, as proved in F.39.
Thus the profinite completion of the displayed free-by-free sequence is
exact. Comparing its discrete and continuous Hochschild--Serre spectral
sequences proves goodness of $\pi_1(U)$: the cohomology groups of the
fibre with finite coefficients are finite, and only rows $0,1$ occur,
so goodness of the free quotient applies to every term. Links over a
one-dimensional stratum have group $\mathbb Z^r\rtimes F_2$ (the
extension splits since $F_2$ is free), and links over a point have group
$\mathbb Z^r$. Their completions are exact by embedding the normal
lattice into $\widehat{\mathbb Z}^r\rtimes\widehat F_2$; the finite
Koszul resolution of $\mathbb Z^r$ and the same spectral-sequence
comparison prove goodness. Hence all the stratum and link profinite
homotopy types are the classifying types of their profinite groups.
The completed meridian injection now proves, by the same homotopy-fibre
calculation, that every profinite exit mapping object is a profinite
set. Descent on the finitely many full flag diagrams reconstructs a
profinite ordinary EI category with these Hom spaces and compositions.

Here is the finite-congruence point explicitly. A compact totally
disconnected category with finitely many objects is the inverse limit
of its finite categorical quotients. Indeed, choose finite clopen
partitions of its Hom spaces, and put $a\sim b$ when every defined
left and right context $u(-)v$ puts $uav$ and $ubv$ in the same partition
piece. Joint continuity and compactness of the context spaces imply
that this equivalence is open: finitely many neighbourhoods of contexts
suffice to find one neighbourhood of $a$ on which all their partition
values are fixed. Compactness then makes the number of classes finite.
Inserting further contexts shows that the equivalence is compatible
with composition. Thus finite congruences are cofinal among the finite
clopen data. In our case their restrictions to the dense discrete
category $C$ are precisely the strict finite congruences used in LM3
F.2--F.4. Here one also uses the topology assertion in F.39 and the
completed link descriptions above: finite quotient data on the
discrete stratum groups and incident cosets extend continuously to
their completions, and the finitely many composition constraints
extend by density. Density alone would not imply this extension
property. Thus the inverse system is the strict profinite completion.

The same statement works with finite homotopy values, not only finite
sets. Given finitely many $\pi$-finite values, take a degreewise finite
simplicial model of the full infinity-category they span; its mapping
spaces have one common finite truncation bound. Its coherent diagram
data are determined by finitely many simplicial degrees, using a
coskeletal model above that bound. Continuous maps in those degrees
factor through finite clopen partitions of finite products of Hom
spaces. Every such partition admits a finite coordinatewise clopen
refinement: compactness gives a finite cover by clopen rectangles,
and the finitely many possible object sequences cause no additional
infinite choice. Refine the resulting individual Hom partitions by
the context congruence of the preceding paragraph. Thus one finite
categorical congruence suffices for all the diagram data.
The same argument with an additional finite simplex
$\Delta^m$ proves the assertion on mapping spaces, for every $m$.
Consequently continuous representations in $\pi$-finite spaces are the
infinity-categorical filtered colimit of their finite-stage
representations; this is stronger than a statement about finite sets.

Finally, the algebraic-to-analytic comparison on these strata and their
normal links identifies their finite covers and their finite-coefficient
higher descent data. In the notation of Barwick--Glasman--Haine, this
is Proposition~12.6.4 and Corollary~12.6.7, followed by the exodromy
Theorem~10.1.8 \citep{FA_BGH2020}. The preceding paragraphs verify in this
particular geometry that the resulting profinite exit object is the
strict completion of the ordinary category, rather than a profinite
infinity-category with hidden higher mapping spaces. Applying the
comparison and exodromy therefore gives precisely the asserted
stalk-preserving equivalence. The classifying space of this category
still has higher homotopy; for example $H^2(X,\mathbb Z)=\mathbb Z^5$
and $H^4(X,\mathbb Z)=\mathbb Z$. Ordinary exit mapping spaces being
discrete does not mean that this classifying space is a $1$-type.
\end{proof}

\emph{Two explicit completion checks for the preceding argument.}
The flag description can be made completely concrete here. The only
strict flags with three strata are
$p<S<U$, where $p$ is a boundary crossing and $S$ is a
thrice-punctured boundary component. The normal circle bundle over
$S$ is topologically trivial, since $H^2(S,\mathbb Z)=0$.
Consequently its link group is $G_S\times\mathbb Z$ and its map to
$G_S$ is the projection. The other adjacent link has group
$\mathbb Z$, mapping to $G_S$ by the relevant puncture loop. The
three-stratum flag space is their homotopy pullback, hence has group
\[
 \mathbb Z\mathbin{\times}_{G_S}(G_S\times\mathbb Z)
                         \simeq\mathbb Z^2.
\]
After profinite completion the same identity holds, with hats: the
projection is split and all these products carry the full induced
profinite topology. Thus the finitely many nontrivial Segal gluing
identities survive completion in this geometry. This supplies the
hypothesis needed to apply the spatial and profinite decollage
equivalences of \cite[Theorem~2.7.4 and equivalence~(2.8.9)]{FA_BGH2020};
no preservation of arbitrary homotopy pullbacks by completion is
asserted.

It also verifies universality for \emph{strict} finite categorical
quotients, which does not follow from density alone. Choose one point
in each stratum. For each incident pair $S<T$, the link is connected,
so its Hom set is a transitive $(G_T,G_S)$-biset. Upon choosing one
arrow its stabilizer is the image of its link group $H_{S,T}$ in
$G_T\times G_S$ (using the inverse right action). The injectivity
already proved gives the completed Hom space
\[
  (\widehat G_T\times\widehat G_S)
           /\operatorname{im}(\widehat H_{S,T}).
\]
The denominator is compact and therefore closed. A functor from the
discrete exit category to a finite category restricts to homomorphisms
from each $G_S$ to finite automorphism groups; these extend to
$\widehat G_S$. Its map on an incident Hom set is determined by the
image of one arrow and the finite group actions. The stabilizer
condition extends to the closure, so this map extends uniquely and
continuously to the displayed profinite Hom space. Composable pairs
of discrete arrows are dense in the corresponding product of Hom
spaces. Hence the extended maps preserve composition, since the
two resulting continuous maps to a finite discrete set agree on that
dense subset. This proves that every finite categorical quotient
extends, and hence that the compact category constructed from the
completed flags has precisely the strict finite quotients of $C$.
The context-congruence argument then identifies it with their inverse
limit.

Finally, the finite bound used for coherent representations has an
exact model-theoretic justification. The full coefficient category on
finitely many $\pi$-finite spaces has finitely many object classes,
finite homotopy groups in every mapping space, and a common finite
truncation bound on those mapping spaces. A minimal quasicategory
model is degreewise finite by Blom--Moerdijk, Lemma~A.10, and has a
finite coskeletality bound by Lurie, Proposition~2.3.4.18 and
Remark~2.3.4.4; this is exactly the conclusion recorded in
Blom--Moerdijk, Remark~A.11
\citep{FA_BlomMoerdijk2023,FA_LurieHTT2009}.
One takes this equivalent minimal model before using finite coherent
data; no strict coskeletality assertion about an arbitrary homotopy
coherent nerve is necessary. Thus a continuous coherent diagram is
determined by finitely many clopen data on products of Hom spaces.
The bound is taken separately for each finite-object diagram and for
each finite simplex of its mapping space, rather than for the whole
class of all $\pi$-finite spaces.

\Needspace{5\baselineskip}
\paragraph{F.44 (Open). Delimitation beyond the proved reconstruction ranges.}
\textit{Manuscript: MPL149; section2; folio25v; scan52; intermediate-base delimitation.}

The requested base-field delimitation now includes the following
proved positive ranges: infinite finitely generated fields by F.5,
sub-$p$-adic fields and the specified general bases by F.38, and
every algebraic subfield of $F\mathbb Q^{\mathrm{ab}}$, for a
number field $F$, by F.47, and, more generally, every subfield of
a finitely generated extension of $\mathbb Q^{\mathrm{ab}}$ by
F.50. These assertions include arbitrary cyclotomic extensions of
number fields and fields finitely generated over them. In particular, the test
over $\mathbb Q(\mu_{2^\infty})$ considered during this work is
settled, including full site realization. F.46 records its initial
point-function proof and F.47 supplies the additional argument
needed for the full geometric morphism. The proposed unrestricted
fullness is false by F.5. F.49 strengthens the counterexamples
F.6 and F.37: pinned fullness fails over every algebraic extension
of a finite field, every algebraically closed field, and every
infinite perfect field with nontrivial centre of its absolute
Galois group, including every real closed field.
Faithfulness of actual reduced-scheme maps remains
the all-characteristic theorem F.4.

The exact general criterion is Carlson--Haine--Wolf,
Theorem~4.18 \citep{CHW2024}: a field has pinned reconstruction if
and only if every pinned morphism from a connected regular affine
scheme with a rational point to $(\mathbb G_m)_{\mathrm{et}}$
has the geometric-point values of a regular unit (on the absolute
weak normalization in positive characteristic). Its implication to
full site realization is proved through all sliced etale objects,
as explained in F.47. Their Theorem~6.24 further supplies the
positive range of amenable fields admitting a nontrivial discrete
valuation. Amenability means that rational points of every torus
have no nonidentity prime-to-characteristic divisible element and
that the prime-to-characteristic Tate module of the Picard group of
every regular finite-type scheme vanishes (Definition~5.15).
This is an exact verified sufficient condition, not a claim that
every field outside it is a failure or a literature-wide open case.

A concrete test not decided by this return is
\[
 k_{\mathrm{solv}}=\mathbb Q^{\mathrm{solv}},\qquad
 \operatorname{Hom}_{k_{\mathrm{solv}}}(\mathbb G_m,\mathbb G_m)
 \longrightarrow
 \pi_0\operatorname{Hom}^{\mathrm{pin}}_{B G_{k_{\mathrm{solv}}}}
       ((\mathbb G_m)_{\mathrm{et}},(\mathbb G_m)_{\mathrm{et}}),
 \tag{F44}
\]
where $\mathbb Q^{\mathrm{solv}}$ is the maximal solvable algebraic
extension in a fixed algebraic closure. The map is injective by
F.4; its surjectivity is not asserted. Finite-type dimension-zero
schemes, after universal-homeomorphism localization, are finite
etale and their relative Galois topoi already determine their maps.
Thus (F44) is a first positive-dimensional test of the remaining
criterion, rather than a claim of a smallest base field or of
literature-wide minimality. Constant morphisms and the general
nonconstant, possibly noninvertible, locus are part of the question;
the unresolved realization issue concerns the latter locus.

Here are the precise failed extensions of the proof, rather than
an appeal to the absence of a cited theorem. The field
$k_{\mathrm{solv}}$ is radically closed. Indeed, an element lies
in a finite solvable Galois extension $E/\mathbb Q$. Adjoin the
needed roots of unity and then $n$th roots of all its conjugates;
the latter extension is abelian over the enlarged $E$, so its
normal closure is still solvable over $\mathbb Q$. Consequently
\[
       D(k_{\mathrm{solv}}^*)=k_{\mathrm{solv}}^*.
\]
The height argument used in F.47 cannot reduce these elements to
roots of unity: this field contains $2^{1/n}$ for every $n$, of
positive heights tending to zero. The same observation shows that
it cannot be contained in $F\mathbb Q^{\mathrm{ab}}$ for any
fixed number field $F$, by the verified Amoroso--Zannier bound.
It is also outside the broader range F.50: the algebraic constants
in a finitely generated extension of $\mathbb Q^{\mathrm{ab}}$
form a finite extension of that field, to which the same height
bound applies.
It is not amenable. It is not sub-$p$-adic either, since it contains
all roots of unity whereas the algebraic constants in a finitely
generated extension of $\mathbb Q_p$ form a finite extension and
have only finitely many roots of unity.

The Picard input of F.47 also fails here. For an elliptic curve
$E/\mathbb Q$, every $2^r$-division field is solvable: its Galois
group is a subgroup of
$\operatorname{GL}_2(\mathbb Z/2^r)$, an extension of
$\operatorname{GL}_2(\mathbb F_2)\simeq S_3$ by a $2$-group.
Hence $E[2^\infty]\subset E(k_{\mathrm{solv}})$, providing
nonzero infinitely divisible points and a nonzero Tate module.
Neither the integral-residue argument followed by Picard
divisibility nor the translated-function proof can therefore be
copied with its hypotheses unchanged. Kummer classes at the
$k_{\mathrm{solv}}$-points already forget all nonzero constants.

The relative birational isomorphism results checked in
Tsujimura's Theorem~4.7 and Corollary~4.10 \citep{Tsujimura2022}
do not supply this noninvertible realization step; the requisite
AVKF hypotheses also fail in this example. Descent of a scheme to
a number field does not imply descent of an arbitrary geometric
morphism on the whole etale site. No exotic pinned morphism in
(F44) has been constructed here, and no argument using the
nonabelian finite covers has yet eliminated one. The next step
would have to exploit those nonsolvable-cover compatibilities,
rather than the Kummer and Picard invariants just shown to vanish
or to have the wrong divisibility. This is the exact residual
test and obstruction left by this return; no broader assertion
about the status of all other fields is made.

A further literature check does not settle (F44).
Mochizuki's \emph{Global Solvably Closed Anabelian Geometry}
\citep{Mochizuki2006SolvablyClosed}, Definitions~1.1 and~2.3 and
Theorem~2.5, concerns quotient groups
$Q_F=\operatorname{Gal}(\widetilde F/F)$ for a solvably closed
extension of a number field, together with the required
$\Sigma$-compatible quotients of hyperbolic-curve fundamental
groups. Its conclusion reconstructs compatible isomorphisms.
The absolute group $G_{\mathbb Q^{\mathrm{solv}}}$ in (F44)
is instead a kernel in $G_{\mathbb Q}$, and (F44) asks about
all pinned morphisms, including noninvertible ones. Neither
changing the quotient to this kernel nor replacing the Isom
conclusion by Hom is justified by that theorem.

\medskip
\noindent\emph{Finite extensions do not repair these particular inputs.}
One possible approach would pass to a proper finite extension
$L/k_{\mathrm{solv}}$, which is Hilbertian by the Weissauer theorem
used explicitly in the proof of \cite[Proposition~2.5]{F44Jarden2011}.
Descent of a reconstruction theorem would indeed be valid here.
More generally, in characteristic zero pinned reconstruction
descends along a finite extension $L/k$. To prove this, apply the
criterion of \cite[Theorem~4.18]{CHW2024} to a connected regular
affine $X/k$ with a rational point and a pinned morphism to
$(\mathbb G_m)_{\mathrm{et}}$. The curve-free criterion applies
after finite etale base change to $L$, where it gives a unit
$f_L$ with the prescribed geometric-point values. On
$X_{L\otimes_kL}$ the two pullbacks of $f_L$ have the same values
at every geometric point, by the descent data of the original
morphism. This scheme is reduced, so the units agree.
Faithfully flat descent gives $f\in\Gamma(X,\mathcal O_X)^*$
with the required values over $k$; the same criterion proves
reconstruction. The source remains connected after base change:
regularity in characteristic zero makes it normal and smooth,
and its rational point forces its finite algebraic constant field
to be $k$.

The obstacle is therefore the arithmetic input over $L$, not this
descent argument. For \emph{every algebraic} extension
$L/k_{\mathrm{solv}}$, including each proper finite one, one has
\[
 k_{\mathrm{solv}}^*\subset D(L^*),\qquad
 E[2^\infty]\subset E(L).
\]
Thus $L^*/\mu_\infty$ contains the nonzero divisible subgroup
$k_{\mathrm{solv}}^*/\mu_\infty$ and cannot be free abelian;
the same elliptic curve still violates the AVKF and Tate-module
conditions. Nor can such an $L$ admit a nontrivial discrete
valuation. A homomorphism $v:L^*\to\mathbb Z$ arising from a
valuation vanishes on the divisible subgroup $k_{\mathrm{solv}}^*$.
For any element algebraic over this trivially valued field, its
minimal polynomial has all nonzero coefficients of value zero.
A positive or negative value of the element would make the
minimum value among its nonzero polynomial terms occur exactly
once, which is impossible. Hence $v$ vanishes on $L^*$.
Hilbertianity of a finite extension therefore does not restore
the multiplicative, Picard, or discrete-valuation hypotheses used
in the reconstruction proofs verified here. This eliminates that
specific route without producing an exotic pinned morphism or
deciding (F44).


\medskip
\noindent\emph{Additional audit of the residual test
$k=\mathbb Q^{\mathrm{solv}}$.}
Neither of the elementary nongeometric constructions in F.49 settles
this test. In fact, with $G=G_{\mathbb Q}$ and
$H=G_{\mathbb Q^{\mathrm{solv}}}$, one has
\[
                         C_G(H)=1.                         \tag{F44.C1}
\]
We give the proof, verifying the precise scope of the recovered
\texttt{CYCLOTOMIC\_KERNEL} row B2 instead of importing a stronger
unqualified assertion about arbitrary closed subgroups.

The following argument works for every Hilbertian field $F$ and
every nontrivial closed normal subgroup $N\triangleleft G_F$.
Two primary results suffice. Jarden's Proposition~2.5 says that if
$E/F$ is algebraic Galois and $E$ is not separably closed, then
$G_E$ is almost locally free; in particular it contains a free
profinite subgroup of rank two
\citep{F44Jarden2011}. Haran--Jarden, Corollary~2.5, says that
$F^{\mathrm{sep}}$ cannot be the compositum of two proper subfields
which are Galois over $F$ \citep{F44HaranJarden1991}. Both statements
allow infinite Galois extensions.

There is no nontrivial closed abelian normal subgroup $A$ of $G_F$:
its fixed field is Galois over $F$ and is not separably closed, so
the first result would place a nonabelian free profinite rank-two
group inside $A$. Put $C=C_{G_F}(N)$. It is closed and normal,
and $N\cap C=Z(N)$ is a closed abelian normal subgroup of $G_F$;
hence $N\cap C=1$. If $C\ne1$, the two proper fixed fields of
$N$ and $C$ would be Galois over $F$, and their compositum would
have absolute Galois group $N\cap C=1$. This contradicts the
second result. Thus $C_{G_F}(N)=1$.

To apply it to $F=\mathbb Q$, we check that
$\mathbb Q^{\mathrm{solv}}\ne\overline{\mathbb Q}$.
The polynomial $P=X^5-X-1$ is irreducible modulo $3$: the attached
exact program divides it by all twelve monic polynomials of degree
one or two and finds no divisor. A reducible quintic must have a
factor of one of these degrees. Modulo $2$ its squarefree
factorization is
\[
           (X^2+X+1)(X^3+X^2+1),
\]
with both factors irreducible. Its rational Galois group is therefore
transitive of degree five and contains a permutation of cycle type
$(2)(3)$, whose cube is a transposition. A transitive group of prime
degree is primitive. The graph whose edges are the conjugates of
that transposition has components forming blocks, and is therefore
connected. Its edge transpositions generate $S_5$, so the Galois
group is $S_5$. The program also computes its derived-subgroup
orders $120,60,60$, certifying nonsolvability. Its splitting field
cannot lie in the maximal solvable extension. Thus $H\ne1$, and
Hilbert irreducibility for $\mathbb Q$ gives (F44.C1).
In particular $Z(H)=1$, so F.49 supplies no nonidentity central
comparison twist over this field.

The semilinear field-automorphism construction also does not apply.
Since $k/\mathbb Q$ is Galois,
$\operatorname{Aut}_{\mathbb Q}(k)=G/H$. For an automorphism
$\sigma$ choose a lift $g\in G$. Its outer action on $H$ is
conjugation by $g$. If this action were inner, say equal to
conjugation by $h\in H$, then
\[
                       h^{-1}g\in C_G(H)=1.
\]
Hence $g=h$ and $\sigma=1$. A nonidentity coefficient-semilinear
automorphism of $\mathbb A^1_k$ consequently cannot be supplied with
the relative comparison used in F.49: such a comparison, restricted
along the rational point $0$, would make its base automorphism of
$BH$ two-isomorphic to the identity and thus its action on $H$
inner. This excludes that construction; it makes no assertion that
all pinned topoi morphisms over $k$ have been classified.

The finite certificate is
\path{code/montpellier/F44_solvable_certificates.py}; its exact output is
included with the deliverable.


\Needspace{5\baselineskip}
\paragraph{F.45 (Open). The remaining literal FLAGGED\_149 audit.}
\textit{Manuscript: Supplemental EXIT task6 source audit; manuscript context MPL149 section4 folios65v-67 scans132-135.}

The required literal comparison with the \texttt{FLAGGED\_149}
register rows F11--F16 and F25 remains incomplete. No exact matching
return was found in the supplied files, the targeted title and content
searches, or the complete 749-item inventory of accessible ZIP
filenames. That inventory does not exclude an unindexed copy nested
inside another archive. The actual register and its proofs are
therefore still needed to verify the intake's row numbers and claims.

Two earlier source gaps were closed during this work. The
\path{LM8_LIE_APPROXIMATIONS_background.zip} archive contains
the actual \texttt{CYCLOTOMIC\_KERNEL} return and the \S53 scans and
transcriptions. F.48 audits its exact rows B5 and C8 and proves the
literal sector-functor comparison with \S53. Earlier claims that
these particular sources had not been located are superseded by
that recovery. The initial unsuccessful searches are retained only
as a record of the work, not as current descriptions of availability.

The EXIT rows P3b-discrete, P3b-completion,
P3b-horizontal-profinite, P3b-vertical,
P3b-vertical-counterexample and P3b-profinite were inspected with
their proofs; F.39 gives the relevant stable-tower theorem and its
precise contact-matrix and cyclotomic qualifications. F.22 refutes
the separate assertion that finite root stacks have one fixed
uncompleted inertia kernel. These independent mathematical results
do not authenticate the unlocated FLAGGED register. Only that
literal cross-return comparison remains open in this source-audit
row; F.48 closes the previously recorded CYCLOTOMIC and \S53 gaps.


\Needspace{5\baselineskip}
\paragraph{F.46 (Proved). The point function of every cyclotomic test morphism.}
\textit{Manuscript: MPL149; section2; folio25v; scan52; prompt Part F remaining-field test.}

Let $k=\mathbb Q(\mu_{2^\infty})$, fix $K=\overline{\mathbb Q}$, and let
\[
 \phi:(\mathbb G_{m,k})_{\mathrm{et}}
       \longrightarrow(\mathbb G_{m,k})_{\mathrm{et}}
\]
be a pinned geometric morphism over $B G_k$, with the relative
identification included in the data. Its induced map on geometric
points is
\[
 f:K^*\longrightarrow K^*,\qquad f(a)=c a^m
 \quad(c\in k^*,\ m\in\mathbb Z).
 \tag{F46}
\]
Thus the point-function obstruction in the cyclotomic test over
$\mathbb Q(\mu_{2^\infty})$ vanishes. This assertion by itself does not identify the full geometric
morphism: the hypothesis of Carlson--Haine--Wolf,
Proposition~4.14, concerns every quasicompact separated etale target
object and its pullback, not only the terminal object.

\emph{Problem locus.} Montpellier cote149, \S2, fol.25v, scan52,
and the supplied prompt's Part~F intermediate-field reconstruction question.
This develops the arithmetic reconstruction question there.
It is a new consequence of the proved elementary arguments
below and the explicitly cited arithmetic input; it is not attributed
to a sentence of the manuscript.

\emph{Proof: arithmetic facts.}
For a field $L$ put $D(L^*)=\bigcap_{n\geq1}(L^*)^n$.
For every finite extension $L/k$ one has
\[
 D(L^*)=\mu_{2^\infty}.                                      \tag{F46.1}
\]
Here is the precise arithmetic input and its application.
Amoroso--Zannier, Theorem~1.1, gives a positive lower bound for the
height of a non-torsion algebraic number in terms of its degree over
an abelian extension of a fixed number field
\citep{AmorosoZannier2000}. Apply it with the abelian extension
$\mathbb Q^{\mathrm{ab}}/\mathbb Q$. The degrees of elements of
$L\mathbb Q^{\mathrm{ab}}$ over $\mathbb Q^{\mathrm{ab}}$ are bounded
by $[L:k]$. Consequently the heights of non-torsion elements of $L^*$
have a uniform positive lower bound. If $a\in D(L^*)$ were not a root
of unity, choosing $b_n\in L$ with $b_n^n=a$ would give
$h(b_n)=h(a)/n>0$, a contradiction. Thus $a$ is torsion.

The field $k$ is unramified at each odd prime $\ell$: every finite
subextension is contained in a $2$-power cyclotomic extension, whose
defining roots of unity are etale over $\mathbb Z_\ell$. The value
group of any extension of the $\ell$-adic valuation to $L$ therefore
has index at most $[L:k]$ over $\mathbb Z$. In particular, the orders
of the $\ell$-primary roots of unity in $L$ are bounded, since
\[
 v_\ell(\zeta_{\ell^r}-1)
       =\frac{1}{\ell^{r-1}(\ell-1)}.
\]
If a torsion element of $D(L^*)$ had nontrivial $\ell$-primary part,
its $\ell^n$-th roots in $L$ would contradict this bound. Only its
$2$-primary part remains. Conversely $\mu_{2^\infty}$ is divisible
by every positive integer and is contained in $k$, proving (F46.1).

We also use that $K/k$ has infinite degree, and that every finite
extension of $k$ admits a discrete valuation. Choose a valuation over
$3$ with $v(3)=1$. Its value group on $k$ is $\mathbb Z$; on a finite
extension it has finite index over $\mathbb Z$ and can be renormalized
to $\mathbb Z$. Finally $T^n-3$ is Eisenstein for this valuation, so
$[K:k]$ is unbounded. Equivalently, a root has valuation $1/n$, which
forces the degree of its extension of $k$ to be at least $n$.

\emph{Proof: a rational-source Kummer lemma.}
Let $L/k$ be finite, and let
\[
 U=\mathbb P^1_L\setminus\{a_1,\ldots,a_r,\infty\},
 \qquad a_i\in L\text{ distinct}.
\]
Suppose a compatible system of Kummer classes
$\alpha\in\varprojlim_n H^1(U,\mu_n)$ has the following property:
at every $x\in U(L')$, for every finite $L'/L$, its evaluation is the
Kummer class of a specified element $F(x)\in(L')^*$.
Then there is a regular unit $g\in\mathcal O(U)^*$ such that
\[
 \alpha=[g],\qquad F(x)/g(x)\in D((L')^*)
       \quad(x\in U(L')).                                  \tag{F46.2}
\]
Indeed, $\operatorname{Pic}(U)=0$, and unique factorization gives
\[
 \alpha=\widehat c\prod_{i=1}^r(t-a_i)^{\lambda_i},\qquad
 \widehat c\in\widehat{L^*},\quad\lambda_i\in\widehat{\mathbb Z}.
\]
Choose a discrete valuation $v:L^*\to\mathbb Z$ and an element
$\pi$ of value one. For fixed $j$, put $x_N=a_j+\pi^N$.
For all sufficiently large integers $N$, these are points of $U(L)$,
$v(x_N-a_j)=N$, and
$v(x_N-a_i)=v(a_j-a_i)$ for $i\ne j$.
Evaluating each Kummer quotient and subtracting the equations for
$N+1$ and $N$ yields
\[
 \lambda_j=v(F(x_{N+1}))-v(F(x_N))
       \quad\text{in }\widehat{\mathbb Z}.
\]
Its right side is an integer, so write $\lambda_j=m_j\in\mathbb Z$.
Choose $x_0\in U(L)$ and set
\[
 c=F(x_0)\Big/\prod_i(x_0-a_i)^{m_i}.
\]
Evaluation at $x_0$ identifies $\widehat c$ with the completion of
this actual $c\in L^*$. Thus $g=c\prod_i(t-a_i)^{m_i}$ proves
(F46.2), including its assertion after every finite extension.
No finite-presentation descent for the topos morphism is used here.

\emph{Proof: apply the lemma to every translate.}
A relative pinned morphism sends $L$-rational points to
$L$-rational points, for every finite $L/k$. One way to check this
is to compose the point's section $G_L\to\pi_1$ with the given map.
If its image is a closed point with residue field $E$, the resulting
map into $G_E$ lies over the inclusion $G_L\subset G_k$; hence
$G_L\subset G_E$ and the specified embedding of $E$ lies in $L$.
The induced function $f$ is therefore $G_k$-equivariant and has the
stated finite-extension compatibility. In particular $f(1)\in k^*$;
divide by this constant and assume from now on that $f(1)=1$.

If the generic point maps to a closed point, continuity and pinning
make $f$ constant, necessarily equal to $1$, and (F46) follows with
$m=0$. Otherwise its image is generic. For each $b\in K$, after a
finite extension containing $b$, pull back the open
$\mathbb G_m\setminus\{b\}$ of the target. This is a nonempty
open $U_b$ of the source, with finitely many missing points. Its
geometric points are exactly the $x$ with $f(x)\ne b$.
After a further finite extension $L/k$ all the missing points are
rational. Pullback of the compatible Kummer classes of $t-b$ to
$U_b$ satisfies the preceding lemma, with $F(x)=f(x)-b$.
Thus there is $g_b\in L(t)^*$, a unit on $U_b$, such that
\[
 f(x)-b=d_b(x)g_b(x),\qquad d_b(x)\in\mu_{2^\infty}\subset k^*,
       \quad x\in U_b(K).                                  \tag{F46.3}
\]
Each $x$ is contained in a finite extension of $L$, so (F46.1)
applies. For $b=0$ the domain is all of $\mathbb G_m$. The lemma
and the normalization at $1$ give
\[
 f(x)=d_0(x)x^m,\qquad m\in\mathbb Z,\quad d_0(x)\in k^*.
                                                               \tag{F46.4}
\]
The stronger assertion $d_0(x)\in\mu_{2^\infty}$ will not be needed
until the end.

\emph{Proof: the case $m=0$.}
In this case all values of $f$ lie in $k$. Choose $b\in K\setminus k$.
Then $U_b=\mathbb G_m$, and (F46.3) gives
$f(x)-b=d_b(x)c_bx^n$ for an integer $n$ and $c_b\in K^*$.
If $n\ne0$, the values of $c_bx^n$ fill $K^*$; the equality would
place all these values in the $k$-vector space $k+kb$, of dimension
at most two. This contradicts $[K:k]=\infty$.
Therefore $n=0$. If $f$ took two distinct values, subtracting their
two equations would give $c_b\in k$ and then $b\in k$, again a
contradiction. So $f$ is constant, and $f(1)=1$ gives $f=1$.
This treats the zero exponent without applying a nonconstant
quasi-rationality criterion to it.

\emph{Proof: a finite-exception quasi-rationality lemma.}
We give the argument needed to use (F46.3) safely at its missing
points. It is the finite-exception form of the mechanism in
Tsujimura, Lemmas~3.3--3.5 and Proposition~3.6
\citep{Tsujimura2022}.
Suppose $K$ is algebraically closed of characteristic zero, algebraic over a subfield $M$,
and $[K:M]=\infty$. Let $C/K$ be a smooth proper connected curve,
let $q$ be a function on a cofinite subset of $C(K)$, and suppose
that for each $b\in K$ one has
\[
 q(x)-b=e_b(x)r_b(x),\qquad e_b(x)\in M^*,\quad r_b\in K(C)^*,
                                                               \tag{F46.5}
\]
outside a finite set depending on $b$. Then $q$ agrees with one
rational function on a cofinite subset. If $r_0$ is constant, this
rational function is constant as well.

First suppose $r_0$ is nonconstant and apply (F46.5) to $q$ and
$1-q$. Write their rational factors as $r,s$. Choose a finite
extension $N/M$ in $K$ over which $C,r,s$ and a finite nonconstant
map $C\to\mathbb P^1$ are defined. For fixed $A,B\in M^*$ the
degree of the pole divisor of $Ar+Bs-1$, if it is nonzero, is
bounded by one integer $d$, independent of $A,B$. Thus a zero
of this nonzero rational function has residue degree at most $d$
over $N$. Choose $x$ outside the exceptional sets whose residue
field over $N$ has degree greater than $d$. Such points exist:
choose a value of the fixed finite map of degree greater than
$d$ over $N$, and any point above it. The values, and hence the
points, can avoid any given finite set because $K/N$ has infinite
degree. Substituting $A=e_0(x)$ and the corresponding coefficient
$B$ of $1-q$ shows that $Ar+Bs-1$ vanishes at $x$.
The degree bound forces the rational-function identity $Ar+Bs=1$.
Rescale $r,s$ by these constants so that $r+s=1$.
The pointwise equation then gives
\[
 (e_0(x)-e_1(x))r(x)=1-e_1(x).
\]
Consequently $q(x)=r(x)$ wherever $r(x)\notin M$, outside the
finite exceptional set.

Choose $a\in K\setminus M$ and repeat this argument for
$q+a$ and $1-q-a$. Its rational factor $r_a$ is nonconstant:
if it were a constant $c$, the equation
$q+a=e_a c$, together with $q=e_0r_0$, would place all cofinite
values of the nonconstant rational function $r_0$ in the
finite-dimensional $M$-vector space $Mc+Ma$. A nonconstant
rational function has cofinite image, so this is impossible.
After rescaling we obtain $q(x)=r_a(x)-a$ wherever
$r_a(x)\notin M$.
The two sets $r(x)\notin M$ and $r_a(x)\notin M$ cover the
cofinite domain under consideration: if both rational values were
in $M$, (F46.5) would put both $q(x)$ and $q(x)+a$ in $M$.
Their intersection is infinite. To check this last point without
a cardinality assumption, enlarge $N$ by their finitely many
coefficients. Points of sufficiently large degree over $N$ have
neither rational value in $M$, since a fibre of either fixed
nonconstant rational function has bounded degree. Infinitely many
such points exist. On this intersection $r=r_a-a$; equality of
rational functions follows. Their common value therefore agrees
with $q$ outside a finite set.

If $r_0=c$ is a nonzero constant, divide $q$ by $c$; the hypothesis
for all translates is preserved, and $q(x)\in M$ outside a finite
set. Choose $b\in K\setminus M$. Its rational factor $r_b$ cannot
be nonconstant, because its cofinite image would be contained in
$M+Mb$ by (F46.5). Hence $r_b$ is constant. If $q$ took two distinct
values on the common cofinite domain, subtracting the equations
$q(x)-b=e_b(x)r_b$ would put $r_b$ and then $b$ in $M$.
Thus $q$ is constant there. This proves the lemma in both cases,
for an arbitrary smooth proper curve, as needed also in F.47.

\emph{Proof: remove the exceptional values and finish.}
For $m\ne0$, apply the lemma with $M=k$, $q=f$, and $r_0=t^m$.
Let $R\in K(t)$ be the resulting nonconstant rational function,
equal to $f$ outside a finite set. For each $b$, equation (F46.3)
implies that
\[
 (R-b)/g_b
\]
takes values in $k^*$ on a cofinite set. A nonconstant rational
function has cofinite image in $K$, whereas $k$ is a proper
subfield; hence this quotient is a constant $c_b\in k^*$.
Thus $R-b=c_b g_b$ as rational functions.

Fix any $x\in K^*$. Taking $b\ne f(x)$ ensures $x\in U_b$,
where $g_b$ is a regular unit. The rational identity just proved
first shows that $R$ has no pole at $x$. For every
$b\notin\{f(x),R(x)\}$ it then gives
\[
 \frac{f(x)-b}{R(x)-b}\in k^*.
\]
If $f(x)\ne R(x)$, the displayed fractional linear expression,
as $b$ varies, takes every value in $K\setminus\{0,1\}$.
This contradicts $k\ne K$. Therefore $f(x)=R(x)$ for every
$x\in K^*$, including the initially exceptional points.
In particular $R$ is a regular unit on $\mathbb G_{m,K}$, so
$R=c t^j$. Alternatively, (F46.4) shows directly that
$R/t^m$ is a rational function whose cofinite values lie in $k$;
it must be constant. Normalization at $1$ makes that constant one.
Undoing the initial normalization gives exactly (F46).
\hfill$\square$


\Needspace{5\baselineskip}
\paragraph{F.47 (Proved). Pinned reconstruction over cyclotomic algebraic bases.}
\textit{Manuscript: MPL149; section2; folio25v; scan52; intermediate-base reconstruction.}

Let $k\subset\overline{\mathbb Q}$ be any subfield contained in
$F\mathbb Q^{\mathrm{ab}}$ for some number field $F$. For all
finite-type $k$-schemes $X,Y$, the natural map
\[
 \operatorname{Hom}_k(X^{\mathrm{awn}},Y)
   \longrightarrow
 \pi_0\operatorname{Hom}^{\mathrm{pin}}_{B G_k}
                 (X_{\mathrm{et}},Y_{\mathrm{et}})
\]
is a bijection, where a relative geometric morphism includes its
commutation isomorphism over $B G_k$. Equivalently, the etale-topos
functor is fully faithful on finite-type schemes after inverting
universal homeomorphisms, with pinned geometric morphisms as the
target morphisms. In characteristic zero the absolute weak
normalization is seminormalization. In particular this settles the
test over $\mathbb Q(\mu_{2^\infty})$ in full, including constants
and noninvertible morphisms, and applies to $\mathbb Q^{\mathrm{ab}}$
and all its finite extensions. We give the derivation from the
verified arithmetic theorems and the reconstruction criterion;
no priority assertion is intended.

\begin{proof}
\subsubsection*{An integrality lemma for cyclotomic algebraic bases}

Let $k\subset\overline{\mathbb Q}$ be contained in
$F\mathbb Q^{\mathrm{ab}}$ for a number field $F$, and let $L/k$ be
finite. Let $b_1,\ldots,b_r\in L$ be distinct and put
$V=\mathbb A^1_L\setminus\{b_1,\ldots,b_r\}$.
Suppose
\[
 \xi\in\varprojlim_n H^1_{\mathrm{et}}(V,\mu_n)
\]
has the following evaluation property: for every $a\in V(L)$,
$a^*\xi$ is the Kummer class of an element of $L^*$.
Then $\xi$ is the Kummer class of an actual unit of $V$.
This is an integrality assertion for a cohomology class; the individual
elements realizing its evaluations need not initially form a rational
function.


Since $\operatorname{Pic}(V)=0$ and
$\mathcal O(V)^*=L^*\times\prod_i(t-b_i)^{\mathbb Z}$, Kummer theory gives
\[
 \xi=\gamma+\sum_{i=1}^r\lambda_i[t-b_i],\qquad
 \gamma\in\widehat{L^*},\quad \lambda_i\in\widehat{\mathbb Z}.
                                                        \tag{F47.R1}
\]
All completions in this argument are $\varprojlim_n A/nA$ (in
multiplicative notation, $\varprojlim_n A/A^n$). We first prove
$\lambda_i\in\mathbb Z$ without assuming that $L$ has a discrete
valuation.

Choose a number field $K_0\subset L$ containing all the $b_i$.
For each $i$, choose a different prime ideal $\mathfrak p_i$ of
$\mathcal O_{K_0}$ such that the $b_j$ are integral there, their pairwise
differences are units there, and the residue field has more than $r$
elements. Only finitely many primes are excluded. The Chinese remainder
theorem supplies $x,y\in K_0$ such that
\[
 v_{\mathfrak p_i}(x-b_i)=1,\qquad
 v_{\mathfrak p_i}(x-b_j)=0\ (j\ne i),\qquad
 v_{\mathfrak p_i}(y-b_j)=0\ (\text{all }j).
                                                        \tag{F47.R2}
\]
For the first condition specify $x$ modulo $\mathfrak p_i^2$; for the
last choose a residue different from each $b_j$. Thus $x,y\in V(L)$.
The elements $r_j=(x-b_j)/(y-b_j)$ are multiplicatively independent
modulo roots of unity: applying the valuations in (F47.R2) to a relation
$\prod_j r_j^{e_j}\in\mu_\infty$ gives $e_i=0$ for every $i$.
This independence persists in every algebraic overfield, since the
relation itself is an equality of the same algebraic numbers.

There is a number field $F_0$ for which $L\subset M=F_0\mathbb
Q^{\mathrm{ab}}$: take $F_0$ to contain $F$ and finitely many algebraic
generators of $L/k$. Theorem~1.1 of Amoroso--Zannier \citep{AmorosoZannier2000} implies
\[
 \bigcap_{n\ge1}(M^*)^n=\mu_\infty.
                                                        \tag{F47.R3}
\]
Indeed $[\mathbb Q^{\mathrm{ab}}(a):\mathbb Q^{\mathrm{ab}}]
\le [M:\mathbb Q^{\mathrm{ab}}]$ for $a\in M$, so their height bound
is uniformly positive on the non-torsion elements of $M$. If
$z_n^n=z\ne0$ for every $n$, then $h(z_n)=h(z)/n$, forcing $z$ to be
a root of unity. Conversely every root of unity is divisible in $M$,
which contains all roots of unity. Tsujimura's Proposition~4.9 \citep{Tsujimura2022}, applied with its parameter $f=1$ and number field
$F_0$, now says that
\[
 A=M^*/\mu_\infty
\]
is a free abelian group. This is the exact specialization of that
theorem used here.

The independent elements $[r_1],\ldots,[r_r]$ of $A$ have finite
support in a basis of the free abelian group. Their coordinate matrix
has rank $r$, so we may choose homomorphisms
$\chi_1,\ldots,\chi_r:A\to\mathbb Z$ such that the integer matrix
$N=(\chi_i([r_j]))_{i,j}$ has nonzero determinant. Choose actual
elements $u_x,u_y\in L^*$ realizing the two evaluations of $\xi$.
Subtracting the two evaluations of (F47.R1), then mapping to
$\widehat A$ and applying the completed characters, gives
\[
 N(\lambda_1,\ldots,\lambda_r)^{\mathrm t}
   =\bigl(\chi_i([u_x/u_y])\bigr)_i\in\mathbb Z^r.
                                                        \tag{F47.R4}
\]
The adjugate matrix shows that
$d\lambda_i\in\mathbb Z$, where $d=\det N\ne0$.
For each prime dividing $d$, membership in $\mathbb Z_p$ forces the
corresponding divisibility of that integer by the prime power in $d$.
Consequently $\lambda_i\in\mathbb Z$. Equivalently,
$\mathbb Q\cap\widehat{\mathbb Z}=\mathbb Z$ inside the finite
adeles. Put $m_i=\lambda_i$.

Finally choose any $a\in V(L)$ and an element $u_a\in L^*$ realizing
$a^*\xi$. The constant class in (F47.R1) is then the Kummer class of
$c=u_a/\prod_i(a-b_i)^{m_i}$. Thus
$\xi=[c\prod_i(t-b_i)^{m_i}]$, as required. When $r=0$ this last
evaluation argument proves the result directly. The actual unit is
unique modulo $\bigcap_n(L^*)^n$, a subgroup of the roots of unity;
the assertion does not identify different pointwise choices before
the translated-function argument is applied.


\subsubsection*{The smooth-curve step in cyclotomic reconstruction}

Fix an algebraic subfield $k\subset\overline{\mathbb Q}$ for which
$k\subset F\mathbb Q^{\mathrm{ab}}$ for some number field $F$.
Every finite extension of $k$ has the same property, after enlarging
$F$. Write $K=\overline{\mathbb Q}$. We use the rational-source
integrality lemma proved above, with all multiplicative completions
taken as $\varprojlim_n A/A^n$.

\emph{Arithmetic input on abelian varieties.}
For each finite extension $L/k$ and each abelian variety $A/L$,
\[
 \bigcap_{n\geq1} nA(L)=0.                                  \tag{F47.1}
\]
Indeed, the finite presentation of $A$ and a finite set of generators
of $L/k$ allow a number field $F_0$ such that
$L\subset F_0\mathbb Q^{\mathrm{ab}}$ and $A$ is obtained by base
change from an abelian variety over $F_0$. Tsujimura's
Theorem~3.1 \citep{Tsujimura2020} says exactly that the Kummer map
on $A(F_0\mathbb Q^{\mathrm{ab}})$ is injective. Its kernel is the
intersection of the multiplication-by-$n$ images, so the injection
$A(L)\hookrightarrow A(F_0\mathbb Q^{\mathrm{ab}})$ proves
(F47.1). The compatible total Tate module of $A(L)$ is also zero:
each coordinate of a compatible torsion system is divisible by every
integer, and (F47.1) kills it. Separately, the height argument in the
rational-source lemma gives
\[
 D(L^*):=\bigcap_n(L^*)^n\subset\mu_\infty.                  \tag{F47.2}
\]
Neither assertion requires that $L$ have a discrete valuation.

\emph{A smooth-curve Kummer specialization lemma.}
Let $C/L$ be a smooth geometrically connected affine curve and let
\[
 \xi\in\varprojlim_n H^1_{\mathrm{et}}(C,\mu_n).
\]
Assume that for every finite $L'/L$ and every $z\in C(L')$ the
specialization $z^*\xi$ is the Kummer class of an actual element of
$(L')^*$. After a finite extension of $L$ there is an actual regular
unit $g$ on $C$ such that $\xi=[g]$. This extension is only an
extension of constants used in the proof; no descent of the given
topos morphism is assumed.

To prove the lemma, let $\overline C$ be the smooth proper completion
and $S=\overline C\setminus C$. Enlarge $L$ so that the points of
$S$ and one point $z_0\in C$ are rational. Purity for a smooth curve
and the Kummer exact sequence give compatible residue maps
\[
 H^1(C,\mu_n)\longrightarrow
       \bigoplus_{s\in S}\mathbb Z/n,
 \qquad
 \operatorname{res}(\xi)=(\lambda_s)_{s\in S},
 \quad\lambda_s\in\widehat{\mathbb Z}.                    \tag{F47.3}
\]
We first prove that every $\lambda_s$ is an integer.

Fix $s\in S$. By Riemann--Roch and interpolation of values and a
first derivative at finitely many points, there is a nonconstant
rational function
\[
 \pi:\overline C\longrightarrow\mathbb P^1_L
\]
which has a simple zero at $s$ and takes nonzero values at every
other point of $S$. For example allow a sufficiently large pole
divisor at a point different from $S$, prescribe the first-order
jet of a uniformizer at $s$, and prescribe the value $1$ at the
other boundary points. The map is finite. Remove from $\mathbb P^1$
the values $0,\infty$, its branch values, and $\pi(S)$, and enlarge
$L$ once more to split this finite set if necessary. This gives
\[
 W=\pi^{-1}(V)\subset C,\qquad W\longrightarrow V
 \quad\hbox{finite etale},
\]
where $V$ is a punctured affine line. Apply the compatible norm
maps to obtain
\[
 \eta=\operatorname{Cor}_{W/V}(\xi|_W)
        \in\varprojlim_n H^1(V,\mu_n).
\]
The specialization property is preserved. If a fibre above
$a\in V(L')$ is the disjoint union of $\operatorname{Spec}L_i$,
etale base change for the norm and the usual Kummer norm identity
give the Kummer class of
\[
        \prod_i N_{L_i/L'}(u_i)\in(L')^*,
\]
where $u_i$ realizes the evaluation of $\xi$ at the corresponding
point of $W$. Thus the rational-source integrality lemma applies to
$\eta$.

The residue at $0$ of a corestriction is
$\sum_{z\mid0}[\kappa(z):L]\operatorname{res}_z(\xi)$.
This is the divisor-of-norm formula, and can also be read directly
from the norm of a local parameter in the Kummer boundary map.
Points of $C$ contribute zero because $\xi$ extends over them;
among the points of $S$ only the rational point $s$ lies over $0$.
Consequently this residue is exactly $\lambda_s$, with coefficient
one. The lemma on $V$ makes it integral. The choice of a simple
zero also makes explicit that no ramification factor has been
silently inserted.

Put $D=\sum_s\lambda_s[s]$, an integral divisor. The next arrow in
the localization sequence sends the residue vector to the Kummer
Chern class of $D$. Therefore
\[
 c_1(D)=0\quad\hbox{in }H^2(\overline C,\mu_n)
       \quad\hbox{for every }n,
\]
and the Kummer injection
$\operatorname{Pic}(\overline C)/n\hookrightarrow
H^2(\overline C,\mu_n)$ implies
$[D]\in\bigcap_n n\operatorname{Pic}(\overline C)$.
Its degree is zero. Any divisor class $E$ with $nE=[D]$ likewise
has degree zero. The rational point $z_0$ identifies degree-zero
line bundles with $J(L)$ for the Jacobian $J$ and removes the
Brauer obstruction to representing a Picard-scheme point by a line
bundle. Equation (F47.1) therefore gives $[D]=0$.
Choose $g_0\in L(C)^*$ with $\operatorname{div}(g_0)=D$;
it is a unit on $C$.

All residues of $\xi-[g_0]$ now vanish. For every $n$ the map
$H^1(\overline C,\mu_n)\to H^1(C,\mu_n)$ is injective: purity
identifies its preceding support group with a negative-degree
cohomology group of $S$. Thus its residue-free lifts are unique,
and hence compatible as $n$ varies. The Kummer sequences give
\[
 0\longrightarrow\widehat{L^*}\longrightarrow
  \varprojlim_n H^1(\overline C,\mu_n)
  \longrightarrow T\operatorname{Pic}(\overline C)
  \longrightarrow0.
\]
Here the left inverse system has surjective transition maps, so
there is no $\varprojlim^1$ obstruction. The total Tate module on
the right is zero by (F47.1). Thus $\xi-[g_0]$ is a constant
completed Kummer class. Evaluate at $z_0$: the hypothesis identifies
it with the Kummer class of an actual constant $c\in L^*$.
Then $g=cg_0$ proves the lemma.

\emph{Algebraicity of the point function for a smooth curve.}
Let
$\phi:C_{\mathrm{et}}\to(\mathbb G_m)_{\mathrm{et}}$
be pinned and relative to $B G_L$, and write
$f:C(K)\to K^*$ for its point function. It is Galois equivariant
and sends $L'$-points to $L'$-points for every finite extension,
as checked in F.46. Its pulled-back Kummer class satisfies the
specialization lemma. After a finite extension choose
$g_0\in\mathcal O(C)^*$ representing it. Equation (F47.2) gives
\[
 f(x)/g_0(x)\in\mu_\infty\qquad(x\in C(K)).                 \tag{F47.4}
\]
If the generic image is a closed point, the inverse-image of its
open complement is empty by continuity and the Jacobson property.
The morphism factors through that closed subtopos. After splitting
its residue field it has a constant point function. We may therefore
assume the generic image is generic.

For each $b\in K$, base change to contain $b$ and pull back the
open $\mathbb G_m\setminus\{b\}$. It is a dense open
$C_b\subset C$ whose omitted geometric points are finite, and
$x\in C_b(K)$ exactly when $f(x)\ne b$. Applying the curve lemma
to the Kummer classes of $t-b$ gives a rational function $g_b$, a
unit on $C_b$, with
\[
 f(x)-b=d_b(x)g_b(x),\qquad d_b(x)\in\mu_\infty,
       \quad x\in C_b(K).                                 \tag{F47.5}
\]
Use the fixed proper field $M=F\mathbb Q^{\mathrm{ab}}\subset K$.
In any one argument, a finite extension $N/M$ contains the finitely
many coefficients needed there; the field $M$ itself stays fixed.
All the coefficients
$d_b(x)$ lie in $M^*$. The extension $K/M$ has infinite degree,
as does $K/N$ for every finite extension $N/M$ in $K$.
For clarity, the latter assertion can also be seen from the height
bound: if $K=N$ were finite over $M$, that bound would be uniformly
positive on its nontorsion elements, whereas $K$ contains all roots
of the rational number $2$.

The finite-exception quasi-rationality argument in F.46 applies to
an arbitrary smooth curve, with the following explicit replacement
for its polynomial-degree step. Descend the curve and the finitely
many rational functions under consideration to a finite extension
$N/M$. The poles of any fixed linear combination of two such
functions lie in one fixed divisor. If the combination is nonzero,
each zero has residue degree over $N$ bounded by the degree of that
divisor. Points of arbitrarily large residue degree exist: choose a
finite map of the proper curve to $\mathbb P^1_N$, choose values of
arbitrarily large degree over $N$, and take points above them.
They avoid any prescribed finite exceptional set. Thus a vanishing
at such a point forces the same rational-function identity as in
F.46. The two shifted outside-$M$ sets cover the cofinite domain,
and have infinite intersection by this same degree bound.

If $g_0$ is nonconstant, this proves that $f$ equals a rational
function $R\in K(C)$ on a cofinite subset. If $g_0$ is constant,
divide $f$ by one of its values and use (F47.4) to put its values
in $\mu_\infty\subset M$. Choose $b\notin M$. If $g_b$ were
nonconstant, (F47.5) would put its cofinite image in the vector
space $M+Mb$, which cannot contain a cofinite subset of $K$.
Hence $g_b$ is constant. Two different values of $f$ would, by
subtracting their equations, force this constant into $M$, and
then force $b\in M$. Thus $f$ is constant in this case.

We remove the finite exceptional set in the nonconstant case.
For every $b$, the rational function $(R-b)/g_b$ takes its values
in $\mu_\infty$ on a cofinite set. A nonconstant rational function
on a smooth curve has cofinite image in $K$, so this quotient is a
constant $c_b\in\mu_\infty$. Fix $x\in C(K)$ and choose
$b\ne f(x)$. Since $g_b$ is a unit at $x$, the identity
$R-b=c_b g_b$ first excludes a pole of $R$ at $x$. For every
$b\notin\{f(x),R(x)\}$, it then gives
\[
 \frac{f(x)-b}{R(x)-b}\in\mu_\infty.
\]
If $f(x)\ne R(x)$ this fractional linear expression takes all
values of $K\setminus\{0,1\}$ as $b$ varies, a contradiction.
Therefore $f=R$ on every geometric point. It follows that $R$ is
a regular unit on $C_K$. Galois equivariance and equality on the
Zariski-dense geometric points imply that $R$ descends to a regular
unit on the original curve over $L$. This proves the full
point-function assertion for every smooth affine curve and every
finite extension of the base field.

It also gives the equality of compatible Kummer classes that will
be used below. Over the finite extension $L'/L$ employed in the
specialization lemma, the rational function $R/g_0$ has values in
$\mu_\infty$ at every geometric point and is therefore a constant
$\zeta\in(L')^*$. Evaluation at a rational point of the curve gives
$\zeta\in D((L')^*)$, so $[R]=[g_0]=\xi$ over $L'$.
If $C(L)$ is nonempty, equality already holds over $L$. Indeed,
for every $n$ the residual class in $H^1(C,\mu_n)$ is geometrically
zero, and hence is a constant class by the low-degree
Hochschild--Serre sequence. Its evaluation at an $L$-point is zero,
because $R$ is the actual point function there. The same sequence,
or the rational-point section, shows that this constant class is
zero. Without an $L$-point, equality after a finite extension is all
that is asserted or needed for the geometric residue tests.

\subsubsection*{From the curve calculation to arbitrary dimension}

We give the higher-dimensional step independently, so that the passage
from point values to a geometric morphism does not rest on a
fundamental-group calculation. Let $k\subset\overline{\mathbb Q}$ be
contained in $F\mathbb Q^{\mathrm{ab}}$ for some number field $F$.
We use the preceding curve calculation in precisely the following
form: after every finite extension $L/k$, a pinned relative geometric
morphism $C_{\mathrm{et}}\to(\mathbb G_{m,L})_{\mathrm{et}}$, with
$C$ a smooth geometrically connected affine curve having an $L$-point,
has the geometric-point values of a unit of $C$. Its compatible Kummer
class is the class of that unit. The conclusion below is conditional
only on that already proved curve calculation and the arithmetic
inputs explicitly recorded here.

It is enough that the curve calculation initially represents its
Kummer class after a further finite extension. Indeed, suppose the
pointwise unit is $f_C\in\mathcal O(C)^*$ and, over a finite
extension $L'/L$, the Kummer class is $[g_0]$. At every algebraic
point the value of $f_C/g_0$ is divisible in the multiplicative group
of a finite extension of $L'$ and hence is a root of unity, by the
height input used in the curve calculation. A nonconstant rational
function on a geometrically integral curve has cofinite image in
$\mathbb G_m(\overline L)$, so this ratio is constant. Evaluation at
an $L$-point shows that its constant value is itself divisible in
$L'^*$; hence its Kummer class is zero. Thus the two classes agree
after $L'$. At each finite Kummer level, their difference on $C$
lies in the image of $H^1(L,\mu_n)$ by the Hochschild--Serre exact
sequence, since it is zero geometrically. Evaluation at the $L$-point
is a retraction onto this image and is zero on the difference.
Consequently the classes already agree over $L$. This proves exactly
the curve input used below, without requiring a chosen function to
descend merely because its Kummer class descends.

\emph{The arithmetic Picard lemma.}
For every finite extension $L/k$ and every abelian variety $A/L$,
\[
                 \bigcap_{n\geq1}nA(L)=0.                 \tag{F47.H1}
\]
Indeed, choose a number field $F_0$ over which $A$ is defined and large
enough that $L\subset M=F_0\mathbb Q^{\mathrm{ab}}$. Theorem~3.1 of
Tsujimura (2020) states that the Kummer map on $A(M)$ is injective
\citep{Tsujimura2020}. Its proof explicitly factors this map as
\[
 A(M)\longrightarrow\varprojlim_n A(M)/nA(M)
       \lhook\joinrel\longrightarrow H^1(G_M,T A),
\]
so its kernel is exactly $\bigcap_n nA(M)$. Thus (F47.H1) follows by
inclusion. Equivalently, the finite nonempty division fibres of a point
in $\bigcap_n nA(M)$ have a compatible inverse-limit point: use the
factorial levels and the finite intersection property. There is no
distinction here between unrelated divisions and coherent divisions.

If $Y/L$ is smooth, projective, and geometrically integral, put
$P=\operatorname{Pic}(Y)$ and let $P^0$ be the subgroup whose geometric
classes are algebraically trivial. Then
\[
                 \bigcap_{n\geq1}nP=0,
                 \qquad T P=\varprojlim_n P[n]=0.         \tag{F47.H2}
\]
Here $P^0$ injects into the group of $L$-points of the abelian variety
$\operatorname{Pic}^0_{Y/L}$, and
$N=P/P^0$ is a subgroup of the finitely generated geometric
Neron--Severi group; for the latter finiteness statement see
\cite[Theorem~7.3 and Lemma~7.2]{FA_ChengEtAl2021}.
These assertions do not require the Picard scheme
to represent every rational Picard class: the possible Brauer
obstruction only changes the image of $P^0$ inside
$\operatorname{Pic}^0_{Y/L}(L)$. Injectivity of the passage from actual
line bundles follows from Hilbert~90.

For completeness, let $x\in\bigcap_n nP$. Its image in $N$ vanishes,
since $N$ is finitely generated. Let $e$ annihilate the torsion subgroup
of $N$. For every $m\geq1$, choose $y\in P$ with $em y=x$.
Then the image of $y$ in $N$ is torsion, so $z=ey$ lies in $P^0$ and
$mz=x$. Thus $x$ is infinitely divisible in
$\operatorname{Pic}^0_{Y/L}(L)$, and (F47.H1) gives $x=0$. Finally, each
coordinate of a compatible torsion system in $T P$ is divisible by
every positive integer, so the first equality in (F47.H2) implies the
second. This proves the lemma, including the torsion in $N$.

\emph{Curves with prescribed points and transverse directions.}
We will use the following elementary consequence of Bertini. In a
smooth geometrically integral projective variety over an infinite
characteristic-zero field, a finite set of rational points lies on a
smooth geometrically integral complete-intersection curve. At any
specified point one may require its tangent line to be transverse to
a given smooth divisor. The construction is over the ground field;
one first makes the chosen points rational by a finite extension when
necessary. To check the imposed conditions, take sufficiently high
powers of an ample bundle. Serre vanishing makes evaluation on the
finite first-order neighbourhoods surjective. The linear systems
vanishing at the prescribed points embed the complement of those
points, and a general member is smooth there and geometrically
irreducible. Smoothness and the required independent differentials at
the finitely many prescribed points are nonempty open conditions on
the first jets. Intersect these opens and repeat until dimension one.
This is the usual prescribed-point version of Bertini; the needed
linear-system and irreducibility statements are Stacks, Lemmas
33.47.1, 33.47.3, and 37.32.3 \citep{FA_StacksBertini}.

\emph{Integral boundary residues and the global unit.}
Let $X/k$ be a connected regular affine scheme of finite type with
$x_0\in X(k)$, and let
\[
       \phi:X_{\mathrm{et}}\longrightarrow
                        (\mathbb G_{m,k})_{\mathrm{et}}
\]
be pinned over $B_k$. Write $\phi(a)$ for its value at a geometric
point. Over a finite extension $L/k$, a point $a\in X(L)$ has
$\phi(a)\in L^*$: the relative morphism of closed-point Galois topoi
forces the residue field of the image to embed in $L$. Base change of
$\phi$ to $L$ is obtained by slicing by the finite etale constant-field
object. These observations justify every evaluation used below.

The scheme $X$ is smooth and geometrically integral. In fact, in
characteristic zero regularity implies smoothness, and the algebraic
closure of $k$ in $k(X)$ is a finite field extension contained in all
normal local rings of $X$; evaluation at $x_0$ embeds it in $k$.
Choose a smooth projective compactification $Y$ of $X$ whose boundary
is a divisor with normal crossings, by characteristic-zero embedded
resolution. Concretely, take a projective closure, resolve it while
preserving the smooth open $X$, and principalize its boundary ideal;
the blowups are projective and supported outside $X$. The strong
resolution and monomialization theorem in
\cite[``The result'' and ``Idea of proof'', pp.~1--3]{FA_EncinasHauser2002}
supplies these operations. Write the boundary prime components over
$k$ as $D_1,\ldots,D_s$.
Only smooth points of individual boundary components will be used.

Pulling back the Kummer torsors of the coordinate $t$ gives
\[
       h=(h_n)_n\in\varprojlim_n H^1_{\mathrm{et}}(X,\mu_n).
                                                               \tag{F47.H3}
\]
Let $K=k(Y)$. The map
$H^1(X,\mu_n)\to K^*/K^{*n}$ is injective: a generically trivial
finite etale torsor on a normal integral scheme has a section, by
normalization. For every prime divisor $D_i$, the valuation of the
generic Kummer representative of $h_n$ defines a residue
$r_{i,n}\in\mathbb Z/n$. Compatibility gives
$r_i\in\widehat{\mathbb Z}$.

We claim $r_i\in\mathbb Z$. After a finite extension $L/k$, choose
an $L$-rational point $q$ on a geometric component above $D_i$, away
from the other boundary components, and a smooth complete-intersection
curve $\overline C\subset Y_L$ passing through $q$ transversely to
that component and through $x_0$. Its open part
$C=\overline C\cap X_L$ is an affine smooth geometrically integral
curve with an $L$-point. The curve calculation makes $h|_C$ the Kummer
class of an actual unit; its residue at $q$ is therefore an integer.
That residue is exactly $r_i$.

One can check this equality without an inverse-limit purity argument.
At $q$, the regular local ring $R$ is factorial and the only boundary
has one prime equation $z$. Hence
$\operatorname{Pic}(R[1/z])=0$ and
$R[1/z]^*=R^*z^{\mathbb Z}$. For each $n$, the restricted Kummer
class has a representative $u z^j$, with $u\in R^*$, and residue
$j\bmod n$. On the transverse curve $z$ is a local parameter and
$u$ remains a unit, giving the same residue for every $n$. The base
change from $D_i$ to a geometric component has multiplicity one in
characteristic zero. This proves the claim.

Set $D=\sum_i r_iD_i$, now an integral Cartier divisor on $Y$.
Its class lies in $\bigcap_n n\operatorname{Pic}(Y)$. Indeed, choose
a rational function $a_n\in K^*$ representing the generic class of
$h_n$. Its valuations at divisors inside $X$ are multiples of $n$, by
the Kummer description over a discrete valuation ring; at $D_i$ its
valuation is $r_i\bmod n$. Consequently
\[
                    \operatorname{div}(a_n)=D+nE_n
                                                               \tag{F47.H4}
\]
for an integral divisor $E_n$ on the regular scheme $Y$.
Thus $[D]=-n[E_n]$. By (F47.H2), $[D]=0$, and there is a rational
function $g\in k(Y)^*$ with $\operatorname{div}(g)=D$. In
particular, $g|_X$ is an actual regular unit.

Subtracting $[g]$ from $h$ leaves a compatible class unramified on
$Y$. Explicitly, the divisor of $a_n/g$ is divisible by $n$, so its
associated line bundle with an $n$th-power trivialization defines a
$\mu_n$-torsor on $Y$ extending $h_n-[g]$. Its uniqueness follows
from injectivity into $K^*/K^{*n}$, so these extensions are compatible
as $n$ varies. The Kummer exact sequences on $Y$ and (F47.H2) give
\[
  \varprojlim_n H^1(Y,\mu_n)=\widehat{k^*}.
                                                               \tag{F47.H5}
\]
To see that no derived-limit term is being suppressed, the maps
$k^*/k^{*mn}\to k^*/k^{*n}$ are surjective; alternatively, since the
compatible Picard-torsion class is zero, each class is represented
uniquely in $k^*/k^{*n}$, and these representatives are automatically
compatible. Here $\Gamma(Y,\mathcal O_Y)^*=k^*$ by geometric
integrality and properness.

Evaluation at $x_0$ identifies the constant class in (F47.H5) with the
actual element $c=\phi(x_0)/g(x_0)\in k^*$. Therefore
\[
        f=cg|_X\in\Gamma(X,\mathcal O_X)^*,\qquad
        h=[f],\qquad f(x_0)=\phi(x_0).                 \tag{F47.H6}
\]

\emph{Pointwise equality in every dimension.}
Let $a\in X(\overline k)$. Choose a finite extension $L/k$ making
$a$ rational and a smooth geometrically integral curve
$\overline C\subset Y_L$ through both $x_0$ and $a$; put
$C=\overline C\cap X_L$. In dimension one take $C=X_L$ itself.
The curve theorem gives a unit $f_C$ whose values equal those of
$\phi|_C$. Since both $f_C$ and $f|_C$ have Kummer class $h|_C$,
the unit $u=f_C/(f|_C)$ has zero Kummer class for every $n$. The
injection
\[
  \Gamma(C,\mathcal O_C)^*/\Gamma(C,\mathcal O_C)^{*n}
       \lhook\joinrel\longrightarrow H^1(C,\mu_n)
\]
shows that $u$ has an $n$th root among regular units for every $n$.
Its divisor on $\overline C$ is therefore divisible by every positive
integer and is zero. Thus $u\in L^*$ is constant. Evaluation at
$x_0$ gives $u=1$, and consequently
\[
                         \phi(a)=f(a)
                  \quad\hbox{for all }a\in X(\overline k).
                                                               \tag{F47.H7}
\]
In dimension zero $X=\operatorname{Spec}k$ and the assertion is
immediate. Notice that this final argument does not infer equality
from a finite collection of Kummer quotients: it uses the integral
curve realization and equality in every Kummer degree. It also does
not require the residual roots of unity to form a chosen compatible
system.

\emph{Recovery of the entire geometric morphism.}
The assertion just proved holds for every connected regular affine
finite-type source with a rational point, and after every finite
extension of the ground field. It is exactly condition (4.18.5) of
Carlson--Haine--Wolf, Theorem~4.18 \citep{CHW2024}. Their equivalence
with (4.18.1) therefore gives the full pinned reconstruction statement
for this field class, with absolute weak normalization as in their
definition. This application addresses all etale objects, not just
the point map of one morphism.

Here is the categorical part of that implication. Two copies of
$\mathbb G_m$ cover $\mathbb A^1$, so the pointwise units just
constructed give pointwise algebraic maps to $\mathbb A^1$. Affine
coordinates and equality on geometric points of a reduced source
then reconstruct the addition and multiplication relations and hence
a map to every affine target. Apply this construction also to each
sliced pinned morphism over a quasicompact separated etale target
object $V$. Its inverse image is again represented by a
quasicompact separated etale source scheme, by CHW Corollary~A.9;
one works on its affine components and, when needed, after a finite
extension supplying a rational point. Uniqueness on geometric points
glues these maps and makes them compatible with all morphisms of
etale objects. Proposition~4.14 then identifies the original inverse
image functor with the functor of the reconstructed scheme map.
Finally, the $h$-descent part of Theorem~4.18 handles arbitrary
finite-type sources. Thus no openness assertion for an unconstructed
birational Galois homomorphism, and no inference from finite set
coefficients alone, is needed.

\end{proof}

\Needspace{5\baselineskip}
\paragraph{F.48 (Proved). Recovered Section 53 comparison and the exact cyclotomic rows.}
\textit{Manuscript: Longue Marche, \S53.5A--C, folios 768--779,
cote 140-4 scans 242--262, especially (47), (56), (81)--(94);
Montpellier cote 149, \S1, folios 6--7, scans 13--15, and \S4,
folios 63--65, scans 128--131; cote 147, \S11, scans 143--148.}
The one-dimensional specialization of F.19 on a smooth surface is the sector-filter
functor actually constructed in \S53, not just an abstract equivalence
between two groupoids having fundamental group $\mathbb Z$. Its finite
completion gives the complex comparison of the tangential fibre functor.
The recovered \texttt{CYCLOTOMIC\_KERNEL} register has exact row identifiers
\texttt{B5} and \texttt{C8}. B5 is a slimness statement; C8 is diagonal
boundary nonvanishing together with a counterexample to unrestricted
local injectivity. Their mathematical contents agree with the corrected
statements below, and do not assert a universal converse recognizing
cuspidal sections.

\begin{proof}
\emph{The germ and its coverings.}
Let $x$ be a point of a smooth real surface $X$ and write
$D_x^*$ for the germ of punctured neighbourhoods of $x$.
Let $S_x=(T_xX\setminus0)/\mathbb R_{>0}$ be its circle of positive
half-rays. A covering of a germ means the filtered colimit of the
categories of coverings of its representatives, exactly (47) on scan
242. Small punctured discs form a cofinal system; restriction between
them induces an equivalence on all covering spaces. Therefore
\[
 \operatorname{Rev}(D_x^*)\simeq\mathbb Z\text{-}\mathbf{Set},
 \qquad \operatorname{Rev}_{\mathrm{uni}}(D_x^*)\simeq B\mathbb Z.
 \tag{F48.1}
\]
Here the latter is the groupoid of connected simply connected covering
germs and their isomorphisms. Choosing an orientation identifies its
automorphism group with $\mathbb Z$; intrinsically it is the orientation
module denoted $\overline T$ in (56), scan 248.

The poset $\mathcal A_x$ in (82), scan 256, consists of nonempty
\emph{closed proper connected arcs} of $S_x$, including single points
and excluding the whole circle. Its groupoid completion is
$\Pi_1(S_x)$. Indeed, choose a point in every arc. An inclusion of
arcs determines a unique homotopy class of paths in the larger arc.
Conversely, subdividing any path into finitely many arc neighbourhoods
expresses it by these generators. A homotopy of paths has a finite
subdivision into triangles lying in such neighbourhoods, so the only
relations required are compositions of inclusions and their inverses.
These constructions are inverse up to the evident changes of the
chosen points. This proves (83)--(85), scans 257--258, without assuming
that the arc poset is simply connected.

\emph{Strict thickening and coordinate independence.}
For $A\in\mathcal A_x$, let $C_A$ be its cone in $T_xX$. Choose a
local differentiable chart from the tangent plane to $X$ inducing the
identity on $T_xX$. The germ associated to $A$ is represented by the
filter generated by
\[
  \{v:0<\lVert v\rVert<\epsilon,\ [v]\in A'\},
  \qquad A\subset\operatorname{int}(A'),\quad A'\subsetneq S_x,
 \tag{F48.2}
\]
where $A'$ runs through strict angular enlargements and $\epsilon>0$.
One may take $A'$ open or take interiors of closed strict enlargements;
the generated filters are cofinal. Every such representative is
contractible, by polar coordinates on an angular interval. No
convexity assertion is needed when the interval is wider than a
semicircle.

If $h$ changes the chosen chart and $d h_0=1$, then
$h(v)=v+o(\lVert v\rVert)$ uniformly in direction. Given a strict
enlargement $A'$ choose an intermediate enlargement $A''$ whose closure
is contained in $\operatorname{int}(A')$. For sufficiently small $\delta$,
$h$ sends the sector for $(A'',\delta)$ into the sector for
$(A',\epsilon)$. Apply the same observation to $h^{-1}$.
Thus both charts define the same filter germ. This is precisely the
strict-enlargement construction (89)--(90) on scans 259--260.
Using just a fixed ray would fail: in coordinates $(s,t)$ the
diffeomorphism $(s,t)\mapsto(s,t+s^2)$ has identity derivative at zero
but does not preserve the germ of the positive $s$-axis.

Pulling a covering back to this simply connected sector germ gives a
set of sheets and hence a fibre functor $F_A$ on
$\operatorname{Rev}(D_x^*)$. Inclusions of arcs give canonical
isomorphisms between these fibre functors, by restriction to the
smaller sector germ; their compatibility is the compatibility of
restriction. The associated universal covering germs therefore define
\[
 \Theta:\Pi_1(S_x)\longrightarrow
              \operatorname{Rev}_{\mathrm{uni}}(D_x^*),
 \tag{F48.3}
\]
the functor (81), scan 256. Rotating the direction once lifts, in
angular coordinates, to a translation by one full turn on the
universal covering of a punctured disc. Thus $\Theta$ sends the
generator of $\pi_1(S_x)$ to a generator of the automorphism group
in (F48.1), with the sign fixed by the chosen fibre-functor and
representing-cover convention. Both groupoids are connected and nonempty, so $\Theta$ is
fully faithful and essentially surjective. This proves the
equivalence by the winding criterion in (86bis), scan 258.

\emph{Arcs and tangential points.}
An outgoing differentiable arc with nonzero derivative and tangent
half-ray $\Delta$ eventually lies in every strict sector enlargement
of $\Delta$. It consequently gives a map of germs into (F48.2).
On coverings this identifies its fibre functor with $F_{\{\Delta\}}$,
canonically and compatibly with morphisms of universal covers. These
are (91)--(94), scans 261--262. The normal-direction construction in
F.19 uses these same simply connected sectors and restriction maps;
its rank-one specialization is consequently (F48.3). This verifies
agreement of the actual functors, including their path morphisms.

For a complex algebraic curve, apply the finite covering comparison
and complete the automorphism group to $\widehat{\mathbb Z}(1)$.
In a local parameter $z$, the finite connected covers are cofinally
the Kummer covers $z^{1/n}$; continuation along the sector picks their
compatible sheets. This is the fibre functor in Deligne's
\S15.13--15.27, as identified in F.19. Replacing $z$ by $u z$ changes
the compatible lifts by choices of $n$th roots of the unit $u$;
their ratios are exactly the Kummer transition torsor. Thus the
comparison retains the parameter-change identifications. In
particular it does not invent a canonical map of spaces
$D_x^*\to S_x$, whose absence is explicitly discussed on scan 262.

There is a separate correction on scans 242--243: the speculative
identification with the inverse limit of the \emph{full small topoi}
of all shrinking punctured neighbourhoods is false. Their localic
intersection is empty, so that inverse limit is the degenerate topos,
whereas (F48.1) is not degenerate. The correct objects are the
pro-topos of neighbourhoods, or the category of covering germs.
This is the exact LM8 E.10/E.11 distinction proved in F.12, now
checked against the original wording, including the inserted word
\emph{petits} on scan 243.

\emph{The recovered register B5.}
Write $\mathsf{Slim}(G)$ for the condition
$C_G(H)=1$ for every open $H\subseteq G$. It is preserved by
extensions: the image of a centralizer element is trivial in a slim
quotient, and its remaining element in the kernel centralizes the
open intersection there. It is also preserved by inverse limits of
slim profinite groups with surjective transition maps. A centralizer
element has trivial coordinates, since projections are open and
send the given open subgroup to open subgroups.
The nonabelian finite-rank free profinite case is the case established
in the free-group argument of F.39. More explicitly, an open normal
subgroup has free rank $1+[F:H](\operatorname{rank}F-1)>1$ and trivial
centre. Its centralizer injects into the finite group $F/H$ and is
trivial by torsion-freeness. Rank one is excluded because
$F=\widehat{\mathbb Z}$ is its own centralizer. This proves exactly
the three assertions of the recovered B5, including its correction
to scan 15. The identifier does not designate an inertia-recognition
or cyclotomic-normalization theorem.

Here torsion-freeness can be checked without a separate group-theoretic
black box. If a nontrivial finite-order element $z\in F$ survives in a
finite quotient $F/N$, put $V=N\overline{\langle z\rangle}$. The open
group $V$ is free profinite, so its abelianization is a finite product
of copies of $\widehat{\mathbb Z}$, a torsion-free group. But the
nonzero cyclic quotient $V/N$ generated by $zN$ factors through that
abelianization, while the finite-order element $z$ maps to zero there.
This contradiction proves the assertion used above.

\emph{The recovered register C8.}
Let $U/k$ be a tower of relative hyperbolic curves in characteristic
zero, $U\subset X$ a smooth proper normal-crossings compactification,
and $x$ a rational boundary point with geometric branch set $I$.
The local geometric group is $T^I$, $T=\widehat{\mathbb Z}(1)$,
with action given by cyclotomic scaling and branch permutation
\citep[Proposition~3.4]{CarlsonCorwin2026}. For its map
$j_x:T^I\to\pi_1(U_{\bar k})$, the diagonal $T\to T^I$ has nonzero
image. Here is the additional argument, beyond local normalization.
Hoshi's valuative criterion for a relative hyperbolic curve says that
a map from the fraction field of a strictly henselian trait extends
over the trait precisely when its fundamental-group image is trivial
\citep[Lemma~2.8, pp.~175--176]{Hoshi2014}. Induction up the tower
gives that criterion for $U$: extend the map to the previous stage
first, and then apply the relative-curve criterion at the last stage.

In formal branch coordinates at a geometric point over $x$, take
the arc $z_i=t$ for every boundary parameter, and take the remaining
parameters to be zero. It has generic point in $U$ and closed point
$x$. Pullback of $z_i^{1/n}$ to $t^{1/n}$ shows that its fundamental
group map is exactly the diagonal followed by $j_x$. If this were
zero the criterion would extend the generic-point map to $U$.
Separatedness of $X$ would identify this extension with the original
arc, contradicting its boundary closed point. This proves
nonvanishing. The cyclotomic kernel fixes the diagonal regardless
of branch permutation, and fixes all of $T^I$ when the branches
are rational. The stated cyclotomic centralizer conclusion follows.

Full injectivity fails in the blown-up product already computed in
F.39: at the new crossing the map is
$T^2\to T$, $(a,b)\mapsto a+b$, with kernel generated by $(1,-1)$.
The diagonal maps by multiplication by two and is nonzero. Thus the
stable-model hypotheses in F.39 and the counterexample in C8 are
compatible. A converse boundary-local factorization of all sections
with nonzero cyclotomic centralizer would imply the point-or-cusp
section assertion for curves: apply it to the product of an arbitrary
curve section and a rational cusp section of the triply punctured
line, and project back to the curve. Neither B5 nor C8 supplies this
converse. Carlson and Corwin's Theorem~3.7 and Proposition~4.12
concern closures of the corresponding section sets; dropping those
closure bars would change their statements.
\end{proof}

\emph{Audit evidence.}
The recovered register, the B5 and C8 proof files, the two \S53
transcriptions, and their file hashes are retained in the source-audit
appendix of the package. The decisive original scans 242, 243, 256,
259--262 were inspected visually. Earlier unsuccessful-search notes
are preserved as dated history and superseded here. The separate
\texttt{FLAGGED\_149} register remains the source obligation in F.45.

\Needspace{5\baselineskip}
\paragraph{F.49 (Proved). Further field ranges where pinned reconstruction fails.}
\textit{Manuscript: MPL149; section2; folio25v; scan52; prompt Part F field-range classification.}

The functor from finite-type $k$-schemes with universal homeomorphisms
inverted to pinned etale topoi over $B G_k$ is not full in each of
the following cases:
\begin{enumerate}
\item $k$ is any algebraic extension of a finite field, including
      $\overline{\mathbb F}_p$;
\item $k$ is infinite and perfect and $Z(G_k)\ne1$; in particular,
      $k$ is any real closed field;
\item $k$ is any algebraically closed field.
\end{enumerate}
In every case an endomorphism of $(\mathbb A^1_k)_{\mathrm{et}}$
already gives a counterexample. The relative comparison with the
base topos is part of the morphism. Forgetting that comparison would
erase the counterexamples constructed in the second case.

\emph{Manuscript problem locus.} Montpellier cote149, \S2,
fol.25v, scan52, including the margin asking about characteristic
zero; these statements sharpen the field-range delimitation in F.44.

\emph{Proof: the central comparison twist and its point function.}
Let $k$ be perfect, choose a separable closure $K$, put $H=G_k$,
$X=\mathbb A^1_k$, and let $p:X_{\mathrm{et}}\to BH$ be the
structure morphism. A central element $z\in H$ acts naturally and
$H$-equivariantly on every discrete continuous $H$-set. Pulling back
this natural automorphism to $X_{\mathrm{et}}$ gives an invertible
two-cell of $p$. Define
\[
 T_z=(\operatorname{id}_{X_{\mathrm{et}}},\alpha_z)
       \quad\text{in the category of topoi over }BH,
 \qquad \alpha_z:p\circ\operatorname{id}\Longrightarrow p.
                                                               \tag{F49.1}
\]
Choose the orientation of this comparison so that on the fibre of a
constant finite etale scheme $E/k$ it is
\[
 \operatorname{Hom}_k(E,K)\longrightarrow
 \operatorname{Hom}_k(E,K),\qquad u\longmapsto z\circ u.
                                                               \tag{F49.2}
\]
The underlying geometric morphism is the identity, hence is pinned.

Its induced function on framed geometric points is
\[
                         a\longmapsto z(a),\qquad a\in K.
                                                               \tag{F49.3}
\]
Here is an explicit check of the relative data. A standard geometric
point at $a$ has the standard identification of its restriction to
constant etale schemes with the fixed fibre functor of $BH$.
Composition with (F49.1) changes that identification to (F49.2).
Applying $z$ to every lift of the point through an etale neighbourhood
sends a lift above $a$ to a lift above $z(a)$. On constant covers
this is exactly (F49.2), so it identifies the twisted framed point
with the standard framed point at $z(a)$.

Distinct elements of $K$ give nonisomorphic relative closed geometric
points. For completeness, over a closed point with residue field
$E\subset K$ their groupoid is the homotopy fibre of
$B G_E\to B G_k$ at the fixed fibre functor. Since $G_E\to G_k$
is injective, this groupoid has trivial automorphism groups, and its
isomorphism classes are the embeddings $E\hookrightarrow K$.
Different closed supports have no isomorphisms between them.
This is also the case $L=K$ of Carlson--Haine--Wolf,
Proposition~3.6 \citep{CHW2024}. Consequently (F49.3) distinguishes
the relative morphisms $T_z$ without any assertion about the centre
of the generic Galois group of $X$.

\emph{Proof: infinite perfect fields with nontrivial Galois centre.}
Take $1\ne z\in Z(G_k)$. Formula (F49.3) fixes every $a\in k$
and moves some $a\in K$. No endomorphism from the scheme
localization can have this point function. We make this assertion
explicit, including the inverted universal homeomorphisms.
By the absolute-weak-normalization description of that localization
(F.4; Carlson--Haine--Wolf, Proposition~2.28), its endomorphisms
from $\mathbb A^1_k$ to itself are represented by
\[
 \operatorname{Hom}_k((\mathbb A^1_k)^{\mathrm{awn}},\mathbb A^1_k).
                                                               \tag{F49.4}
\]
In characteristic zero $\mathbb A^1_k$ is already absolutely weakly
normal, so a member is given by $P(t)\in k[t]$. If $P(a)=a$ for
every $a\in k$, the infinitude of $k$ gives $P=t$.
In characteristic $p$, (F49.4) is the set
\[
                  k[t^{1/p^\infty}]
                    =\bigcup_{n\geq0}k[t^{1/p^n}].
\]
For an element $h$ of this ring choose $n$ so that
$h^{p^n}=P(t)\in k[t]$. If $h(a)=a$ for every $a\in k$, then
$P(a)=a^{p^n}$ on an infinite set; hence $P=t^{p^n}$.
The ring is reduced, so Frobenius is injective and $h=t$.
In either characteristic the only possible scheme morphism is the
identity. Its geometric-point function fixes all of $K$, whereas
(F49.3) does not. Thus $T_z$ is an explicit morphism outside the
image. A real closed field has algebraic closure $k(i)$ and
$G_k=C_2$; complex conjugation supplies the stated example.

\emph{Proof: all algebraic extensions of finite fields.}
Let $k\subset\overline{\mathbb F}_p$ be algebraic over
$\mathbb F_p$ and write
\[
 H=\operatorname{Gal}(\overline{\mathbb F}_p/k)
       \subset\operatorname{Gal}(\overline{\mathbb F}_p/\mathbb F_p)
       =\widehat{\mathbb Z}.
\]
If $H\ne1$, it is abelian and has cardinality $2^{\aleph_0}$.
Indeed a nonzero element has a nonzero projection to some
$\mathbb Z_\ell$; the closure of its integer multiples projects
onto a nonzero subgroup of $\mathbb Z_\ell$, already of that
cardinality. Every $z\in H$ gives (F49.1), and (F49.3) makes
their isomorphism classes distinct. On the scheme side there are
only countably many localized endomorphisms: $k$ is countable and
the ring $k[t^{1/p^\infty}]$ in (F49.4) is countable. Thus the
functor cannot be full. This handles finite fields as well as all
proper infinite algebraic subfields. For the infinite subfields the
preceding argument also gives the stronger explicit obstruction for
every nonidentity $z\in H$.

If $H=1$, then $k=\overline{\mathbb F}_p$ and $BH=\mathrm{Set}$.
For each $\sigma\in\widehat{\mathbb Z}$ the ring automorphism
\[
 k[t]\longrightarrow k[t],\qquad
    \sum_i a_i t^i\longmapsto\sum_i\sigma^{-1}(a_i)t^i
                                                               \tag{F49.5}
\]
gives a semilinear automorphism of the underlying scheme.
Its map on closed points is $a\mapsto\sigma(a)$.
The resulting topos automorphism is automatically relative to
$BH=\mathrm{Set}$, and it is pinned. Distinct $\sigma$ induce
distinct point maps, so they give $2^{\aleph_0}$ distinct morphisms
again, compared with the same countable ring (F49.4). This proves
the first assertion, including the algebraically closed endpoint.

\emph{Proof: arbitrary algebraically closed fields.}
Let $k$ now be any algebraically closed field. It suffices to find
a nonidentity field automorphism $\sigma$ fixing an infinite subfield
$k_0\subset k$. The semilinear construction (F49.5), relative to
$\mathrm{Set}$, has point function $a\mapsto\sigma(a)$. It fixes
all $k_0$-points and moves another point. Exactly the polynomial or
finite-$p$-denominator calculation used above shows that a morphism
in (F49.4) fixing this infinite subset must be the identity.
Therefore (F49.5) is outside the image.

In characteristic zero, choose a copy of $\overline{\mathbb Q}$ in
$k$, take a nonidentity automorphism of it fixing $\mathbb Q$, and
extend it to $k$. Such extension is obtained by fixing a
transcendence basis over $\overline{\mathbb Q}$ and extending the
resulting automorphism of the rational-function field to its
algebraic closure. The fixed subfield contains $\mathbb Q$.
In characteristic $p$, choose
$\delta\in\widehat{\mathbb Z}=\prod_\ell\mathbb Z_\ell$ whose
$2$-adic coordinate is zero and whose odd-primary coordinates are
one. The automorphism $\operatorname{Frob}_p^\delta$ of
$\overline{\mathbb F}_p$ is nonidentity and fixes the infinite field
\[
                    \bigcup_{n\geq0}\mathbb F_{p^{2^n}}.
\]
Extend it to $k$ by the same transcendence-basis argument. This
supplies $\sigma$ and $k_0$ in every characteristic, and completes
all three assertions. The attached program \path{code/montpellier/F49_checks.py}
checks the framed-point distinction over $\mathbb F_9$ and 2726 finite
compatibilities of the displayed profinite exponent; the infinitary
claims follow from the proofs above. \hfill$\square$

\Needspace{5\baselineskip}
\paragraph{F.50 (Proved). Reconstruction over subfields of finitely generated cyclotomic extensions.}
\textit{Manuscript: MPL149; section2; folio25v; scan52; intermediate-base reconstruction.}

Let $k$ be any subfield of a field finitely generated over
$\mathbb Q^{\mathrm{ab}}$. Then $k$ has pinned etale reconstruction:
for all finite-type $k$-schemes $X,Y$,
\[
 \operatorname{Hom}_k(X^{\mathrm{awn}},Y)
  \ \xrightarrow{\ \sim\ }\
 \pi_0\operatorname{Hom}^{\mathrm{pin}}_{B G_k}
             (X_{\mathrm{et}},Y_{\mathrm{et}}).
 \tag{F50}
\]
The relative commutation isomorphism is part of the morphism data.
This includes F.47 and also every subfield of a finitely generated
extension of an arbitrary cyclotomic extension of a number field.
We prove the additional arithmetic and quasi-rational steps that
extend F.47; its geometric and categorical deductions then apply
with all their hypotheses verified.

\begin{proof}
Choose $k\subset M$, where $M/\mathbb Q^{\mathrm{ab}}$ is finitely
generated, and choose an algebraic closure $\Omega$ of $M$. Let
$K\subset\Omega$ be the algebraic closure of $k$. In particular
$\overline{\mathbb Q}\subset K$. For every finite extension
$L/k$ inside $K$, the compositum $M_L=ML$ is finite over $M$,
hence again finitely generated over $\mathbb Q^{\mathrm{ab}}$.
Thus the field class in the statement is stable under every finite
extension needed for points, boundary components, or coefficient
descent.

\emph{1. Divisibility and the free multiplicative quotient.}
Let $E/\mathbb Q^{\mathrm{ab}}$ be finitely generated, and let
$E_0$ be the algebraic closure of $\mathbb Q^{\mathrm{ab}}$ in $E$.
This is a finite extension (Stacks, Lemma~9.26.11
\citep{StacksFGConstants}); hence $E_0=F_0\mathbb Q^{\mathrm{ab}}$
for a number field $F_0$. The extension $E/E_0$ is regular in
characteristic zero. Choose a proper normal geometrically integral
$E_0$-model $Z$ of $E$. Its divisor map gives an injection
\[
 E^*/E_0^*\lhook\joinrel\longrightarrow
       \operatorname{Div}(Z),                              \tag{F50.1}
\]
because a rational function with zero divisor on a proper normal
variety is a global unit, and its constant field is $E_0$.
The divisor group is free abelian, so its subgroup $E^*/E_0^*$ is
free abelian. By Tsujimura's Proposition~4.9 with parameter $f=1$,
$E_0^*/\mu_\infty$ is free abelian; the identification of its
divisible subgroup with $\mu_\infty$ is the Amoroso--Zannier
height argument in F.47. Consequently the exact sequence
\[
 0\longrightarrow E_0^*/\mu_\infty
 \longrightarrow E^*/\mu_\infty
 \longrightarrow E^*/E_0^*\longrightarrow0
\]
splits as a sequence of abelian groups, because its quotient is
free. Therefore $E^*/\mu_\infty$ is free abelian.

Also $D(E^*)=\bigcap_n(E^*)^n=\mu_\infty$. If $u$ is infinitely
divisible, its divisor is divisible by every integer, hence zero.
Every chosen $n$th root of $u$ has zero divisor as well, so $u$ and
all these roots lie in $E_0$. The height theorem applies there and
forces $u$ to be a root of unity. The converse holds since $E_0$
contains all roots of unity. In particular, for every finite $L/k$,
\[
               D(L^*)\subset\mu_\infty.                  \tag{F50.2}
\]
This follows by embedding $L$ and all its division witnesses into
$M_L$; no assertion that $L$ itself contains all roots of unity is
needed.

\emph{2. Abelian-variety Kummer faithfulness.}
For each finite $L/k$ and every abelian variety $A/L$,
\[
                    \bigcap_n nA(L)=0.                    \tag{F50.3}
\]
This is also the stability statement in Tsujimura,
Proposition~1.2(ii) \citep{Tsujimura2022}, combined with his
Theorem~3.1 over finite extensions of $\mathbb Q^{\mathrm{ab}}$
\citep{Tsujimura2020}. We give the specialization argument, including
the uniformity needed for all divisions.

For a finitely generated field $E/\mathbb Q^{\mathrm{ab}}$, an
abelian variety $A/E$, and a point $P\in\bigcap_n nA(E)$, spread
the abelian variety and this one point to an abelian scheme
$\mathcal A\to U$ and a section $P_U$, where $U$ is a normal
integral finite-type scheme over $\mathbb Q^{\mathrm{ab}}$ with
function field $E$. For every $n$, choose $Q_n\in A(E)$ with
$nQ_n=P$. The finite etale scheme $[n]^{-1}(P_U)\to U$ has the
generic section $Q_n$. Normality of $U$ extends it to a section on
the same $U$: the generic section selects a finite birational
component, which is isomorphic to the normal base. Thus all the
divisions extend without an infinite sequence of shrinkings.
At every closed point $x\in U$, the residue field is finite over
$\mathbb Q^{\mathrm{ab}}$, and $P_U(x)$ is divisible by all $n$.
Theorem~3.1 gives $P_U(x)=0$. Closed points are dense and the
abelian scheme is separated, so $P_U=0$ and $P=0$.
Apply this to $E=M_L$ and the base change of $A$; the injection
$A(L)\hookrightarrow A(M_L)$ proves (F50.3). The Picard and total
Tate-module conclusions of F.47 follow exactly as there, including
the treatment of Neron--Severi torsion.

\emph{3. Rational-source integrality without algebraic coefficients.}
The rational-source lemma of F.47 must now allow distinct punctures
$b_1,\ldots,b_r\in L$ which can be transcendental over $\mathbb Q$.
Put $E=\mathbb Q(b_1,\ldots,b_r)$, a finitely generated field.
There is an integral domain $A\subset E$, smooth and of finite
type over $\mathbb Z[1/N]$, with fraction field $E$, containing
all $b_i$ and all $(b_i-b_j)^{-1}$ for $i\ne j$. To construct it,
take a smooth affine model over $\mathbb Q$, shrink it to contain
these finitely many functions as regular functions and units, and
spread the finitely presented smooth algebra over $\mathbb Z[1/N]$;
Stacks, Lemma~10.138.14 \citep{StacksSmoothSpread} supplies this
last step. Flatness makes the model inject into its generic fibre,
so it remains a domain. Its smooth image in
$\operatorname{Spec}\mathbb Z[1/N]$ is an open set containing the
generic point, hence contains all but finitely many primes.

Choose different primes $p_i>r$ with nonempty fibres, and choose
a generic point $\mathfrak p_i$ of a component of the fibre over
$p_i$. The local ring $A_{\mathfrak p_i}$ is a discrete valuation
ring and $p_i$ is a uniformizer: the model is regular and flat and
the fibre is smooth, hence generically reduced. Let $v_i$ be this
valuation of $E$. Since the pairwise differences are units in
$A$, their residues at $\mathfrak p_i$ are nonzero. Among the
$p_i$ elements of $\mathbb F_{p_i}$ choose $c_i$ different from
all the at most $r$ residues of $b_j$ that lie in that prime field.
The Chinese remainder theorem in $A$ now supplies $x,y\in A$
such that
\[
 x\equiv b_i+p_i\pmod{p_i^2A},\qquad
 y\equiv c_i\pmod{p_iA}\quad(1\le i\le r).
\]
The ideals for different rational primes are comaximal. Therefore
\[
 v_i(x-b_j)=\delta_{ij},\qquad v_i(y-b_j)=0.                \tag{F50.4}
\]
In particular $x,y$ lie in the punctured affine line over $L$,
and the ratios $(x-b_j)/(y-b_j)$ are multiplicatively independent
modulo roots of unity. Independence persists in the overfield
$M_L$, since it is a statement about equalities of these same
elements.

Part~1 makes $M_L^*/\mu_\infty$ free abelian. The character-matrix
argument (F47.R4) therefore applies without alteration: finitely
many integer characters give a nonsingular integer matrix $N$ on
these independent ratios; two actual Kummer evaluations give
$N\lambda\in\mathbb Z^r$. The adjugate matrix and
$\mathbb Q\cap\widehat{\mathbb Z}=\mathbb Z$ force
$\lambda\in\mathbb Z^r$. A single further evaluation makes the
constant completed class an actual element of $L^*$.
Thus every compatible Kummer class on a punctured affine line
with actual specializations is an actual regular-unit class, in
the same generality as required for the norm calculation in F.47.
No valuation on the possibly much larger field $L$ was asserted
to be discrete.

\emph{4. The quasi-rational degree argument in the larger algebraic closure.}
Use the fixed subfield $B=\mathbb Q^{\mathrm{ab}}\subset K$.
All pointwise ambiguities furnished by (F50.2) lie in
$\mu_\infty\subset B^*$. The proof of the finite-exception lemma
in F.46 uses high-degree closed points after adjoining finitely
many coefficients. Here those coefficients generate a finitely
generated field $N/B$ inside $K$, which need not be algebraic.
We verify the required high-degree property directly.

Let $N_0=N\cap\overline{\mathbb Q}$. It is finite over $B$ by
the constants lemma, and $N/N_0$ is regular. Hence $N$ and
$\overline{\mathbb Q}$ are linearly disjoint over $N_0$; this is
the geometric irreducibility assertion in Stacks,
Lemma~10.47.13 \citep{StacksFGConstants}, using characteristic zero
for geometric reducedness. Thus, for algebraic numbers $a$,
\[
          [N(a):N]=[N_0(a):N_0].                           \tag{F50.5}
\]
The right side is unbounded. If
$\overline{\mathbb Q}/N_0$ had finite degree, its elements would
have uniformly bounded degree over $B$, and the Amoroso--Zannier
height bound would contradict the elements $2^{1/n}$.
Finite extensions of the infinite algebraic extension therefore
have unbounded degrees; characteristic zero supplies primitive
elements.

Descend the smooth curve and the finitely many rational functions
in a given translate argument to such an $N$. A finite map to
$\mathbb P^1_N$ and values $a$ as in (F50.5) give closed points
of arbitrarily large residue degree. A nonzero linear combination
of the fixed rational functions has a uniformly bounded zero
divisor, so cannot vanish at those points. The same choice avoids
any finite exceptional set: its points algebraic over $N$ have
bounded degrees, and its points transcendental over $N$ cannot
coincide with the chosen closed points. This is exactly the
degree argument used to obtain the rational identities, the two
covering outside-$B$ sets, and their infinite intersection in
F.46. Also $K/B$ is infinite dimensional, since it contains
$\overline{\mathbb Q}/B$, so the constant-factor case of that
proof applies. Consequently the lemma holds here without assuming
that $K/B$ is algebraic. We have verified the proof's actual
hypothesis, rather than using the algebraic-extension version of
Tsujimura's criterion beyond its stated range.

\emph{5. Completion of the reconstruction argument.}
The curve argument of F.47 now applies with its three inputs
verified for every finite $L/k$: the norm class on the punctured
line has integral residues by part~3; the integral divisor class
and the total Picard Tate module vanish by (F50.3); and the
pointwise quotient is a root of unity by (F50.2). Part~4 gives a
rational point function from all translates, including every
exceptional value. Galois equivariance descends that function,
and the rational-point Hochschild--Serre argument identifies its
entire compatible Kummer class.

The higher-dimensional argument of F.47 uses only characteristic
zero, these curve assertions after finite extensions, and
(F50.3). Smooth compactification and transverse curves make the
boundary residues integral; the Picard/Neron--Severi argument
makes their divisor principal; the remaining class is constant;
and a curve through the fixed rational point and an arbitrary
geometric point proves pointwise equality with one global unit.
These hypotheses all hold here. Finally Carlson--Haine--Wolf,
Theorem~4.18, applies to every regular affine source and all its
sliced etale objects, giving (F50). Thus the conclusion concerns
the full relative geometric morphisms after localization, including
constants and noninvertible maps, not only their point functions
or their fundamental-group homomorphisms.
\end{proof}

\Needspace{5\baselineskip}
\paragraph{F.51 (Refuted). Independent atlas germs need not be intrinsic stack neighbourhoods.}
\textit{Manuscript context: MPL147, \S\S11--12, folios72,74,76,
scans146,150,155; the unrestricted stack extensions of the gluing and
tangential assertions.}
The following two constructions must be distinguished. For a closed
substack $Z\subset\mathcal X$, the \emph{intrinsic neighbourhood
pro-topos} is
\[
 \mathfrak V_{Z,\mathcal X}
   =\bigl(\operatorname{Sh}(\mathcal U)\bigr)_{
           Z\subset\mathcal U\subset\mathcal X\ {\rm open}}.
                                                        \tag{F51.1}
\]
Given an etale presentation $X_\bullet\to\mathcal X$ with
$Z_q=Z\times_{\mathcal X}X_q$, one can instead first take the full
neighbourhood germ of $Z_q$ in each $X_q$, independently, and then
form its simplicial diagram in pro-topoi. The canonical cocone from
this diagram to (F51.1) is \emph{not} always a two-colimit. The
failure occurs even for an analytic etale groupoid with trivial
isotropy. Moreover, for a nonproper presentation the intrinsic
punctured neighbourhood can have a different covering groupoid from
the punctured normal stack.

The strongest unconditional replacement retains the entire atlas
diagram over one actual open-substack stage before passing to the
pro-object. Ordinary intrinsic two-open gluing remains valid for
every such stack. Independent chart-germ descent and F.13's controlled
flag construction agree with the intrinsic object for a proper
etale groupoid with paracompact locally compact orbit space, or under
the explicit common-stage synchronization hypothesis of F.13.
Finite isotropy alone does not replace properness.

\begin{proof}
\emph{1. A full-topos counterexample.}
Let $\Gamma=\mu_\infty$ be the countable group of all complex roots
of unity, given the discrete topology, acting on the unit disc
$\Delta$ by rotation. Put
$\mathcal X=[\Delta/\Gamma]$ and $Z=[\{0\}/\Gamma]$.
This is an analytic etale quotient stack; its groupoid is not proper.
Its atlas nerve is $X_q=\Gamma^q\times\Delta$, with
$Z_q=\Gamma^q\times\{0\}$. A neighbourhood of $Z_q$ permits an
independent disc radius on each of the discrete components
$\Gamma^q$. By contrast an intrinsic neighbourhood is $[U/\Gamma]$
for a single invariant open $U\subset\Delta$ containing $0$.

For $n\ge1$, put
\[
 A_n=\{2^{-n-1}<|z|<2^{-n}\},\qquad
 b_n(r)=2^{-n}\max\{0,1-16|2^nr-3/4|\}.
\]
The support of $b_n$ is strictly inside the radial interval of
$A_n$. Define a continuous real function by
\[
 f(0)=0,\qquad
 f(re^{i\theta})=b_n(r)\cos(n!\theta)\quad(r\in A_n),
                                                        \tag{F51.2}
\]
and set it to zero elsewhere. Continuity at the annular boundaries
comes from the support condition; continuity at zero follows from
$|f|\le2^{-n}$ on $A_n$.
If $\zeta$ has order $m$, then for every $n\ge m$ the integer
$n!$ is divisible by $m$. Hence
\[
 f(\zeta z)=f(z)\quad\text{for }|z|<2^{-m}.                \tag{F51.3}
\]
For every finite tuple of roots, choose $n$ at least their maximum
order; the same equality holds for all products in the subgroup
they generate. Consequently the germ of $f$ gives compatible maps
from every term of the atlas-germ diagram to
$\operatorname{Sh}(\mathbb R)$. All face compatibilities and their
cocycles hold: each fixed tuple has a representative radius on
which the required function equalities are literal. The target is
localic and Hausdorff, so the corresponding mapping categories here
are discrete; no unrecorded two-cell choice is required.

This cone does not come from (F51.1). A geometric morphism
$\operatorname{Sh}([U/\Gamma])\to\operatorname{Sh}(\mathbb R)$
pulls back along $U$ to a continuous $\Gamma$-invariant real
function. This follows either from fixed-stage etale descent or from
the inverse-image frame homomorphism into the invariant open subsets
of $U$. Since $\Gamma$ is dense in the circle group, continuity makes
the function constant on each circle contained in $U$.
Every neighbourhood of zero contains a disc. But at
$r_n=3\cdot2^{-n-2}$, the two angles $0$ and $\pi/n!$ give
values $2^{-n}$ and $-2^{-n}$ in (F51.2). Thus $f$ is nonradial on
every disc. It cannot agree as a germ with any such invariant
representative. The canonical mapping-out functor to this fixed
target is not essentially surjective, proving the failure of the
pro-topos colimit universal property.

\emph{2. The failure is not repaired by finite stabilizers.}
Let $\Gamma$ instead act on
$Y=\mathbb C^*\times\Delta$ by
$\zeta(z,w)=(\zeta z,w)$, and put
$Z_0=\mathbb C^*\times\{0\}$.
This action is free. Define
\[
 F(z,w)=b_n(|w|)\cos(n!\arg z)\quad(|w|\in A_n),
 \qquad F(z,0)=0,
                                                        \tag{F51.4}
\]
with zero values in the remaining radial intervals. The expression
is single-valued because $n!$ is integral. Equations (F51.3) hold
uniformly in $z$ after shrinking the $w$-radius for each fixed group
element, so again this is a full coherent atlas-germ cone.

If it came from an invariant neighbourhood representative, equality
of the two germs would hold on an ordinary neighbourhood of $Z_0$.
Compactness of $\{|z|=1\}\times\{0\}$ supplies a single positive
$w$-radius on which that equality holds for the whole circle.
Invariance and continuity force the representative to be constant
in $\arg z$ on that circle, whereas (F51.4) has the same two opposite
values for arbitrarily small $w$. This is a contradiction. All
isotropy groups in this example are trivial; its dense circle
orbits make the groupoid nonproper. Properness, not merely a bound
on stabilizers, is material.

\emph{3. The intrinsic covering comparison can fail as well.}
Let $\mathbb Z$ act on $\mathbb P^1$ by
\[
 g(z)=\frac{z}{1+z},\qquad D=\{0\}.
                                                        \tag{F51.5}
\]
The action is understood projectively at the pole and at infinity.
In the coordinate $w=1/z$, its restriction to
$\mathbb P^1\setminus D=\mathbb C_w$ is $w\mapsto w+1$.
Every such orbit accumulates at $D$ in $\mathbb P^1$.
Therefore an invariant open neighbourhood of $D$ is necessarily
all of $\mathbb P^1$: it meets every orbit in an end of its sequence,
and invariance then contains that entire orbit. The intrinsic
punctured neighbourhood pro-topos is consequently constant, with
value
\[
 \operatorname{Sh}([\mathbb C_w/\mathbb Z])
       \simeq\operatorname{Sh}(\mathbb C^*),
                                                        \tag{F51.6}
\]
where the quotient map is $w\mapsto e^{2\pi i w}$. Its fundamental
group is $\mathbb Z$.

On the other hand $g'(0)=1$, so the action on the normal line at
$D$ is trivial. Its punctured normal stack is
$[\mathbb C^*/\mathbb Z_{\rm trivial}]$, with classifying
homotopy type $S^1\times B\mathbb Z$ and fundamental group
$\mathbb Z^2$. The tangential comparison sends the normal meridian
to zero: the small circle $z=\epsilon e^{i\theta}$ becomes a
closed circle in the simply connected plane $\mathbb C_w$.
The operator generator is represented by the path $w\to w+1$,
which becomes the standard loop of the quotient. Thus on these
fundamental groups the comparison is
\[
 \mathbb Z\oplus\mathbb Z\longrightarrow\mathbb Z,
                    \qquad(a,b)\longmapsto b.              \tag{F51.7}
\]
It is not an equivalence, even on finite covering categories: a
nontrivial finite normal-meridian action does not descend through
(F51.7). Hence F.19's intrinsic-stack equivalence requires its
properness or common-stage hypothesis. Its coherent chart-germ
construction is a different assertion and is not refuted by this
example. The linear unipotent example of F.20 still proves exactly
the stated failure of an invariant differentiable tubular embedding;
that failure alone is not an argument for or against full-topos
germ descent.

\emph{4. The common-stage replacement for every stack.}
Retain an actual open substack $\mathcal U\supset Z$ and its full
atlas nerve $X_\bullet\times_{\mathcal X}\mathcal U$.
At this one stage ordinary etale descent gives
$\operatorname{Sh}(\mathcal U)$ as the bicategorical colimit of
the sheaf topoi of the nerve. Now let \emph{this entire diagram}
vary with $\mathcal U$. Mapping to a fixed target topos uses the
inverse-image functor on degree zero, an invertible comparison on
degree one, and its cocycle equality on degree two. Each item is
represented at one actual $\mathcal U$; their finitely many stage
indices have a common refinement. The same argument treats natural
transformations and their equalities. Thus filtered two-colimits
of these mapping categories commute with this descent pseudolimit.
For a pro-target take the further limit over target stages. This
proves the colimit assertion for the synchronized pro-diagram.
It does not claim that its degree-$q$ term is cofinal among all
independent neighbourhoods of $Z_q$ in $X_q$; Parts~1--2 exhibit
exactly the failure of that claim.

Likewise, at each actual stage $\mathcal U$, the opens
$\mathcal X\setminus Z$ and $\mathcal U$ cover $\mathcal X$.
Their two-open descent is a finite pseudolimit on mapping
categories. The same common-refinement argument proves
\[
 \operatorname{Sh}(\mathcal X)\simeq
 \operatorname{Sh}(\mathcal X\setminus Z)
   \amalg_{\mathfrak V^*_{Z,\mathcal X}}
                 \mathfrak V_{Z,\mathcal X}
                                                        \tag{F51.8}
\]
in pro-topoi, with the left side constant, without a properness
assumption. All neighbourhoods in this formula are actual open
substacks throughout.

\emph{5. Why proper etale presentations admit the chartwise repair.}
A proper etale groupoid admits local charts by finite isotropy
groups, without an effectiveness assumption; see Moerdijk (2002,
\S3.4) \citep{FA_Moerdijk2002}. Choose a locally finite relatively
compact shrinking of these charts on the paracompact locally
compact orbit space. Over a pair or triple of compact chart
closures, properness makes the space of relevant arrows or
composable arrows compact. Etaleness therefore covers it by
finitely many local bisections. Thus the chart datum, overlap
isomorphisms, cocycle equalities, and comparisons of two such data
impose only finitely many radius bounds near any point of this
locally finite system.

First take the minimum over the finite isotropy group in each
chart. Next shrink the compact chart cores so every representative
bound is positive there. A partition of unity on the orbit space
supplies one positive continuous minorant of these locally finite
bounds. Its pullbacks give invariant compatible tube radii; the
finite control identities for incident branches are incorporated
by the successive shrinking already proved in F.13. Descent on
the resulting common neighbourhood realizes the full datum.
Applying the same procedure to comparison maps and their
equalities proves full faithfulness as well as essential
surjectivity. This proves the proper-stack version of F.13 and
identifies its independently formed chart germs with the intrinsic
construction, including all flag data. The argument retains the
finite stabilizers; it does not replace a quotient stack by its
coarse space.

This procedure does not assert that an arbitrary neighbourhood in
an unbounded atlas contains an invariant neighbourhood. Even for
the proper translation action of $\mathbb Z$ on
$\mathbb R_x\times\mathbb R_y$, the neighbourhood
$|y|<e^{-x^2}$ of $y=0$ has no invariant refinement. One first
restricts the germ datum to the relatively compact quotient charts,
performs the finite synchronizations there, and glues a stack
neighbourhood. Comparison with the original atlas datum is then by
its individual common refinements. These comparisons are precisely
the morphisms in the atlas-germ descent category; no uniform lower
bound along an entire noncompact orbit is used.

For a nonproper presentation one may instead require exactly this
common-stage realization property as part of its controlled atlas
data. The synchronized statement of Part~4 then applies. No
arbitrary interchange of a filtered germ colimit with an infinite
atlas descent diagram is included in either replacement.

The exact certificates in \path{code/foundations/F51_checks.py} check the finite
rotation divisibilities, tuple compatibilities, and parabolic composition
identities used above; their actual output is attached.
\end{proof}
